\documentclass[openany]{book}

\makeatletter
\let\Title\@title
\let\Author\@author
\makeatother

%\usepackage{amssymb}
\usepackage{amsmath}
\usepackage{mdwlist}
\usepackage[compact,sc]{titlesec}
\usepackage{graphicx}
\usepackage[geometry,weather,misc,clock]{ifsym}
\usepackage{array}
\usepackage{hanging}
\usepackage{hyphenat}
%\usepackage{epsfig}
%\usepackage{egame}
\usepackage{booktabs}
\usepackage{nicefrac}
%\usepackage{bm}

\usepackage[pagebackref]{hyperref}
\hypersetup{colorlinks=true,linkcolor=black,citecolor=black}
\usepackage{etoolbox}
\makeatletter
\patchcmd{\BR@backref}{\newblock}{\newblock(Page~}{}{}
\patchcmd{\BR@backref}{\par}{)\par}{}{}
\makeatother


\usepackage{needspace}
\usepackage{textcomp}
%\usepackage{substr}
\usepackage{natbib}
\usepackage{nicefrac}
\usepackage{enumitem}
\setlist{noitemsep}
\usepackage{multirow}
\usepackage{hyphenat}
\usepackage{tabulary}
%\usepackage{booktabs}

\usepackage[overload]{textcase}

\renewcommand{\phi}{$\varphi$}
\newcommand{\SP}{$S ^\prime$}
\DeclareMathAlphabet{\mathpzc}{OT1}{pzc}{m}{it}
\newcommand{\eqm}{equilibrium}
\newcommand{\Eqm}{Equilibrium}
\newcommand{\NE}{Nash Equilibrium}
\newcommand{\NEs}{Nash Equilibria}
\newcommand{\sth}{$^{\text{th}}$}
\newcommand{\snd}{$^{\text{nd}}$}
\newcommand{\sst}{$^{\text{st}}$}
\newcommand{\ol}[1]{$\langle #1 \rangle$}
\usepackage{pict2e}
\usepackage{color}
\usepackage{cclicenses}


%\usepackage[log-declarations=false]{xparse}

\usepackage{chngcntr}
\counterwithout{footnote}{chapter}


\renewcommand{\labelitemi}{$\bullet$}

%\usepackage[garamond]{mathdesign}
\usepackage[no-math]{fontspec}
%\setmainfont[Ligatures=TeX, BoldFont = Minion Pro Semibold]{Minion Pro}
%\setmainfont[Ligatures=TeX, BoldFont = Gill Sans SemiBold]{Gill Sans}
%\setmainfont{Perpetua}
%\setmainfont[Ligatures=TeX, Mapping=tex-text]{Constantia}
%\setmainfont{Adobe Garamond Pro}
%\setmainfont[Ligatures=TeX, Mapping=tex-text, BoldFont = Adobe Caslon Pro Semibold]{Adobe Caslon Pro}
\setmainfont[Mapping=tex-text, BoldFont = Baskerville SemiBold]{Baskerville}
%\setmainfont{Old Standard TT}
\usepackage[defaultmathsizes,italic]{mathastext}

\titlespacing*{\subsection}{0pt}{9pt}{3pt}
\titleformat{\subsection}{\normalfont \normalsize \itshape}{\thesubsection}{1em}{}
\titlespacing*{\section}{0pt}{9pt}{3pt}
\titleformat{\section}{\normalfont \large \itshape}{\thesection}{1em}{}
\titlespacing*{\chapter} {0pt}{30pt}{20pt} 
\titleformat{\chapter}[display]{\huge \centering}{
%\chaptertitlename\ 
\thechapter}{1em}{}

\widowpenalty=150
\clubpenalty=150

\setlength{\parskip}{0ex plus 2ex minus 0ex}
\setlength{\headsep}{16pt}
\setlength{\headheight}{14pt}
%\setlength{\textheight}{540pt}
%\setlength{\textwidth}{360pt}
\setlength{\textheight}{517pt}
\setlength{\textwidth}{320pt}

\setlength{\marginparwidth}{151pt}
\setlength{\oddsidemargin}{75.5pt}
\setlength{\evensidemargin}{75.5pt}

%\setlength{\marginparwidth}{111pt}
%\setlength{\oddsidemargin}{55.5pt}
%\setlength{\evensidemargin}{55.5pt}

\raggedbottom

\usepackage{fancyhdr}
\pagestyle{fancy}
\renewcommand{\chaptermark}[1]{\markboth{#1 }{}}
\renewcommand{\sectionmark}[1]{\markright{#1}}
\fancyhf{}
\fancyfoot{}
%\fancyfoot[C]{\thepage}
\fancyhead[OR]{\thepage}
\fancyhead[EL]{\thepage}
\fancyhead[EC]{\small \MakeUppercase{\leftmark}}
\fancyhead[OC]{\small \MakeUppercase{\rightmark}}
\renewcommand{\headrulewidth}{0pt}
\renewcommand{\footrulewidth}{0pt}
\addtolength{\headheight}{10.5pt}
\fancypagestyle{plain}{
  \fancyhead{}
  \renewcommand{\headrulewidth}{0pt}
}

\def\thesection{\arabic{chapter}.\arabic{section}}

%\newcommand{\PPI}{Reprinted in \textit{Philosophical Papers}, Volume I, pp. \hbox{}}   
%\newcommand{\PPII}{Reprinted in \textit{Philosophical Papers}, Volume II, pp. \hbox{}}    
%\newcommand{\PPL}{Reprinted in \textit{Papers in Philosophical Logic}, pp. \hbox{}}    
%\newcommand{\PME}{Reprinted in \textit{Papers in Metaphysics and Epistemology}, pp. \hbox{}}    
%\newcommand{\PESP}{Reprinted in \textit{Papers in Ethics and Social Philosophy}, pp. \hbox{}}

\DeclareSymbolFont{symbolsC}{U}{txsyc}{m}{n}
\DeclareMathSymbol{\boxright}{\mathrel}{symbolsC}{128}
\DeclareMathAlphabet{\mathpzc}{OT1}{pzc}{m}{it}

%\newcommand{\numbex}[2]{
%\begin{enumerate*}
%\setcounter{enumi}{#1}
%\renewcommand{\labelenumi}{(\arabic{enumi})}
%#2
%\end{enumerate*}
%}

%\newcommand{\starttab}{\begin{center}
%\vspace{6pt}
%\begin{tabular}}
%
%\newcommand{\fintab}{
%\end{tabular}
%\end{center}
%}
%
%\newcommand{\stoptab}{\end{tabular}
%\vspace{6pt}
%\end{center}
%\noindent}
%
%\newcommand{\gamelab}[1]{\newcounter{#1}\setcounter{#1}{\value{games}}\textbf{Game \arabic{games}}\addtocounter{games}{1}}
%
%\newcommand{\regame}[1]{
%\textbf{Game \arabic{#1}}}
%
%\newcommand{\citegame}[1]{Game \arabic{#1}}
%
%\newcommand{\stratcomp}[2]{
%\needspace{10pt}
%\smallskip
%\noindent \textbf{#1}
%\vspace{-6pt}
%\begin{quote}
%#2
%\end{quote}}
%
%\newcommand{\pictext}[3]{
%\put(#1, #2){\makebox(0, 0)[b]{#3}}}
%
%\newcommand{\drawsquare}[5]{
%\put(#1, #2){\line(0, 1){#5}}
%\put(#1, #2){\line(1, 0){#5}}
%\put(#3, #2){\line(0, 1){#5}}
%\put(#1, #4){\line(1, 0){#5}}
%}
%
%\newcommand{\tol}[1]{$\langle #1 \rangle$}
%
%\newcounter{games}
%\setcounter{games}{1}

\DeclareSymbolFont{symbolsC}{U}{txsyc}{m}{n}
\DeclareMathSymbol{\boxright}{\mathrel}{symbolsC}{128}

\hyphenation{Bayes-ian-ism de-prag-mat-ised prag-mat-ised Schlen-ker rath-er Gil-bert cele-bri-ties thou-gh O-phe-lia grok-king-ness Gend-ler Al-iens non-fund-a-men-tal grok-king Dog-gett phe-nom-e-non e-pi-stem-olog-ist Will-iam-son Witt-gen-stein meta-ethics Key-nes third-er loc-ation-al third-er sp-ell-ings Des-cartes con-glom-er-able}


















\newcommand{\pictext}[3]{
\put(#1, #2){\makebox(0, 0)[b]{#3}}}

\newcommand{\drawsquare}[5]{
\put(#1, #2){\line(0, 1){#5}}
\put(#1, #2){\line(1, 0){#5}}
\put(#3, #2){\line(0, 1){#5}}
\put(#1, #4){\line(1, 0){#5}}
}

\begin{document}

\frontmatter
{\pagestyle{plain}
\fancyhf{}
\markboth{ }{}

%\newcommand*{\FSfont}[1]{\fontencoding{T1}\fontfamily{#1}\selectfont}
\newlength{\drop}% for my convenienc
{\begingroup% Robert Frost, T&H p 149
%\FSfont{5bp} % FontSite Bergamo (Bembo)
\drop = 0.2\textheight
\centering
\vfill
{\Huge Lecture Notes on}\\[\baselineskip]
{\Huge GROUPS AND CHOICES}\\[\baselineskip]
{\large Brian Weatherson}\\[0.5\drop]
\vfill
\includegraphics{cclicence} \\
%{\Large \plogo}\\[0.5\baselineskip]
%{\Large \cc \ccnc \ccsa}\\[0.5\baselineskip]
%{\Large The Publisher}\par
{\large\scshape Winter 2017}\par
\vfill\null
\endgroup}

\setcounter{tocdepth}{1}
\tableofcontents
\cleardoublepage}

\mainmatter

\part{Game Theory}

\chapter{Basics of Game Theory}

\section{Nature of Games}

We'll think of a game as any situation where multiple players have to make some decisions, and the outcome they get is a function of (perhaps among other things) the decisions that each make. So understood, that's a very large class of situations. Indeed, game theory has been used to analyse any number of real-world interactions, as you might expect from a definition so broad.

We're going to start our focus on a special class of games, two player simultaneous move games. The term `two player' probably obvious; it means that we're looking at games where just two players are involved. Simultaneous move games are where each player only gets to make one decision, or move. Each player makes that decision in ignorance of what the other players have decided. And once the decisions are made, the payoffs for each player can be revealed.

This feels like a very narrow class of games. After all, most interactions involve multiple moves by at least one player, or at least players making moves at different times. But there is a good sense in which this is not a particularly tight restriction. Think of a game where the players make multiple moves, like chess. We can imagine a game between two chess super-computers that have planned, in advance, what they will do at every possible stage of the game. Call such plans \textbf{strategies}. If two such computers played, they would in a sense each just have to make one, simultaneous, move, namely to announce their strategies. Then it would be relatively trivial to work out the result of the game from the strategies announced.

In practice, this isn't really possible with current technology. There are too many possible states of the board, even in a restricted game like chess, for a computer to have planned for every possible eventuality. But we can imagine such complete strategies for simple games like tic-tac-toe, and even for moderately complicated games like checkers. 

For these reasons, when we are discussing simultaneous move games, we will often refer to the players' options as \textit{strategies}. The representation of the game as an $n \times m$ matrix, where one player (represented on rows) has $n$ options, and the other player (represented on columns) has $m$ options, is sometimes called a \textbf{normal form} representation of a game and sometimes called a \textbf{strategic form} representation. And we'll start off with these matrix representations. When we complicate things by explicitly representing time in the game, we'll call the representation the \textbf{extensive form} representation of the game.

\section{Utility}
Here is a representation of an important kind of two player simultaneous move game.

%Prisoners' Dilemma - Specific
\begin{center}
\begin{tabular}{l r | c c}
 & & \multicolumn{2}{c}{\textbf{Col}} \\
& & X & Y \\ \hline
\multirow{2}{*}{\textbf{Row}} & X & 3, 3 & 0, 5 \\ & Y & 5, 0 & 1, 1 \\
\end{tabular}
\end{center}

\noindent Each player, Row and Col, has two options: X and Y. The payoffs, given their moves, are as follows:

\begin{itemize}
\item If both play X, both get 3.
\item If both play Y, both get 1.
\item If one plays X and the other Y, then the one who plays X gets 0, and the one who plays Y gets 5.	
\end{itemize}

\noindent What does it mean to say that they `get' 0, 1, 3 or 5? It means that the \textbf{utility} they get from the payout is that amount. For any given game, the results for the players can be specified in two related, but importantly distinct, ways.

\begin{itemize}
\item We can specify their \textbf{payouts}. These are the concrete things the players get. In the first example in Leyton-Brown and Shoham, the payout the players get is that their average packet delay is a certain amount. Payouts might be specified in dollars, or grades, or goals, or any thing we like.
\item Alternatively, we can specify the \textbf{utility} each player gets from their payout.
\end{itemize}
Utility is a somewhat technical, somewhat theoretical, notion. 

The idea is to numerically measure how good an outcome is to an agent. It is meant to sum up every aspect of outcomes that agents care about. For instance, if an agent cares a lot about their friend's payout, then their utility in an outcome will be a function not just of their own payout, but of their friend's. So if her payout is \$$x$ and her friend's payout is \$$y$, then her utility might equal $x + 0.5y$. This would represent that she cares about her friend's payout, but cares a bit more (well, twice as much) about her own outcome.

The substantive assumption here is that we can sensibly map the goodness of outcomes onto the real numbers. This assumes that goodness is, in a sense, one-dimensional. For any two numbers $x, y$, $x$ is either greater than, less than, or equal to $y$. So one assumption we're making is that any two outcomes are comparable in this way. We're also assuming that we can compare the distance between pairs of outcomes. For any numbers $x_1, y_1, x_2, y_2$, we can ask which of $x_1 - y_1$ and $x_2 - y_2$ is greater. So if we can map outcomes onto numbers, then for any four outcomes, we can ask if the difference between the first two in goodness is more than, less than, or the same as, the difference in goodness between the last two. This isn't an absurd assumption, though it can seem a little arbitrary in some cases.

Although we assume that the goodness to an agent of outcomes can be mapped onto numbers, we do not assume that there is a unique way to do this mapping. In fact there are two degrees of freedom in choosing a mapping. The outcome that gets utility zero can be any point is convenient to use. In practice, it is often convenient to give the status quo utility zero. But note this is arbitrary, we could use anything. The scale of the utility function is also arbitrary; we can multiply the utilities by a positive constant and it won't change their meaning.

Technically, if one utility function correctly represents an agent, then so will any other function generated from it by a positive affine transformation. So if $u_1$ is a utility function that correctly represents the agent, then so will the function $u_2$ defined as follows.

\begin{equation*}
u_2(X) = au_1(X) + b \text{, where }a\text{ is positive}.
\end{equation*}
What does it mean to say that a function $u$ correctly represents the agent? First, it means that the agent prefers $X$ to $Y$ just in case $u(X) > u(Y)$. Second, it means that for any $X, Y, Z$ the agent prefers $Y$ to a gamble that returns $X$ with probability $p$ and $Z$ with probability $1-p$ iff $u(Y) > pu(X) + (1-p)u(Z)$.

When it isn't stated otherwise, we will describe the outcomes of games in terms of the utility each agent gets. Sometimes we will have to just state what the payout is each agent gets, and the utility function will be either unknown, or will have to be figured out. But the default is that we will list the utility of each outcome.

\section{Four Examples of Games}
We'll start with a very simple class of games: 2 player simultaneous move games where each player has 2 choices for their move. As we'll see, there is plenty of complexity within even this class.

\subsection{Prisoners' Dilemma}

We've already seen a version of this game.

%Prisoners' Dilemma - Specific
\begin{center}
\begin{tabular}{l r | c c}
 & & \multicolumn{2}{c}{\textbf{Col}} \\
& & X & Y \\ \hline
\multirow{2}{*}{\textbf{Row}} & X & 3, 3 & 0, 5 \\ & Y & 5, 0 & 1, 1 \\
\end{tabular}
\end{center}
The general form of the game is given here, assuming $a > b > c > d$, and $2b > a + d$.

%Prisoners' Dilemma - General
\begin{center}
\begin{tabular}{l r | c c}
 & & \multicolumn{2}{c}{\textbf{Col}} \\
& & X & Y \\ \hline
\multirow{2}{*}{\textbf{Row}} & X & b, b & d, a \\ & Y & a, d & c, c \\
\end{tabular}
\end{center}
One characteristic of the game is that for each player, playing $Y$ is better than playing $X$ no matter what the other person does. That is, playing $Y$ \textbf{dominates} playing $X$; it is better no matter what. But, and this is what gives the game its distinctive character, it is better for both players if they both play $X$ than if they both play $Y$. (And in a repeated game it is better to always play $X$ than to alternate between the outcomes where players make opposite moves.) It is an interesting way to think about modeling situations where there are rewards to cooperation, but incentives to individualist behavior.

\subsection{Stag Hunt}
The next game we'll look at is called Stag Hunt. We'll again do a version with commonly used numbers, then the general version.
%Stag Hunt - Specific
\begin{center}
\begin{tabular}{l r | c c}
 & & \multicolumn{2}{c}{\textbf{Col}} \\
& & X & Y \\ \hline
\multirow{2}{*}{\textbf{Row}} & X & 5, 5 & 0, 4 \\ & Y & 4, 0 & 2, 2 \\
\end{tabular}
\end{center}
The general version has $a > b > c > d$ and $b + c > a + d$
%Stag Hunt - General
\begin{center}
\begin{tabular}{l r | c c}
 & & \multicolumn{2}{c}{\textbf{Col}} \\
& & X & Y \\ \hline
\multirow{2}{*}{\textbf{Row}} & X & a, a & d, b \\ & Y & b, d & c, c \\
\end{tabular}
\end{center}
Here the best outcome for all parties is the cooperative one. But if a player has no idea what the other player will do, and thinks it is 50/50 whether they will do the cooperative thing, then their best bet is to not cooperate.

The cute story behind this game comes from the eighteenth century political philosopher Jean-Jacques Rousseau. Imagine two people are setting out on the day's hunt, and are going to hunt stag or rabbits. Catching a stag will lead to more meat, but it is impossible with just one person. One person can catch a lot of rabbits, especially if no one else is hunting rabbits. In this case, X is joining the stag hunt, and Y is setting off on a solo rabbit hunt. It isn't as hard to get cooperative outcomes in this game as in Prisoners' Dilemma, but it takes some background trust.

\subsection{Battle of the Sexes}
I don't particularly like the name of this one, but it is common enough that you probably should know it.
%Battle of Sexes - Specific
\begin{center}
\begin{tabular}{l r | c c}
 & & \multicolumn{2}{c}{\textbf{Col}} \\
& & X & Y \\ \hline
\multirow{2}{*}{\textbf{Row}} & X & 2, 1 & 0, 0 \\ & Y & 0, 0 & 1, 2 \\
\end{tabular}
\end{center}
Here's the general version, with $a > b > c$,

%Battle of Sexes - General
\begin{center}
\begin{tabular}{l r | c c}
 & & \multicolumn{2}{c}{\textbf{Col}} \\
& & X & Y \\ \hline
\multirow{2}{*}{\textbf{Row}} & X & a, b & c, c \\ & Y & c, c & b, a \\
\end{tabular}
\end{center}
The thought is that both parties prefer coordination - making the same pick - to doing different things. But they differ on what they would like to coordinate on; Row would like to coordinate on X, and Column would like to coordinate on Y. The name comes from the version of the game where Row and Column are husband and wife, and X and Y are stereotypically male/female things to prefer doing. You can probably fill in your own example here.

\subsection{Chicken}
Our last game is not for the faint of heart.
%Chicken - Specific
\begin{center}
\begin{tabular}{l r | c c}
 & & \multicolumn{2}{c}{\textbf{Col}} \\
& & X & Y \\ \hline
\multirow{2}{*}{\textbf{Row}} & X & -100, -100 & 1, -1 \\ & Y & -1, 1 & 0, 0 \\
\end{tabular}
\end{center}
The idea comes from a James Dean movie. Two cars are lined up facing each other at the end of a long straight. They start driving at each other. The drivers can either stay on the road (that's X), or swerve to the right (that's Y). If they both stay on the road, they die in a fiery inferno. If one swerves and the other stays, the one who stays wins the game, and the other loses. If they both swerve, the game is a draw. 

Unlike the previous games we've looked at, there is no pure \textbf{equilibrium} to this game. That is, there is no point in the game where no one regrets their choice on finding out what the other players choose. There is an equilibrium to the game, it turns out, but getting to it requires a bit of background.

\section{Types of Games}

Let $u(x, o)$ be the utility that player $x$ gets in outcome $o$. Then a two player game, with players $p_1, p_2$ is \textbf{zero-sum} if there is some constant $c$ such that for any outcome $o$, $u(p_1, o) + u(p_2, o) = c$. That is, however good an outcome is for $p_1$, it is equally bad for $p_2$, and vice versa. Pure competitions are zero-sum, though real zero-sum interactions are rare in the real world.

By convention, when a game is zero-sum, we rescale the utilities so the constant $c$ is 0. We always can do this - remember any positive affine transformation of the utility function doesn't change its meaning. That's why we can talk about zero-sum games, and not just constant sum games.

You might worry that whether a game is zero-sum or not depends on completely arbitrary choices about how to scale the utility functions. Strictly speaking, a game is zero-sum if there is \textit{any} pair of utility functions that correctly represents the agents, and which makes the utility each gets from any outcome sum to a constant. 

A game is \textbf{cooperative} if in any outcome, each player gets the same utility. As with the definition of zero-sum, it doesn't matter whether the game looks cooperative on every way of representing the utilities of the players, just if it looks cooperative on some representation.

\section{Dominant Strategies}
Say that one player has available to them strategies $s_1, \dots, s_n$, and the other players have available to them collectively strategies $v_1, \dots, v_m$. In the case of a two player game, the strategies available to other players will just be strategies. But in games with more than two players, we'll think of the alternatives available to other players as \textit{vectors}, with each dimension being a choice available to one of the other players. So $v_17$ might be that player 2 plays X, and player 3 plays Y, and player 4 plays Z, and player 5 also plays Z. Every one of these vectors specifies the choices of all other players in the game. 

Let $u_x(s_i, v_k)$ be the utility that $x$, the player we're focussing on, gets if she plays $s_i$ and the others play $v_k$. Then strategy $s_i$ \textbf{strongly dominates} strategy $s_j$ if and only if the following condition is satisfied.

\begin{equation*}
\forall m: u_x(s_1, v_m) > u_x(s_2, v_m)
\end{equation*}
That is, no matter what the other players end up playing, $x$ will do better playing $s_1$ than playing $s_2$.

We'll say that strategy $s_i$ \textbf{weakly dominates} strategy $s_j$ if and only if the following condition is satisfied.

\begin{equation*}
\forall m: u_x(s_1, v_m) \geq u_x(s_2, v_m) \& \exists m: u_x(s_1, v_m) > u_x(s_2, v_m)
\end{equation*}
That is, playing $s_1$ can't do any worse than playing $s_2$, no matter what the others do, and it could do better.

\section{Dominance and Solving Games}
Here is one simple way to make progress on a game. 

\begin{enumerate*}
\item Assume that no player will play a strongly dominated strategy.
\item Assume that every player will make a rational choice given that no player will play a dominated strategy.
\end{enumerate*}
The first step here lets us solve Prisoners' Dilemma. Since $Y$ dominates $X$, neither player will play $X$. So the solution to the game is that both players play $Y$. Here is another game that isn't solved by the first step, but is solved by the second step.

%Dominant Strategy Example
\begin{center}
\begin{tabular}{l r | c c}
 & & \multicolumn{2}{c}{\textbf{Col}} \\
& & X & Y \\ \hline
\multirow{2}{*}{\textbf{Row}} & X & 2, 1 & 3, 0 \\ & Y & 0, 0 & 1, 1 \\
\end{tabular}
\end{center}
Here Column has no dominant strategy; what's best for them to play is a function of what Row plays. But Row does have a dominant strategy. Whatever Column does, they are better off playing Y. So we should assume they will play Y. And given they play Y, Column should also play Y. So the solution to the game again is that both players play Y.

In later chapters, we'll look what happens when we iterate this idea. But first we need an introduction to probability theory.

% Crash course from St Andrews notes plus Arsenal example
\chapter{Probability and Decision Theory}
\section{Probability}
We'll start by going over the notion of probability that is used in decision theory. The best way to think about probability functions is by thinking about \textbf{measure} functions. Intuitively, a measure function is a non-negative \textbf{additive} function. That is, it is a function which only takes non-negative values, and where the value of a whole is equal to the sum of the value of the (non-overlapping) parts. In principle the domain of a measure function could be the size of any set, but we'll only consider functions with finite domains for now.

Such functions are frequently used in the real world. Consider, for instance, the functions from regions to their population and land mass. This table lists the (very approximate) population and land mass for England, Scotland, Wales and Northern Ireland.

\begin{center} \begin{tabular}{r c c}
& Population & Land Mass \\
& (Millions) & (Thousands of mi$^2$) \\
England & 50 & 50\\
Scotland & 5 & 30 \\
Wales & 3 & 8\\
Northern Ireland & 2 & 5\\
\end{tabular} \end{center} The nice thing about tables like this is that they give you a lot of information implicitly. You don't need a separate table to tell you that the combined population of England and Scotland is 55 million, or that the combined land mass of England and Wales is 58,000 square miles. Those (approximate) facts follow from the facts listed on the table, as well as the fact that populations and land masses are additive. Once you know the population of the parts, you know the population of the whole.

The next notion we need is that of a \textbf{normalised} measure function. A normalisted measure is a measure where every function takes the value assigned to the largest whole is 1. We often use this notion when talking about proportions. Here is a table listing, again approximately, the portion of the UK population and land mass in each of the four countries.

\begin{center} \begin{tabular}{r c c}
& Population & Land Mass \\
& (Proportion) & (Proportion) \\
England & 0.84 & 0.54\\
Scotland & 0.08 & 0.32\\
Wales & 0.05 & 0.09\\
Northern Ireland & 0.03 & 0.05\\
\end{tabular} \end{center} The way we turn the regular measure function into a normalised measure is by dividing each value on the earlier table by the value of whole. So the land mass of England is (about) 50,000 mi$^2$, and the land mass of the UK is (about) 93,000 mi$^2$, and $\nicefrac{50,000}{93,000}$ is about 0.54. So that's why we right 0.54 in the top right cell of the normalised measure table.

The last notion we need is that of a possibility space, or more formally, a \textbf{field} of propositions. Although these terms might be unfamiliar, the notion should be very familiar; it is just what we represent with a truth table. We will identify, for these purposes, propositions with sets of possible worlds. Then a field of propositions is a set of propositions such that whenever $A$ and $B$ are in the set, so are $\neg A$ and $A \wedge B$. This is, of course, just what we see with truth tables. Whenever you can represent two propositions on a truth table, you can represent their conjunction and each of their negations. Just like with truth-tables, we'll only consider finite fields. It's possible to extend this to cover any set-sized partition of possibility space, but we'll ignore such complications for now. 

A \textbf{Probability function} is just a normalised measure on a possibility space. That is, it is a function satisfying these two conditions, where $A$ and $B$ are any propositions, and $T$ is a logical truth.

\begin{itemize}
\item $0 \leq \Pr(A) \leq 1 = \Pr(T)$
\item If $A$ and $B$ can't both be true, then $\Pr(A \vee B) = \Pr(A) + \Pr(B)$.
\end{itemize}

\noindent We normally suppose that any rational agent has \textbf{credences} over the space of (salient) epistemic possibilities, and these credences form a probability function.

In any finite field of propositions, there will be some \textbf{atoms}. These are propositions such that no stronger consistent proposition is in the field. (Some, but not all, infinite fields are atomic.) We will use these atoms a lot. I'll use the variable $w$ to range over atoms, reminding us that for practical purposes, atoms behave a lot like \textit{worlds}. They aren't worlds, of course; they only settle the truth values of a handful of propositions. But they are divisions of possibility space that are maximally precise for present purposes. That is, they are being treated as worlds. On a pragmatist approach to modality, there's no more to being a world than to be treated as a world in the right circumstances, and thinking about probability encourages, I think, such a pragmatist approach.

\section{Conditional Probability}
So far we've talked simply about the probability of various propositions. But sometimes we're not interested in the absolute probability of a proposition, we're interested in its \textbf{conditional} probability. That is, we're interested in the probability of the proposition \textit{assuming} or \textit{conditional on} some other proposition obtaining.

It isn't too hard to visualise how conditional probability works if we think of a possibility space simply as a truth table, so a measure over a possibility space is simply a measure over on the truth table. Here's how we might represent a simple probability function, defined over the Boolean combinations of three propositions, if we're thinking about things this way. 

\begin{center} 
\begin{tabular}{c c c | c}
$p$ & $q$ & $r$ \\ \hline
T & T & T & 0.0008\\
T & T & F & 0.008\\
T & F & T & 0.08\\
T & F & F & 0.8\\
F & T & T & 0.0002\\
F & T & F & 0.001\\
F & F & T & 0.01\\
F & F & F & 0.1
\end{tabular} 
\end{center} 
The way to read the table is that each row defines a world. The characteristics of the world are just the truth values of three salient propositions: $p, q, r$. If T is shown underneath the name of a proposition, that means the proposition is true in the world, and if an F is shown, it mans the proposition is false. So the third line is the world where $p$ and $r$ are true, and $q$ is false. To the right of the vertical line is the probability of that world. We can work out the probability of any proposition that is true at multiple worlds by adding up the probabilities of the worlds where it is true. So $\Pr(q) = 0.0008 + 0.008 + 0.0002 + 0.001 = 0.01$.

Now let's see what happens when we look at conditional probabilities, i.e., probabilities you get from assuming some particular proposition is true. If we assume that $B$ is true, then we should `zero out', i.e. assign probability 0, to all the possibilities where $B$ doesn't obtain. We're now left with a measure over only the $B$-possibilities. The problem is that it isn't a normalised measure. The values will only sum to $\Pr(B)$, not to 1. We need to renormalise. So we divide by $\Pr(B)$ and we get a probability back. In a formula, we're left with

\begin{equation*}
Pr(A|B) = \frac{\Pr(A \wedge B)}{\Pr(B)}
\end{equation*}
Assume now that we're trying to find the conditional probability of $p$ given $q$ in the above table. We could do this in two different ways. First, we could set the probability of any line where $q$ is false to 0. So we will get the following table.

\begin{center} \begin{tabular}{c c c | c}
$p$ & $q$ & $r$ \\ \hline
T & T & T & 0.0008\\
T & T & F & 0.008\\
T & F & T & 0\\
T & F & F & 0\\
F & T & T & 0.0002\\
F & T & F & 0.001\\
F & F & T & 0\\
F & F & F & 0
\end{tabular} \end{center} 
The numbers don't sum to 1 any more. They sum to 0.01. So we need to divide everything by 0.01. It's sometimes easier to conceptualise this as multiplying by $\nicefrac{1}{\Pr(q)}$, i.e. by multiplying by 100. Then we'll end up with:

\begin{center} \begin{tabular}{c c c | c}
$p$ & $q$ & $r$ \\ \hline
T & T & T & 0.08\\
T & T & F & 0.8\\
T & F & T & 0\\
T & F & F & 0\\
F & T & T & 0.02\\
F & T & F & 0.1\\
F & F & T & 0\\
F & F & F & 0
\end{tabular}
\end{center} 
And since $p$ is true on the top two lines, the `new' probability of $p$ is 0.88. That is, the conditional probability of $p$ given $q$ is 0.88. As we were writing things above, $\Pr(p | q) = 0.88$.

Alternatively we could just use the formula given above. Just adding up rows gives us the following numbers.
\begin{eqnarray*}
\Pr(p \wedge q) &=& 0.0008 + 0.008 = 0.0088 \\
\Pr(q) &=& 0.0008 + 0.008 + 0.0002 + 0.001 = 0.01
\end{eqnarray*} Then we can apply the formula.
\begin{eqnarray*}
\Pr(p | q) &=& \frac{\Pr(p \wedge q)}{\Pr(q)} \\
&=& \frac{0.0088}{0.01} \\
&=& 0.88
\end{eqnarray*}

\section{Bayes Theorem}
This is a bit of a digression from what we're really focussing on in game theory, but it's important for understanding how conditional probabilities work. It is often easier to calculate conditional probabilities in the `inverse' direction to what we are interested in. That is, if we want to know $\Pr(A | B)$, it might be much easier to discover $\Pr(B | A)$. In these cases, we use Bayes Theorem to get the right result. I'll state Bayes Theorem in two distinct ways, then show that the two ways are ultimately equivalent.
\begin{eqnarray*}
\Pr(A|B) &=& \frac{\Pr(B|A) \Pr(A)}{\Pr(B)} \\
&=& \frac{\Pr(B|A) \Pr(A)}{\Pr(B|A)Pr(A) + \Pr(B|\neg A)Pr(\neg A)}
\end{eqnarray*}
These are equivalent because $\Pr(B) = \Pr(B|A)Pr(A) + \Pr(B|\neg A)\Pr(\neg A)$. Since this is an independently interesting result, it's worth going through the proof of it. First note that
\begin{eqnarray*}
\Pr(B|A)\Pr(A) &=& \frac{\Pr(A \wedge B)}{\Pr(A)} \Pr(A) \\
&=& \Pr(A \wedge B) \\
\\
\Pr(B|\neg A)\Pr(\neg A) &=& \frac{\Pr(\neg A \wedge B)}{\Pr\neg (A)} \Pr(\neg A) \\
&=& \Pr(\neg A \wedge B)
\end{eqnarray*} Adding those two together we get
\begin{eqnarray*}
\Pr(B|A)\Pr(A) + \Pr(B|\neg A)\Pr(\neg A) &=& \Pr(A \wedge B) + \Pr(\neg A \wedge B) \\
&=& \Pr((A \wedge B) \vee (\neg A \wedge B)) \\
&=& \Pr(B)
\end{eqnarray*} 
The second line uses the fact that $A \wedge B$ and $\neg A \wedge B$ are inconsistent. And the third line uses the fact that $(A \wedge B) \vee (\neg A \wedge B)$ is equivalent to $A$. So we get a nice result, one that we'll have occasion to use a bit in what follows.
\begin{equation*}
\Pr(B) = \Pr(B|A)\Pr(A) + \Pr(B|\neg A)\Pr(\neg A)
\end{equation*} So the two forms of Bayes Theorem are the same. We'll often find ourselves in a position to use the second form. 

One kind of case where we have occasion to use Bayes Theorem is when we want to know how significant a test finding is. So imagine we're trying to decide whether the patient has disease D, and we're interested in how probable it is that the patient has the disease conditional on them returning a test that's positive for the disease. We also know the following background facts.
\begin{itemize}
\item In the relevant demographic group, 5\% of patients have the disease.
\item When a patient has the disease, the test returns a position result 80\% of the time
\item When a patient does not have the disease, the test returns a negative result 90\% of the time
\end{itemize}
\noindent So in some sense, the test is fairly reliable. It usually returns a positive result when applied to disease carriers. And it usually returns a negative result when applied to non-carriers. But as we'll see when we apply Bayes Theorem, it is very unreliable in another sense. So let $A$ be that the patient has the disease, and $B$ be that the patient returns a positive test. We can use the above data to generate some `prior' probabilities, i.e. probabilities that we use prior to getting information about the test.
\begin{itemize}
\item $\Pr(A) = 0.05$, and hence $\Pr(\neg A) = 0.95$
\item $\Pr(B | A) = 0.8$
\item $\Pr(B | \neg A) = 0.1$
\end{itemize}
\noindent Now we can apply Bayes theorem in its second form.
\begin{eqnarray*}
\Pr(A|B) &=& \frac{\Pr(B|A) \Pr(A)}{\Pr(B|A)\Pr(A) + \Pr(B|\neg A)\Pr(\neg A)} \\
&=& \frac{0.8 \times 0.05}{0.08 \times 0.05 + 0.1 \times 0.95} \\
&=& \frac{0.04}{0.04 + 0.095} \\
&=& \frac{0.04}{0.135} \\
&\approx& 0.296
\end{eqnarray*} So in fact the probability of having the disease, conditional on having a positive test, is less than 0.3. So in that sense the test is quite unreliable.

This is actually a quite important point. The fact that the probability of $B$ given $A$ is quite high does not mean that the probability of $A$ given $B$ is equally high. By tweaking the percentages in the example I gave you, you can come up with cases where the probability of $B$ given $A$ is arbitrarily high, even 1, while the probability of $A$ given $B$ is arbitrarily low.

Confusing these two conditional probabilities is sometimes referred to as the \textit{Prosecutors' fallacy}, though it's not clear how many actual prosecutors are guilty of it. The thought is that some prosecutors start with the premise that the probability of the defendant's blood (or DNA or whatever) matching the blood at the crime scene, conditional on the defendant being innocent, is 1 in a billion (or whatever it exactly is). They conclude that the probability of the defendant being innocent, conditional on their blood matching the crime scene, is about 1 in a billion. Because of derivations like the one we just saw, that is a clearly invalid move.

A more striking real world example comes from the analysis of soccer carried out by Charles Reep and Charles Hughes. (Reep and Hughes worked separately, but they reached similar conclusions by similar means.) They watched many games, and tracked each of the possessions (meaning a period of time where one team has control of the ball). They counted how many passes took place in each possession, and whether the possession resulted in a goal. They noted that the vast majority of the goals came from very short possessions, ones involving three or fewer passes. From this they concluded that soccer teams should be direct; kicking the ball to forwards as soon as possible to try to score from a short possession. Hughes became technical director of the Football Association in England, where he spread this idea of how to play soccer across the country.

You may be able to spot the error in their reasoning already, but if not let's work through a fictionalised version to see the problem. (The numbers I'm using are made up, though I think the scale roughly matches what Reep and Hughes found.) Divide all possessions into short (1-3 passes), medium (4-7 passes) and long (8+ passes). Imagine that we observed a bunch of games, and recorded the following data.

\begin{center}
\begin{tabular}{r | c c}
& Goal & No Goal \\ \hline
Short & 78 & 822 \\ %900 
Medium & 12 & 48 	\\ %60
Long & 10 & 30  %40 
\end{tabular}
\end{center}
As you can see, most goals come from short possessions. But you shouldn't get excited by finding your team has had a short possession; over 90\% of short possessions go nowhere. On the other hand, medium to long possessions often end in goals. Let's add another column to make this clear.

\begin{center}
\begin{tabular}{r | c c | c}
& Goal & No Goal & Success \% \\ \hline
Short & 78 & 822 & 8.7 \\ %900 
Medium & 12 & 48 &	20 \\ %60
Long & 10 & 30 & 25 %40 
\end{tabular}
\end{center}
The longer a possession is, the higher the frequency that it ends in a goal. I'm not sure in real soccer it is as dramatic as the numbers I gave you. But I am sure that (at least at high levels, at most times in the history of soccer) the broad pattern I've outlined here is correct. Most goals come from short possessions, but short possessions lead to fewer goals, on average, than long possessions. Anyone who has taken a probability class should know that this suggests aiming for longer possessions rather than shorter ones, but apparently the England hierarchy didn't take that lesson for several decades.

\section{Expected Value}
One very convenient notion to define using a probability function is that of an \textbf{expected value}. Say that a \textbf{random variable} is any function from possibilities to real numbers. We usually designate random variables with capital letters. So expressions like $X=x$, which you'll often see in probability textbooks, means that the random variable $X$ takes real value $x$. 

The expected value of a random variable is calculated using the following formula.

\begin{equation*}
E(X) = \sum_w \Pr(w)X(w)
\end{equation*}

\noindent Remember that $w$ is a variable that ranges over atoms. So the summation here is over the atoms in the field. $\Pr(w)$ is the probability of that atom obtaining, and $X(w)$ is the value of the random variable at $w$. Here's a worked example to make it clearer.

Let our random variable $X$ be the number of children of the next person to walk through the door. Let's assume that the department contains 10 people; 2 of whom have three children, 3 of whom have two children, 3 of whom have one child, and 2 are childless. Each member of the department is equally likely to walk through the door, and no one else is going to. So the probability of any given member of the department walking through the door is $\nicefrac{1}{10}$. Then the expected value of $X$, i.e., $E(X)$, is:

\begin{align*}
\nicefrac{1}{10} \times 3 + \nicefrac{1}{10} \times 3 &+ \\
\nicefrac{1}{10} \times 2 + \nicefrac{1}{10} \times 2 + \nicefrac{1}{10} \times 2 &+ \\
\nicefrac{1}{10} \times 1 + \nicefrac{1}{10} \times 1 + \nicefrac{1}{10} \times 1 &+ \\
\nicefrac{1}{10} \times 0 + \nicefrac{1}{10} \times 0 &= 1.5 \\
\end{align*}

\noindent Note that the expected value of $X$ is \textit{not} the value we `expect', in the ordinary sense, $X$ to take. We can be pretty sure that $X$'s value will be an integer, so it won't be 1.5, for instance! That is to say, `expected value' is a term of art.

As well as `absolute' expected values, we can also define conditional expected values, as follows:

\begin{equation*}
E(X | p) = \sum_w \Pr(w | p)X(w)
\end{equation*}

\noindent I won't go through the proof of this, but it turns out that conditional expected values satisfy a nice theorem. Let $\{p_1, \dots, p_n\}$ be a partition of possibility space. Then the following holds.

\begin{equation*}
E(X) = \sum_{i=1}^n \Pr(p_i) E(X | p_i)
\end{equation*}

\noindent Here are a few of immediate corrollaries of that result.

\begin{align*}
E(X) -E(Y) &= E(X - Y) \\
&= \sum_{i=1}^n \Pr(p_i) (E(X | p_i) - E(Y | p_i)) \\
&= \sum_{i=1}^n \Pr(p_i) (E((X - Y) | p_i)) 
\end{align*}

\noindent Often we know more about the difference between two expected values than we do about the absolute value of either. For instance, if we are comparing the expected values of taking and declining a small bet, we can figure out which action has a higher expected value without figuring out the expected value of either, which presumably depends on all sorts of facts about how our life might (probably) go. But as long as we can work out the difference in expected value between the two acts conditional on each $p_i$, and the probability of each $p_i$, we can work out at least which has a higher expected value, and by how much.

\section{Orthodox Decision Theory}
Given all that, it is fairly easy to summarise orthodox decision theory.

\begin{itemize}
\item Rational agents have a credence function over salient alternatives that is identical to some probability function $\Pr$.
\item Rational agents also have a real-valued utility function $U$ over salient outcomes.
\item The rational action for such an agent to do is the action that maximises \textbf{expected} utility, i.e., maximises the expected value of $U$ relative to that function $\Pr$.
\end{itemize}
Next time we'll start with an application of those ideas that leads to the most central concept in game theory.

%%%%
%
%
%
%%%%%%%%%%%%%%%%%%%%%%%%%%%%%%%%%%
% Start of Chapter 3
%%%%%%%%%%%%%%%%%%%%%%%%%%%%%%%%%%
%
%
%
%%%%%


\chapter{Mixed Strategies and Nash Equilibrium}
\section{Mixed Strategies}

In any game, if player $x$ has strategies $s_1, s_2$ available, then we'll also assume they have available the \textbf{mixed strategy} of playing $s_1$ with probability $p$ and $s_2$ with probability $1-p$ (for arbitrary $p \in (0, 1)$). Such strategies have an expected value given strategies for the other players in the game, and we'll think of that expected value as the return to the player for playing the strategy.

Mixed strategies are contrasted with \textbf{pure strategies}. These are the strategies that are usually listed on the game table. They are the only strategies that we have considered so far.

I'm not sure whether there is a standard terminology for denoting mixed strategies. We'll write $\langle s_1, p; s_2, 1-p \rangle$ for the mixed strategy of playing $s_1$ with probability $p$ and $s_2$ with probability $1-p$. It is possible for mixed strategies to involve more than two pure strategies; in those cases the list of pure strategies and their probabilities will have to be longer.

If the agent has more than two pure strategies available to them, and plays $\langle s_1, p; s_2, 1-p \rangle$, it wouldn't make any difference if they instead played $\langle s_1, p; s_2, 1-p; s_3, 0; \dots; s_n, 0\rangle$. We can treat any mixed strategy as a mixture of all the pure strategies, just that some of the strategies have probability 0 of being played. Indeed, we can treat any strategy as a mixed strategy that perhaps gives some pure strategies probability 0. So we no longer need to talk about pure strategies at all; they are just a limit case of mixed strategies where one pure strategy gets probability 1, and the others all get probability zero.

Still, even if every mixed strategy can be represented as a mixture of all pure strategies, there is an important distinction between the pure strategies that get positive probability, and those that get zero probability. We'll call the set of pure strategies that get positive probability the \textbf{support} of the mixed strategy. What we used to call a pure strategy is a mixed strategy whose support is a singleton.

Now it's a somewhat delicate philosophical question just what this assumption amounts to. While the assumption is more or less universally made by game theorists, there is less agreement about just what it comes to. Here are three natural interpretations:

\begin{itemize}
\item To play the mixed strategy $\langle s_1, p; s_2, 1-p \rangle$ is to use a randomizing device to pick between $s_1$ and $s_2$, with the device set to choose $s_1$ with probability $p$. Call this the \textbf{objective chance} interpretation, since it says that the agent playing that strategy has some defined chance of playing each pure strategy.
\item To say that someone is playing the mixed strategy $\langle s_1, p; s_2, 1-p \rangle$ is to say our evidence doesn't settle what they will play, but it points to them playing either $s_1$ or $s_2$, the first with probability $p$, the second with probability $1-p$. Our evidence could point this way because it is settled that the agent is using a randomizing device, but it could also be because we just don't have enough information to settle what they will play. Call this the \textbf{epistemic} interpretation of a mixed strategy, since it locates a mixed strategy in our level of knowledge about what the agent will do.
\item To say that  someone is playing the mixed strategy $\langle s_1, p; s_2, 1-p \rangle$ is to say that in the long run they will play $s_1$ with frequency $p$, and $s_2$ with frequency $1-p$. Call this, for obvious reasons, the \textbf{frequentist} interpretation of mixed strategies. On this interpretation, there is no such thing as playing a mixed strategy in a one-off game; it only makes sense to think about these strategies in repeated games. In practice, we'll usually deploy mixed strategies in repeated games, so this isn't as big a restriction as it might sound.
\end{itemize}

\section{Best Responses}

A strategy $s$ is a \textbf{best response} for player $x$ iff there is some probability distribution $\Pr$ over the possible strategies of other players such that no strategy has a higher expected payoff than $s$ does, given $\Pr$.

Let's look at an example of this. So far we've used the same names for the options that Row and Column have available. But we'll start to look at asymmetric games, so it helps to use asymmetric terminology. We'll use $U$, for Up, $M$ for Middle, and $D$, for Down, for Row's options, and $L$ for Left, and $R$ for Right, for Column's options.

\begin{center}
\begin{tabular}{l r | c c}
& & \multicolumn{2}{c}{\textbf{Col}} \\
& & $L$ & $R$ \\ \hline
\multirow{3}{*}{\textbf{Row}} & $U$ & 3, 0 & 0, 0 \\
& $M$ & 2, 0 & 2, 0 \\
& $D$ & 0, 0 & 3, 0
\end{tabular}
\end{center}
We will just look at things from the perspective of Row. What we want to show is that all three of the possible moves here are best responses. We will write $E(X) = x$ to mean that the expected value of strategy $X$ is $x$.

It is clear that $U$ is a best response. Set $\Pr(L) = 1, \Pr(R) = 0$. Then $E(U) = 3, E(M) = 2, E(D) = 0$, so $U$ is best. It is also clear that $D$ is a best response. Set $\Pr(L) = 0, \Pr(R) = 1$. Then $E(U) = 0, E(M) = 2, E(D) = 3$, so $D$ is best.

The striking thing is that $M$ can also be a best response. Set $\Pr(L) = \Pr(R) = 0.5$. Then $E(U) = E(D) = 1.5$. But $E(M) = 2$, which is greater than 1.5. So if Row thinks it is equally likely that Col will play either $L$ or $R$, then she maximizes expected utility by playing $M$. Of course, she doesn't maximize actual utility. Maximizing actual utility requires making a gamble on which choice $C$ will make. That isn't always wise; it might be best to take the safe option.

It's very plausible that agents should play best responses. If a move is not a best response, then there is no way that it maximizes expected utility, no matter what one's views of the other players are. And one should maximize expected utility. 

Any strongly dominated option will never be a best response. That's because whatever strategy strongly dominates it will have a higher expected return no matter what the probabilities are. But a weakly dominated response can be a best response. That's because we've defined best responses to allow that something can be a best response when something else ties it. So in the following example, $D$ is weakly dominated by $U$, but it is a best response given $\Pr(R) = 1$.

\begin{center}
\begin{tabular}{l r | c c}
& & \multicolumn{2}{c}{\textbf{Col}} \\
& & $L$ & $R$ \\ \hline
\multirow{2}{*}{\textbf{Row}} & $U$ & 1, 0 & 0, 0 \\
& $D$ & 0, 0 & 0, 0
\end{tabular}
\end{center}
What about the other direction? If something is not dominated, does that mean it is a best response? It turns out the answer to this question is a little more complicated. Start with this game.

%Simple Worst of Both Worlds game
\begin{center}
\begin{tabular}{l r | c c}
& & \multicolumn{2}{c}{\textbf{Col}} \\
& & $L$ & $R$ \\ \hline
\multirow{3}{*}{\textbf{Row}} & $U$ & 3, 0 & 0, 0 \\
& $M$ & 1, 0 & 1, 0 \\
& $D$ & 0, 0 & 3, 0
\end{tabular}
\end{center}
In this game, $M$ is not dominated by either $U$ or $D$. After all, if Col plays $L$, then $M$ beats $D$. And if Col plays $R$ then $M$ beats $D$.

On the other hand, $M$ is not a best response. The possible views Row could have about Col are all characterised by a single variable $x$, where $x$ is the probability (according to Row) of Col playing $L$. This variable $x$ can take any value in $[0, 1]$. But whatever value it takes, $M$ is not the best option. Note that the expected return of $L$ is $3x$, and the expected return of $R$ is $3(1-x)$. From that, we can easily compute that:

\begin{itemize}
\item If $0 \leq x < 0.5$ then $D$ is the best response.
\item If $x = 0.5$ then $U$ and $D$ are tied for best response.
\item If $0.5 < x \leq 1$ then $U$ is the best response.
\end{itemize}
So in no circumstance is $M$ best. It is not a best response.

But on second thought, perhaps it is dominated. Let's add a row to the table for the mixed strategy $\langle U, 0.5; D, 0.5 \rangle$. We'll add this to the table, and note in the table its expected value.

%Simple Worst of Both Worlds game with mix
\begin{center}
\begin{tabular}{l r | c c}
& & \multicolumn{2}{c}{\textbf{Col}} \\
& & $L$ & $R$ \\ \hline
\multirow{4}{*}{\textbf{Row}} & $U$ & 3, 0 & 0, 0 \\
& $M$ & 1, 0 & 1, 0 \\
& $D$ & 0, 0 & 3, 0 \\
& $\langle U, 0.5; D, 0.5 \rangle$ & 1.5, 0 & 1.5, 0
\end{tabular}
\end{center}
And note that $\langle U, 0.5; D, 0.5 \rangle$ dominates $M$; whatever Col plays, that mixed strategy has a higher expected return than $M$.

If we change the payout to $M$ so that it isn't dominated, it becomes a best response. So recall the version of the game where the payout from $M$ from a guaranteed 1 to a guaranteed 2. I've reprinted that here with the mixed strategy made explicit.

%Simple Best of Both Worlds game
\begin{center}
\begin{tabular}{l r | c c}
& & \multicolumn{2}{c}{\textbf{Col}} \\
& & $L$ & $R$ \\ \hline
\multirow{4}{*}{\textbf{Row}} & $U$ & 3, 0 & 0, 0 \\
& $M$ & 2, 0 & 2, 0 \\
& $D$ & 0, 0 & 3, 0 \\
& $\langle U, 0.5; D, 0.5 \rangle$ & 1.5, 0 & 1.5, 0
\end{tabular}
\end{center}
Now $M$ is not dominated. Indeed, it dominates $\langle U, 0.5; D, 0.5 \rangle$. And as we saw above it is also a best response. If Row thinks it is 50/50 whether Col plays $L$ or $R$, then $M$ has the highest expected return.

\section{Iterated Best Responses}

It seems very plausible that one should only play best responses. But not all best responses are worth playing. Some best responses are pretty hard to justify when playing against a rational opponent. Consider the following game from Row's perspective.

\begin{center}
\begin{tabular}{l r | c c}
& & \multicolumn{2}{c}{\textbf{Col}} \\
& & $L$ & $R$ \\ \hline
\multirow{2}{*}{\textbf{Row}} & $U$ &  5, 5 & 0, -5\\
& $D$ & 0, 5 & 2, -5\\
\end{tabular}
\end{center}
In this game, $D$ is a best response. It does best if Col chooses $R$. But why would Col do that? Col gets 5 for sure if she chooses $L$, and -5 for sure if she chooses $R$. Assuming the weakest possible rationality constraint on Col, she won't choose a sure loss over a sure gain. So given that, Row should choose $U$.

Now we could have shown that with just considering domination; indeed domination by pure strategies.The following example is a little more complex.

\begin{center}
\begin{tabular}{l r | c c c}
& & \multicolumn{3}{c}{\textbf{Col}} \\
& & $L$ & $C$ & $R$ \\ \hline
\multirow{2}{*}{\textbf{Row}} & $U$ & 1, 3 & 0, 1 & 1, 0\\
& $D$ & 0, 0 & 1, 1 & 0, 3
\end{tabular}
\end{center}
In this game neither player has a dominated move, as long as we restrict ourselves to pure strategies. And, from Row's perspective, thinking about best responses doesn't help either. $U$ is a best response if Col is going to play $L$ or $R$ with probability at least 0.5, and $D$ is best response if Col is going to play $m$ with probability at least 0.5.

But note that $C$ is not a best response for Col. One of the outer options has a higher expected value, whatever probability she gives to Row playing $U$. Now let's assume, when thinking from Row's perspective, that Col will play a best response. That is, we're assuming Col will play either $L$ or $R$. Given that, the best thing for Row to do is to play $U$.

More carefully, if Row knows that Col is rational, and if Row knows that rational agents always play best responses, then Row has a compelling reason to play $U$. This suggests an important normative status. Call a \textbf{BRBR} strategy one that is a \textbf{B}est \textbf{R}esponse to a \textbf{B}est \textbf{R}esponse.

But we can go further than that. What should Col play, given what we just concluded about Row? She should play $L$. After all, she knows that Row will play $U$, and given that $L$ has the best return.

The idea that players should play BRBR strategies isn't strong enough to get us that result. After all, $D$ is a best response, and $R$ is a best response to $D$. So $R$ is BRBR. What it isn't is a best response to something that is a best response to a best response. Rather than keep counting layers of best responses, let's collapse this into one notion of being a best response to a best response to a best response to etc., to infinity. Call any strategy which has this feature an \textbf{IBR}, or \textbf{I}terated \textbf{B}est \textbf{R}esponse.

\section{Nash Equilibrium}
In a two-player game, for a strategy $s_0$ to be IBR, it must be the best response to some strategy $s_1$, which is the best response to some strategy $s_2$, which is the best response to some strategy $s_3$, etc. But we are assuming that there are only finitely many strategy choices available. So how can we get such an infinite chain going?

The answer, of course, is to repeat ourselves. As long as we get some kind of loop, we can extend such a chain forever, by keeping on circling around the loop. And the simplest loop will have just two steps in it. So consider any pair of strategies $\langle s_0, s_1 \rangle$ such that $s_0$ is the best response to $s_1$, and $s_1$ is the best response to $s_0$. In that case, each strategy will be BRBRI, since we can run around that two-step `loop' forever, with each stop on the loop being a strategy which is a best response to the strategy we previous stopped at.

When a pair of strategies fit together nicely like this, we say they are a \textbf{Nash Equilibrium}. More generally, in an $n$-player game, we use the following definition.

\begin{description}
\item[NE] Some strategies $s_1, ..., s_n$ in an $n$-player game form a \textbf{N}ash \textbf{E}quilibrium iff for each $i$, $s_i$ is a best response to the strategies played by the other $n-1$ players. That is, iff player $i$ cannot do better by playing any alternative strategy to $s_i$, given that the other players are playing what they actually do.
\end{description}

\noindent The `Nash' in Nash \Eqm\ is in honour of John Nash, who developed the concept, and proved some striking mathematical results concerning it.

\section{Nash Equilibrium in Simultaneous Move Games}

Consider the following game:

\begin{center}
\begin{tabular}{l r | c c}
& & \multicolumn{2}{c}{\textbf{Col}} \\
& & $L$ & $R$ \\ \hline
\multirow{2}{*}{\textbf{Row}} & $U$ & 2, 0 & 0, 2 \\
& $D$ & 0, 4 & 1, 0
\end{tabular}
\end{center}
Any of the pure strategies is a best response. Indeed, any of them are iterated best responses. So is all we can say about the game that anything goes? Most game theorists disagree with that claim. They say that once we consider mixed strategies, there is a distinctive solution to the game. Consider the following pair of strategies:

\begin{itemize}
\item Row plays $\langle U, \nicefrac{2}{3}; D, \nicefrac{1}{3} \rangle$
\item Col plays $\langle L, \nicefrac{1}{3}; R, \nicefrac{2}{3} \rangle$
\end{itemize}
Given that Row has that strategy, let's work out the expected return (to Col) of Col's two possible plays.

\begin{align*}
E(L) &= 0 \times \frac{2}{3} + 4 \times \frac{1}{3} \\
&= 0 + \frac{4}{3} \\
&= \frac{4}{3} \\
E(R) &= 2 \times \frac{2}{3} + 0 \times \frac{1}{3} \\
&= \frac{4}{3} + 0 \\
&= \frac{4}{3} 
\end{align*}
And we can do the same calculation for the expected return (to Row) of Row's possible strategies.

\begin{align*}
E(U) &= 2 \times \frac{1}{3} + 0 \times \frac{2}{3} \\
&= \frac{2}{3} + 0\\
&= \frac{2}{3} \\
E(D) &= 0 \times \frac{1}{3} + 1 \times \frac{2}{3} \\
&= 0 + \frac{2}{3} \\
&= \frac{2}{3} 
\end{align*}
Both players would be, if they knew what the other player was doing, indifferent between the two possible strategies that are available to them.

The expected return of Col's playing the mixed strategy $\langle L, x; R, 1-x \rangle$ is $xE(L) + (1-x)E(R)$. So if $E(L) = E(R) = \nicefrac{4}{3}$, then for all $x \in [0, 1]$, the expected return of $\langle L, x; R, 1-x \rangle$ is also $\nicefrac{4}{3}$. So all strategies are equally good in response to Row playing $\langle U, \nicefrac{2}{3}; D, \nicefrac{1}{3} \rangle$. That is, they are all best responses. By similar reasoning, all strategies Row could play in response to Col playing $\langle L, \nicefrac{1}{3}; R, \nicefrac{2}{3} \rangle$ are equally good, so they are all best responses.

So the two strategies $\langle U, \nicefrac{2}{3}; D, \nicefrac{1}{3} \rangle$ and $\langle L, \nicefrac{1}{3}; R, \nicefrac{2}{3} \rangle$ are best responses to each other; they form a Nash Equilibrium. Indeed, this is the only Nash Equilibrium of the game.

We aren't going to prove this, but for any game with bounded payouts, and finite options for each player, there will be some Nash Equilibrium. This equilibria will often be a mixed strategy equilibria, but it is guaranteed to exist.

\section{Finding Nash Equilibria}

As the textbook notes, it isn't easy in general to find Nash Equilibria. But let's start with some general approaches we can use. 

\subsection{Iterated Dominance}

Strongly dominated strategies can't be part of an equilibrium. (I'll just use 'equilibrium' for 'Nash Equilibrium' in this section, as is common in game theory texts.) So strategies that only make sense if they other player plays a dominated strategy can't be part of an equilibrium either. This idea can help us find the equilibrium in some games.

\begin{center}
\begin{tabular}{l r | c c}
& & \multicolumn{2}{c}{\textbf{Col}} \\
& & $L$ & $R$ \\ \hline
\multirow{2}{*}{\textbf{Row}} & $U$ & 1, 2 & 5, 0\\
& $D$ & 0, 0 & 3, 6
\end{tabular}
\end{center}
Neither of Col's strategies is dominated. Col wants to choose whatever Row chooses. But look at the game from Row's perspective. Here there is a strongly dominating strategy, namely choosing $U$. So $D$ cannot be part of any equilbria. Indeed, if we assume that Row is rational, and so won't play strongly dominated strategies, it is really as if the strategy isn't even there. From that perspective, the table really looks like this.

\begin{center}
\begin{tabular}{l r | c c}
& & \multicolumn{2}{c}{\textbf{Col}} \\
& & $L$ & $R$ \\ \hline
\multirow{1}{*}{\textbf{Row}} & $U$ & 1, 2 & 5, 0
\end{tabular}
\end{center}
Now in this game there is a dominating strategy for Col, namely $L$. So the only equilibrium of the game is $\langle U, L \rangle$. And it seems that those are the only rational strategies for the players, assuming they each know the other to be rational. 

We can use the idea of removing dominating strategies to illustrate some puzzling features of a couple of other games. I won't do tables for these games, because they'd be much too big. The first is a location game that has many applications.

\begin{quote}
Two trucks have to choose where they will sell ice-cream on a particular beach. There are 11 locations to choose from, which we'll number 0, 1, \dots, 9, 10. Spot 0 is at the left end of the beach, Spot 10 is at the right end of the beach, and the other spots are equally spaced in between. There are 10 people at each location. Each of them will buy ice-cream. If one truck is closer, they will buy ice-cream from that truck. If two trucks are equally close, then 5 of them will buy ice-cream from one truck, and 5 from the other. Each truck aims to maximise the amount of ice-cream it sells. Where should the trucks end up?
\end{quote}

\noindent Let's start by looking at a fragment of the payoff matrix. The payoffs are numbers of ice-creams sold. We'll call the players Row and Col, as usual, and just use the number $n$ for the strategy of choosing location $n$.

\begin{center}
\begin{tabular}{l r | c c c c c c}
& & \multicolumn{6}{c}{\textbf{Col}} \\
 & & 0 & 1 & 2 & 3 & 4 & \dots  \\ \hline
\multirow{2}{*}{\textbf{Row}}  & 0 & 55, 55 & 10, 100 & 15, 95 & 20, 90 & 25, 85 & \dots \\
& 1 & 100, 10 & 55, 55 & 20, 90 & 25, 95 & 30, 80 & \dots \\
\dots 
\end{tabular}
\end{center}
I'll leave it as an exercise to confirm that these numbers are indeed correct. Note that no matter what Col selects, Row is better off picking 1 than 0. If Col picks 0 or 1, she is a lot better off; she sells 45 more ice-creams. And if Col picks a higher number, she is a bit better off; she sells 5 more ice-creams. So picking 1 dominates picking 0. 

Let's look at the opposite corner of the matrix.

\begin{center}
\begin{tabular}{l r | c c c c c c}
& & \multicolumn{6}{c}{\textbf{Col}} \\
& & \dots & 6 & 7 & 8 & 9 & 10  \\ \hline
\dots \\
\multirow{2}{*}{\textbf{Row}} & 9 & \dots & 30, 80 & 25, 85 & 20, 90 & 55, 55 & 100, 10 \\
& 10 & \dots & 25, 85 & 20, 80 & 15, 95 & 10, 100 & 55, 55 \\
\end{tabular}
\end{center}
Again, there should be a pattern. No matter what Col does, Row is better off picking 9 than 10. In most cases, this leads to selling 5 more ice-creams. If Col also picks 9 or 10, then picking 9 gives Row a big advantage. The argument here was obviously symmetric to the argument about picking 0 or picking 1, so I'll stop concentrating on what happens when both players select high numbers, and focus from here on the low numbers.

So there is a clear conclusion to be drawn from what we've said so far.

\begin{itemize}
\item Spot 0 is dominated by Spot 1, and Spot 10 is dominated by Spot 9. So if Row is rational, she won't pick either of those spots. Since the game is symmetric, if Col is rational, she won't pick either of those spots either. Moreover, neither 0 nor 10 can be part of any equilibrium.
\end{itemize}

\noindent Now let's turn to the comparison between Spot 1 and Spot 2. Again, we'll just look at a fragment of the matrix.

\begin{center}
\begin{tabular}{l r | c c c c c c c}
& & \multicolumn{7}{c}{\textbf{Col}} \\
& & 0 & 1 & 2 & 3 & 4 & 5 & \dots  \\ \hline
\dots \\
\multirow{2}{*}{\textbf{Row}} & 1 & 100, 10 & 55, 55 & 20, 90 & 25, 95 & 30, 80 & 35, 75 & \dots \\
& 2 & 95, 15 & 90, 20 & 55, 55 & 30, 80 & 35, 75 & 40, 70 &  \dots \\
\dots 
\end{tabular}
\end{center}
The pattern is also clear. If Col selects any number above 2, then Row sells 5 more ice-creams by picking 2 rather than 1. If Col selects either 1 or 2, then Row sells 35 more ice-creams by picking 2 rather than 1. But if Col selects 0, then Row sells 5 \textit{fewer} ice-creams by picking 2 rather than 1. So picking 2 does \textit{not} dominate picking 1.

But note the only circumstance when picking 2 is worse than picking 1 is if Col picks 0. And picking 0 is, for Col, a strongly dominated strategy. So picking 2 is sure to do better than picking 1 if Col does not play a strongly dominated strategy. So let's delete that dominated strategy. Here's what the top left corner of the game matrix looks like when we do that.

\begin{center}
\begin{tabular}{l r | c c c c c c}
& & \multicolumn{6}{c}{\textbf{Col}} \\
& & 1 & 2 & 3 & 4 & 5 & \dots  \\ \hline
\multirow{2}{*}{\textbf{Row}} & 1 & 55, 55 & 20, 90 & 25, 95 & 30, 80 & 35, 75 & \dots \\
& 2 & 90, 20 & 55, 55 & 30, 80 & 35, 75 & 40, 70 &  \dots \\
\dots 
\end{tabular}
\end{center}
And now it looks like picking 1 is dominated. This teaches us another lesson about the game.

\begin{itemize}
\item If we `delete' the option of picking either 0 or 10 for Col, then picking 2 dominates picking 1 for Row. In other words, if Row knows Col is rational, and hence won't pick a dominated option, then picking 2 dominates picking 1 for Row, relative to the space of epistemically possible moves in the game. For similar reasons, if Row knows Col is rational, then picking 8 dominates picking 9. And if Col knows Row is rational, then picking either 1 or 9 is dominated (by 2 and 8 respectively). And for similar reasons, neither 1 nor 9 can be part of an equilibrium strategy.
\end{itemize}
Summarising what we know so far,

\begin{itemize}
\item If the players are rational, they won't pick 0 or 10.
\item If the players know the other player is rational, they also won't pick 1 or 9.
\end{itemize}
Let's continue down the matrix. For simplicity, we'll leave off columns 0 and 10, since they are dominated, and we have deleted those as possible options.

\begin{center}
\begin{tabular}{l r | c c c c c c}
& & \multicolumn{6}{c}{\textbf{Col}} \\
& & 1 & 2 & 3 & 4 & 5 & \dots  \\ \hline
\multirow{2}{*}{\textbf{Row}} & 2 & 90, 20 & 55, 55 & 30, 80 & 35, 75 & 40, 70 &  \dots \\
& 3 & 85, 25 & 80, 30 & 55, 55 & 40, 70 & 45, 65 & \dots \\
\dots 
\end{tabular}
\end{center}
Picking 3 doesn't \textit{quite} dominate picking 2. In most circumstances, Row does better by picking 3 rather than 2. But she does a little worse if Col picks 1. But wait! We just had an argument that Col shouldn't pick 1. Or, at least, if Col knows that Row is rational, she shouldn't pick 1. Let's assume that Col does know that Row is rational, and Row in turn knows that fact, so she can use it in reasoning. That means she knows Col won't pick 1. So she can delete it from consideration too. And once she does, picking 3 dominates picking 2, relative to the reduced space of epistemically possible outcomes.

I haven't belaboured the point as much as in the previous paragraphs, but hopefully the following conclusion is clear.

\begin{itemize}
\item If the players know that the players know that the other player is rational, they won't pick 2 or 8. So neither 2 nor 8 can be part of an equilibrium.
\end{itemize}
And you can see where this is going. Once we rule out each player picking 2 or 8, then picking 4 dominates picking 3, so picking 3 should be ruled out. And picking 7 should be ruled out for symmetric reasons. But once 3 and 7 are ruled out, picking 5 dominates picking either 4 or 6. So both players will end up picking 5.

And that is, in the standard economic textbooks, a nice explanation of why we see so much `clustering' of shops. (For instance, why so often there are several gas stations on one corner, rather than spread over town.) Of course the full explanation of clustering is more complicated, but it is nice to see such a simple model deliver such a strong outcome. 

This process is called the \textbf{Iterative Deletion of Dominated Strategies}. More precisely, what we've used here is the strategy of iteratively deleting \textbf{strongly} dominated strategies. This is a powerful technique for solving games that don't, at first glance, admit of any easy solution. If only one strategy pair remains after using this technique, that is the only equilibrium of the game.

We'll return a bit to this point in what follows, but it's worth noting that in a one-shot version of this game, to get to the conclusion that each player should play 5, we had to use a number of very strong assumptions. We had to assume that each player is rational, that they know the other is rational, that they know that the other knows they are rational, and so on through six iterations. In principle, to justify iteratively deleting dominated strategies, we have to assume that there are arbitrarily many iterations of this idea that each player knows that the other player knows etc that they are rational.

\subsection{Finding Best Responses}

It can be useful to mark explicitly what the best responses are to each possible play by the other player. In a two-player game, that can sometimes lead to turning up Nash equilibria. Even when it does not, it can give you a hint about what the support for a mixed strategy equilibria could be. To see an example of this, consider the following game.

\begin{center}
\begin{tabular}{l r | c c c c c}
& & \multicolumn{5}{c}{\textbf{Col}} \\
& & $a$ & $b$ & $c$ & $d$ & $e$ \\ \hline
\multirow{5}{*}{\textbf{Row}}
& $A$ & 9,9 & 5,1 & 6,4 & 3,3 & 1,8\\
& $B$ & 3,1 & 7,3 & 1,7 & 3,9 & 8,2 \\
& $C$ & 4,5 & 8,7 & 2,4 & 4,1 & 2,9\\
& $D$ & 3,2 & 1,5 & 3,7 & 7,7 & 9,1\\
& $E$ & 5,3 & 8,3 & 7,6 & 3,5 & 8,4
\end{tabular}
\end{center}
We'll start by finding the best response for Row to each of Col's plays. We'll underline the payoff for Row that does best, assuming Col makes that choice. Here's how the table looks after we do that.

\begin{center}
\begin{tabular}{l r | c c c c c}
& & \multicolumn{5}{c}{\textbf{Col}} \\
& & $a$ & $b$ & $c$ & $d$ & $e$ \\ \hline
\multirow{5}{*}{\textbf{Row}}
& $A$ & \underline{9},9 & 5,1 & 6,4 & 3,3 & 1,8\\
& $B$ & 3,1 & 7,3 & 1,7 & 3,9 & 8,2 \\
& $C$ & 4,5 & \underline{8},7 & 2,4 & 4,1 & 2,9\\
& $D$ & 3,2 & 1,5 & 3,7 & \underline{7},7 & \underline{9},1\\
& $E$ & 5,3 & \underline{8},3 & \underline{7},6 & 3,5 & 8,4
\end{tabular}
\end{center}
Note that under $b$, we underlined two possible payoffs. That's because they are tied for best, and best response is compatible with ties. Now we'll go through Col's options, finding the best response to each possible play by Row.
\begin{center}
\begin{tabular}{l r | c c c c c}
& & \multicolumn{5}{c}{\textbf{Col}} \\
& & $a$ & $b$ & $c$ & $d$ & $e$ \\ \hline
\multirow{5}{*}{\textbf{Row}}
& $A$ & \underline{9},\underline{9} & 5,1 & 6,4 & 3,3 & 1,8\\
& $B$ & 3,1 & 7,3 & 1,7 & 3,\underline{9} & 8,2 \\
& $C$ & 4,5 & \underline{8},7 & 2,4 & 4,1 & 2,\underline{9}\\
& $D$ & 3,2 & 1,5 & 3,\underline{7} & \underline{7},\underline{7} & \underline{9},1\\
& $E$ & 5,3 & \underline{8},3 & \underline{7},\underline{6} & 3,5 & 8,4
\end{tabular}
\end{center}
Any cell where both values are underlined is an equilibrium; neither player can improve their lot by unilaterally changing their play. And we can see there are three pure strategy equilibria here: $\langle Aa \rangle$, $\langle Dd \rangle$, $\langle Ec \rangle$. I'll leave finding the mixed strategy equilibria as an exercise for the interested reader.
% Just go over underlining

\subsection{Finding Mixed Strategy Equilibria}

It isn't easy in general to find mixed strategy equilibria. But we can start by dividing the problem into two parts.

\begin{enumerate*}
\item Finding the support for each player's part of the equilibria.
\item Finding the probability with which they play each part of their support.
\end{enumerate*}
And it turns out, a bit surprisingly, that it is the first part of the problem that is harder in general. (Recall that the support of a mixed strategy is the set of pure strategies that have a positive probability of being played given that mixed strategy.) And while that isn't an easy problem to solve in general, it isn't something we need to worry about in very simple games. In games with just two options per player, then either there will be a pure strategy equilibria, or the mixed strategy Nash Equilibria will have all strategies in its support. Even if there are three strategies available per player, then it usually isn't too hard to find the support for the mixed strategy equilibria by just noting that strongly dominated strategies cannot be part of an equilibria.

We'll start by considering the very abstractly defined game.

\begin{center}
\begin{tabular}{l r | c c}
& & \multicolumn{2}{c}{\textbf{Col}} \\
& & $L$ & $R$ \\ \hline
\multirow{2}{*}{\textbf{Row}} & $U$ & \underline{a},b & c,\underline{d} \\
& $D$ & e,\underline{f} & \underline{0},0
\end{tabular}
\end{center}
Don't make too much of the fact that I've put 0 for both players in the bottom right. Since utility scales are only defined up to affine transformation, adding a constant to each possible payout a player gets literally doesn't change the game. So any game representation can be changed into an equivalent representation with 0 for each player in an arbitrarily chosen cell. We'll make that cell be the bottom right, but it could be anywhere.

But do notice the underlining. That has meaning. In particular, it means that all of the following must hold.

\begin{itemize}
\item $a > e$, since $a$ is underlined and $e$ is not, and those are the choices if Col chooses $L$.
\item $0 > c$, since $0$ is underlined (for Row) and $c$ is not, and those are the choices if Col chooses $R$.
\item $d > b$, since $d$ is underlined and $b$ is not, and those are the choices if Row chooses $U$.
\item $f > 0$, since $f$ is underlined and $0$ is not (for Col), and those are the choices if Row chooses $D$.
\end{itemize}
Since there is no cell where both are underlined, there is no pure strategy equilibrium. So the equilibria must include at least one mixed strategy.

Recall how we calculated the expected return $E(M)$ of the mixed strategy $\langle x, L; 1-x, R$, given the expected returns of $L$ and $R$. It is

\begin{equation*}
E(M) = xE(L) + (1-x)E(R)
\end{equation*}
If we let $k = E(R) - E(L)$, we can rewrite this as

\begin{align*}
E(M) &= xE(L) + (1-x)E(R) \\
&= xE(L) + (1-x)(E(L) + k) \\
&= xE(L) + (1-x)E(L) + k(1-x) \\
&= (x + (1-x))E(L) + k(1-x) \\
&= E(L) + k(1-x)
\end{align*}
A similar bit of reasoning shows that $E(M) = E(R) + kx$. Now if the strategy is genuinely mixed, then both $x$ and $1-x$ will be greater than 0. And so if $k \neq 0$, then one of $kx$ and $k(1-x)$ will be positive. And if one of them is positive, then $E(M)$ will be less than one of $E(L)$ and $E(R)$, so it won't be a best response.

So the only time when $\langle x, L; 1-x, R$ is a best response for Col is when $E(L) = E(R)$. In that circumstance, all responses are best responses. In any other circumstance, one of the pure responses is best. So,

\begin{itemize}
\item If the equilibrium has Col playing a mixed strategy, then Row must play a strategy that makes Col's expected returns from $L$ and $R$ equal.
\item And that can only obtain if Row herself plays a mixed strategy.
\end{itemize}
I've done all this from Col's perspective, but we could redo the argument from Row's perspective to get a similar result; if a mixed strategy is to be best for Row, then Col must play the mixed strategy that makes each of Row's pure strategies equally good.

So now we have a recipe, or at least something like one, for finding the mixed strategy equilibria. Each player must choose a mixed strategy that makes the expected return of the other player's pure strategies equal. So let's assume that

\begin{itemize}
\item Row plays $\langle U, x; D, 1-x \rangle$; and
\item Col plays $\langle L, y; R, 1-y \rangle$
\end{itemize}
Then we need $x$ to be such that $E(L) = E(R)$, and $y$ to be such that $E(U) = E(D)$. So we can solve for each of these variables in turn.

\begin{align*}
E(L) &= E(R) \\
xb + (1-x)f &= xd \\
xb + f - xf &= xd \\
f &= xd + xf - xb \\
f &= x(d + f -b) \\
\frac{f}{d + f - b} &= x
\end{align*}
So Row must play the mixed strategy $\langle U, \frac{f}{d + f - b}; D, \frac{d-b}{d + f - b} \rangle$ if there is to be an equilibrium. If she does, then Col's expected return from playing either $L$, or $R$, or any mixture of them, will be $\frac{df - bf}{d + f - b}$.

Let's now solve for $y$ using the same technique.

\begin{align*}
E(U) &= E(D) \\
ya + (1-y)c &= ye \\
ya + c - yc &= ye \\
c &= ye + yc - ya \\
c &= y(e + c - a) \\
\frac{c}{e + c - a} &= y
\end{align*}
So Col must play the mixed strategy $\langle L, \frac{c}{e + c - a}; R, \frac{e-a}{e + c - a} \rangle$ if there is to be an equilibrium. If she does, then Row's expected return from playing either $U$, or $D$, or any mixture of them, will be $\frac{ec}{e + c - a}$.

There is something puzzling about this result. Imagine a version of the game where Col is a driver, and Row a police officer. Col doesn't care much about the law, or the safety of other users of the road. Col just wants to get where they are going as fast as possible, but really doesn't want a speeding ticket. Row wants to catch speeding motorists, but it is costly to engage in detection work. Col has a choice, Speed or Not Speed. Row has a different choice, Detect or Not Detect. Here are the payoffs.

\begin{center}
\begin{tabular}{l r | c c}
& & \multicolumn{2}{c}{\textbf{Col}} \\
& & $S$ & $NS$ \\ \hline
\multirow{2}{*}{\textbf{Row}} & $D$ & 10, -10 & -1, 0 \\
& $ND$ & 0, 1 & 0, 0 
\end{tabular}
\end{center}
The equilibrium of this game is that Col speeds with probability $\frac{1}{11}$, and Row detects with probability $\frac{1}{11}$, and each has an expected payoff of 0. 

Now what happens if the legislature decides that there is too much speeding, and doubles the fine for speeding. Now the payoff table looks like this.

\begin{center}
\begin{tabular}{l r | c c}
& & \multicolumn{2}{c}{\textbf{Col}} \\
& & $S$ & $NS$ \\ \hline
\multirow{2}{*}{\textbf{Row}} & $D$ & 10, -20 & -1, 0 \\
& $ND$ & 0, 1 & 0, 0 
\end{tabular}
\end{center}
If nothing else apart from the penalty for speeding, the rational response by police department Row is to reduce monitoring of speeding, and spend time on other crimes. After all, the new equilibrium is that Row detects with probability $\frac{1}{21}$, and Col still speeds with probability $\frac{1}{11}$. If there is a game in a mixed strategy equilibrium, and the result of any change to the payoffs will be to move it to a new equilibrium, then the only way to change Col's behavior is to either create a new dominant strategy, or to change Row's payoffs. Perhaps what the legislature should do here is to not just increase speeding fines, but to give the money to police departments - that would change the equilibrium of the game in a hurry.

\subsection{Combining Strategies}
We'll work out the equilibria for the following game.

\begin{center}
\begin{tabular}{l r | c c c}
& & \multicolumn{3}{c}{\textbf{Col}} \\
& & $a$ & $b$ & $c$ \\ \hline
\multirow{3}{*}{\textbf{Row}}
& $A$ & 5,0 & 0,1 & 1,-1 \\
& $B$ & 0,2 & 1,0 & 2,1\\
& $C$ & 3,4 & -1,2 & 6,3
\end{tabular}
\end{center}
As a first step, we'll underline the best responses, and see if any equilibria turn up that way.

\begin{center}
\begin{tabular}{l r | c c c}
& & \multicolumn{3}{c}{\textbf{Col}} \\
& & $a$ & $b$ & $c$ \\ \hline
\multirow{3}{*}{\textbf{Row}}
& $A$ & \underline{5},0 & 0,\underline{1} & 1,-1 \\
& $B$ & 0,\underline{2} & \underline{1},0 & 2,1\\
& $C$ & 3,\underline{4} & -1,2 & \underline{6},3
\end{tabular}
\end{center}
That doesn't look like it helps. Row doesn't have any dominated strategies either. But Col does have a dominated strategy, since $a$ dominates $c$. It doesn't matter that $c$ might outperform $b$; there is no reason to play $c$ since $a$ will always do 1 point better. So let's delete that.

\begin{center}
\begin{tabular}{l r | c c}
& & \multicolumn{2}{c}{\textbf{Col}} \\
& & $a$ & $b$  \\ \hline
\multirow{3}{*}{\textbf{Row}}
& $A$ & \underline{5},0 & 0,\underline{1} \\
& $B$ & 0,\underline{2} & \underline{1},0 \\
& $C$ & 3,\underline{4} & -1,2 
\end{tabular}
\end{center}
We don't have to change the underlining, because a dominated strategy is never a best response. 

Now compare options $A$ and $C$. If we rule out Col playing $c$, it turns out there is never a reason for Row to play $C$ over $A$. Option $A$ does better than $C$ whether Col plays $a$ or $b$. So let's delete option $C$ as well.

\begin{center}
\begin{tabular}{l r | c c}
& & \multicolumn{2}{c}{\textbf{Col}} \\
& & $a$ & $b$ \\ \hline
\multirow{2}{*}{\textbf{Row}}
& $A$ & \underline{5},0 & 0,\underline{1} \\
& $B$ & 0,\underline{2} & \underline{1},0 \\
\end{tabular}
\end{center}
Now we have a familiar kind of game. Each player has two options, and there is no pure strategy equilibrium. The problem has reduced to the form of the problem we tackled in the previous subsection. The only equilibrium will be a mixed strategy equilibrium. And we find it by making each player leave the other indifferent between their options. I won't repeat the algebra for finding it, but you can verify yourself that the following works.

\begin{itemize}
\item Row plays $\langle A, \frac{2}{3}; B, \frac{1}{3} \rangle$.
\item Col plays $\langle a, \frac{1}{6}; b, \frac{5}{6} \rangle$.
\end{itemize}
If you prefer to have the mixed strategies range over all of the options that were originally available, with the strategies outside the support getting probability 0, we can write the equilibrium of the game as.

\begin{itemize}
\item Row plays $\langle A, \frac{2}{3}; B, \frac{1}{3}; C, 0 \rangle$.
\item Col plays $\langle a, \frac{1}{6}; b, \frac{5}{6}; c, 0 \rangle$.
\end{itemize}

\section{Why Care about Equilibria?}
There is a philosophical question, alluded to at the end of Chapter 2 of the textbook, about why we should care about Nash equilibria. Thinking about the mixed strategy equilibria can make these questions particularly vivid.

Imagine I'm Col in one of the driving games from the last section. Why should I ever play a mixed strategy? Here is an argument that I should not.

\begin{itemize}
\item If police Row is playing the mixed strategy that brings the game into equilibrium, then my expected return from playing either Speed or Not Speed is 0. So I don't have any incentive to play the particular mixed strategy that brings the game into equilibrium. 
\item Figuring out how to make it that I speed with probability $\frac{1}{11}$ seems like hard work, when it will have no positive return.
\item If police Row is not playing the mixed strategy that brings the game into equilibrium, then one of the pure strategies open to me will have a higher expected return than playing the mixed strategy.
\end{itemize}
So either playing the mixed strategy will involve work for no gain, or it will actually lower my expected return compared to an available alternative. This isn't seeming very promising. But we can imagine some circumstances when the mixed strategy would be more appealing. Here are three related circumstances where it seems useful.

\begin{enumerate*}
\item Imagine that I'm too incompetent to make decisions for myself in games like this. You have to advise me. But the advice will be public; police Row will see it as well. Then the only advice you can rationally offer is that I should play the mixed strategy. Any other advice will backfire because the police will respond to it in ways that make it worthless.
\item Imagine that my strategy will leak to the police, no matter what. That could be because I'm so incompetent that I need to rely on public advice, or it could be because the police have a way of reading my moves. And imagine that the police are in the same boat; their strategy will leak to me. Then we should each play the equilibrium strategy; it is immune to doing worse in light of leaks. That's an implausible situation in most cases, but there is one special case worth considering.
\item Imagine the game will be repeated so many times that eventually our strategies will become visible enough just from the record of our play. Then we are (eventually) back into the `leaked strategy' situation. And in that situation, we should play equilibrium strategies.
\end{enumerate*}
Personally, I think it's the last possibility that makes the concept of an equilibrium play so important. Equilibria (in one shot games) are interesting because they alone are sensible plays in repeated games.

This matches up nicely with our intuitive judgments about games like the driving game. What would happen, intuitively, if one player deviated dramatically from the equilibrium in a one-shot game? The answer is that that player would probably do fine. If Col speeds, then they'll probably get away with it. But think about what would happen (again intuitively) if a player repeated non-equilibrium play over a long period of time. I think it would end up being counterproductive. If Col speeds all the time, Row will adjust and catch them. If Row detects all the time, Col will slow down, but Row will end up wasting detecting resources. It's in repeated games, that non-equilibrium play leads to disaster.

Recall at the start of the chapter we said that it is hard to sometimes make sense of what it even means to play a mixed strategy outside of the context of a repeated game. The reasoning of the last paragraph is why I don't find that worrying. Just like it is hard to define a mixed strategy in a genuinely one-shot game, it is hard to motivate such a strategy. Mixed strategies become easier to define in repeated games, and they become easier to motivate in those games too.

\chapter{Other Equilibrium Concepts}

The primary concept we'll use in this course is Nash Equilibrium. But there are many many other solution concepts for games. Many of them have practical applications. Many of them have a case to be the one true constraint on rational agents playing games.

\section{Rationalizability}

We looked at this briefly last chapter, under the label of `Iterated Best Response'. But it's worth stating the idea more carefully, and saying a bit about why it is an important concept.

A class of strategies is rationalizable if:

\begin{itemize}
\item A class of strategies is rationalizable if every member of the class is a best response, assuming that the probability that every other player plays some strategy from the class is 1;
\item A strategy is rationalizable if it is a member of some class of rationalizable strategies.
\end{itemize}
This is strictly weaker than the notion of a Nash equilibrium. Note that for any Nash equilibrium, the class consisting of just the strategies that form the equilibrium will be rationalizable, so every member of the class will be rationalizable. But some strategies that are not part of any Nash equilibrium are rationalizable. Consider, for example, this game (which is similar to some we've looked at in the previous chapter).

\begin{center}
\begin{tabular}{l r | c c}
& & \multicolumn{2}{c}{\textbf{Col}} \\
& & $a$ & $b$  \\ \hline
\multirow{2}{*}{\textbf{Row}}
& $A$ & 5,0 & 0, 1 \\
& $B$ & 0,1 & 1, 0
\end{tabular}
\end{center}
There is a unique Nash equilibrium for the game. Row plays $\langle A, 0.5; B, 0.5 \rangle$, and Col plays $\langle a, \frac{1}{6}; b, \frac{5}{6} \rangle$. But any strategy whatsoever, including all mixtures, is rationalizable. For any strategy you like, for either player, there is some opposing strategy it is a best response to.

Now compare the following two claims about rationality and game playing.

\begin{enumerate*}
\item Rational players always play some part of a Nash equilibrium strategy.
\item Rational players always play rationalizable strategies.
\end{enumerate*}
These two theories make different assessments of the game I just described. Principle 1 says that rationality rules out all but 1 strategy for each player. Principle 2 says anything goes. Which of them get it right?

If the game was being repeated many times, the answer is easy. The pure strategy $A$ will lead to disaster for Row, since after a while, Col will pick up what's going on and respond by playing $b$. In a repeated game, randomizing as Principle 1 requires is the right thing to do.

But in a one shot game it is harder to say what is correct. There is a nice argument that Principle 1 is self-undermining. Assume that the only rational thing to do is to play the Nash Equilibrium strategy. Then Row will flip a coin to decide what to play. So Col has the same expected return either way, given Row's strategy. Moreover, since Col knows what rationality is, Col knows that Row will flip a coin. So now Col is just facing a decision problem; not a game. And in that decision problem, all responses have equal expected utility. So all responses are equally rational. So Col can rationally do anything. But that's just what Principle 2 says.

The argument of the last paragraph is perfectly general. In any game where the only Nash equilibrium is a mixed strategy equilibrium, the following principles cannot all be true.

\begin{enumerate*}
\renewcommand{\labelenumi}{(\Alph{enumi})}
\item Only strategies that are part of Nash equilibria are rational, and all players know this.
\item All players are rational, and all players know this.
\item Any strategy that is known to have the same expected return as a rational strategy can be rationally played.
\end{enumerate*}
By (A) and (B), the other player will play a mixed strategy. But as we've seen, this will mean that every strategy in the support for one's own mixed strategy will have equal expected return. And by (C), that means that the pure strategies are just as rational as the mixed strategies. But that contradicts the assumption in (A) that only mixed strategies are rational. 

Conventional wisdom in game theory texts is that (C) fails. Principle 2 says that (A) fails. I'm not going to take a stand on this, just to note that there is a clash between the idea that Nash equilibria have a distinctive role to play in rational strategy choice, and the principle that rational agents maximize expected utility.

In games where the players only have two choices to pick from, the difference between the Principles comes down to whether it is obligatory to randomize, or whether it is permissible to play one of the pure strategies from the support of the mixed strategy Nash equilibrium. This not a trivial question to even have intuitions about, because it is tied up with the metaphysics of what a mixed strategy is. But in three way choices there can be more separation between the two principles. Consider either of the following games.

\begin{center}
\begin{tabular}{l r | c c c}
& & \multicolumn{3}{c}{\textbf{Col}} \\
& & $a$ & $b$ & $c$ \\ \hline
\multirow{3}{*}{\textbf{Row}}
& $A$ & 3,0 & 0,3 & 0,2 \\
& $B$ & 0,3 & 3,0 & 0,2 \\
& $C$ & 2,0 & 2,0 & 1,1
\end{tabular}
\end{center}

\begin{center}
\begin{tabular}{l r | c c c}
& & \multicolumn{3}{c}{\textbf{Col}} \\
& & $a$ & $b$ & $c$ \\ \hline
\multirow{3}{*}{\textbf{Row}}
& $A$ & 3,0 & 0,3 & 2,2 \\
& $B$ & 0,3 & 3,0 & 2,2 \\
& $C$ & 2,2 & 2,2 & 2,2
\end{tabular}
\end{center}
In both games, there is a unique Nash equilibrium, $Cc$. But all strategies are rationalizable. Note that $A$ is a best response to $a$, which is a best response to $B$, which is a best response to $b$, which is a best response to $A$. Once you have a cycle like this, every member of the cycle is rationalizable. But there is no mixed strategy Nash equilibrium with those pure strategies as supports. For reasons that should be familiar by now, the only mixed strategy equilibrium that would be available is for each player to play a 50/50 split between their first two options. But if one player plays that mixed strategy, the other player has a return of 1.5 if they reciprocate, and a return of 2 if they play the third option. So there is only one equilibrium, pure or mixed, for the game, while all options are rationalizable.

Ideally, we would have clear intuitions here about what strategies are rationally permissible, and this would let us decide between Principle 1 and Principle 2. Unfortunately, I don't think the intuitions about the cases are at all clear. For what it is worth, I think intuitively that playing one of $A$ or $B$ is rational when the alternative is the equilibrium that returns 1 to each player, but not when the alternative is a guaranteed 2. So these cases don't ultimately settle the question. But they do make more vivid what is at stake. Could it ever be rational, under conditions of common knowledge of rationality, to play $A$ or $B$ in either of these games?

\section{Correlated Equilibrium}
To get Nash equilibrium, we have to assume that every player has a way of letting their play be chosen by a random process. But we assume that these random choices are independent. If Row has a 40\% chance of playing Up, and Column has a 30\% chance of playing Left, then the probability of the Up-Left outcome is 12\%. (That is, 40\% times 30\%.) The idea behind correlated equilibrium is that the players may have a way to coordinate on a random strategy.

Start with a game where we might want such a coordination. Imagine that Row and Col are drivers, approaching an intersection at right angles. Each of them has a choice to go (written as G/g), or stop (written as S/s). They would each prefer to get through the intersection faster. But the other person going and them stopping isn't terrible; there is no crash and the intersection is cleared. And both going is a disaster. So we can approximate the game table as this. (In reality it's more like a many move game, but this is a first approximation.)

\begin{center}
\begin{tabular}{l r | c c}
& & \multicolumn{2}{c}{\textbf{Col}} \\
& & $g$ & $s$  \\ \hline
\multirow{2}{*}{\textbf{Row}}
& $G$ & 0,0 & 10, 9 \\
& $S$ & 9,10 & 8, 8
\end{tabular}
\end{center}
The game has two pure strategy equilibria, $Gs$ and $Sg$. It also has a rather annoying mixed strategy equilibria, where each player plays goes with probability $\nicefrac{2}{11}$, and stays otherwise, and they crash about 3\% of the time. 

Those are all the mixed strategy equilibria. Two of them are asymmetric, and intuitively unfair, and it isn't easy to see how they would get chosen. The only symmetric equilibrium is much worse than either pure strategy equilibria; it even allows for the possibility of crashes.

The correlated equilibria are designed to avoid those problems. Imagine that instead of choosing their own strategy at random, each player could see a common randomising device. Maybe we could have a coin that was flipped in public, and Row goes if it lands heads, and Col goes if it lands tails. That would give each player an expected return of 9.5. And, crucially, once each player sees the (common) signal, each player has no incentive to defect. If Row knows that Col will go if the coin comes up tails, and Row sees that the coin comes up tails, and that Col sees this, Row will prefer to Stay than Go.

This is not completely unlike how we actually solve these problems in real life. Traffic lights are ways of correlating our actions without having one party always getting what they want. I'm simplifying over a lot of details here, but the main thing to stress is that a traffic light (a very public signal) is a way of introducing a new kind of solution to the game drivers face.

In this case, the coin flip is the best symmetric strategy that we can have. But that is to some extent a feature of the values. I won't try to do a just so story to make the following table make sense, but imagine that players faced the following option. (Perhaps it comes about if they really dislike waiting at lights while the other person goes through.)

 \begin{center}
\begin{tabular}{l r | c c}
& & \multicolumn{2}{c}{\textbf{Col}} \\
& & $g$ & $s$  \\ \hline
\multirow{2}{*}{\textbf{Row}}
& $G$ & 0,0 & 10, 7 \\
& $S$ & 7,10 & 9, 9
\end{tabular}
\end{center}
The Nash equilibria don't change much with this change. And the correlated equilibrium of the previous problem still works. That is, the solution: flip a coin in public, if Heads, Row goes, if Tails, Col goes, is still an equilibrium. Each player has an incentive to do their part in the equilibrium once they see the coin. But there is a yet better equilibrium available.

Imagine an outside judge makes three cards: $Gs$; $Sg$ and $Ss$. Each card has a strategy for Row (in Capital letters) and a strategy for Col (in small letters). The judge selects a card, and then tells each player the strategy that is recommended for them, but not what's recommended for the other player. Each player still has an incentive to follow the judge's recommendations. If the judge says go, they know the other will stop, so they will want to go and get the best result. If the judge says stop, it is 50/50 whether the judge has told the other player go or stop. And in that circumstance, going is a bad gamble. It has an expected return of 5; stopping has an expected return of 8. So both parties will do what the judge says.

And if the parties agree in advance to this process, which is fair and symmetric, their expected return is even higher than if they agree to the coin flip. If they flip a coin to decide who will go, their expected return is 8.5. Under this strategy, their expected return is $8\nicefrac{2}{3}$. That's because they have equal chances of getting 7, 9 or 10 as their return.

Let's take stock. A \textbf{correlated equilibrium} is a way of randomly choosing instructions for each player to follow such that each player is motivated to follow the instructions assuming (a) the probabilities of different instructions being chosen are common knowledge, and (b) each other player will follow the instructions. Every correlated equilibrium will consist entirely of rationalizable strategies, but some will not be Nash equilibria. Every Nash equilibria will be a correlated equilibrium.

Given a correlated equilibrium, we can work out the expected return to each player. We do this, as always, by working out their probability of getting each outcome, multiplying the probability by the utility of that outcome, and adding up. We can also find the correlated equilibrium that maximises expected returns, but that takes some more time. Let's look again at the previous table.

 \begin{center}
\begin{tabular}{l r | c c}
& & \multicolumn{2}{c}{\textbf{Col}} \\
& & $g$ & $s$  \\ \hline
\multirow{2}{*}{\textbf{Row}}
& $G$ & 0,0 & 10, 7 \\
& $S$ & 7,10 & 9, 9
\end{tabular}
\end{center}
First, let's work out the mixed strategy equilibrium for this game. We want to find a probability for Row to Go that makes Col indifferent between going and staying. Let that probability be $x$. Then we have all the following lines being equivalent.

\begin{align*}
E(g) &= E(s) \\
10(1-x) &= 7x + 9(1-x) \\
10 - 10x &= 7x + 9 - 9x \\
10 - 10x &= 9 - 2x \\
10 &= 9 + 8x \\
1 &= 8x \\
x &= \frac{1}{8}
\end{align*}
So the mixed strategy equilibrium is that each player goes with probability $\nicefrac{1}{8}$. That leaves them in the crash outcome in $\frac{1}{64}$ of cases. Despite that, the strategies have a surprisingly high expected return of 8.75.

Note one other consequence of the calculation we just did. If the probability of the other player going is $\nicefrac{1}{8}$, the other player will be indifferent between their options. If it is higher than $\nicefrac{1}{8}$, the other player will maximise expected return by stopping; if it is lower, the other player will maximise expected return by going.

Now look at the simplest correlated equilibrium. A coin is flipped, and if it is Heads, Row Goes, and Col stops, while if it is Tails, Row Stops and Col goes. It is easy to confirm that each player will conform to this instruction; if you know what the other player is doing, you do the opposite. But this isn't a great equilibrium. Each player gets 10 with probability 0.5, and 7 with probability 0.5, for an expected return of 8.5. That's worse even than the mixed strategy equilibrium.

If the judge selects one of the three cards $Gs$; $Sg$ and $Ss$ at random, then each player will follow the instructions. That's obvious if they are told to go. If they are told to stop, the probability that the other player was told to go is 0.5. Just to confirm that,

\begin{align*}
&Pr(\text{Other player told to go}|\text{I'm told to stop}) \\
= &\frac{Pr(\text{Other player told to go and I'm told to stop})}{Pr(\text{I'm told to stop})} \\
= &\frac{Pr(\text{Other player told to go and I'm told to stop})}{Pr(\text{They told go and I'm told stop}) + Pr(\text{They told go and I'm told stop})} \\
= &\frac{\frac{1}{3}}{\frac{1}{3} + \frac{1}{3}} \\
= &0.5
\end{align*}
And if the probability that the other player goes is greater than $\nicefrac{1}{8}$, a player will stop. But as we noticed above, the expected return of this strategy is 8$\nicefrac{2}{3}$, which is still lower than the expected return of the mixed strategy. Can we do better?

It turns out that we can. The expected return of each player is just half the sum of the expected returns of each player (in any symmetric equilibrium). And the sum of the expected returns is maximised when we maximise the probability of ending up in the bottom left cell, where the players combine for 18 utils, compared to just 17 utils in the other two non-crash cells. So we want to make the probability of ending up there as high as possible. But we can't just tell the players to go to that cell; they won't have an incentive to cooperate. 

What we want to do is to maximise the probability of getting the $Ss$ instruction, consistent with each player being motivated to follow that instruction. And that will happen as long as each player thinks the probability the other player goes (conditional on them being told to stop) is at least $\nicefrac{1}{8}$. So let's assume that the judge makes the $Gs$ instruction with probability $x$, the $Sg$ instruction with probability $x$, and the $Ss$ instruction with probability $1 - 2x$. Now do the above conditional probability calculation with those values.

\begin{align*}
&Pr(\text{Other player told to go}|\text{I'm told to stop}) \\ 
= &\frac{Pr(\text{Other player told to go and I'm told to stop})}{Pr(\text{I'm told to stop})} \\
= &\frac{Pr(\text{Other player told to go and I'm told to stop})}{Pr(\text{They told go and I'm told stop}) + Pr(\text{They told go and I'm told stop})} \\
= &\frac{x}{x + (1-2x)}
\end{align*}
And we want that last value to be $\nicefrac{1}{8}$. A little calculation shows that this will be true if $x = \nicefrac{1}{9}$. So consider the following instructions:

\begin{itemize}
\item The judge makes the $Gs$ instruction with probability $\nicefrac{1}{9}$.
\item The judge makes the $Sg$ instruction with probability $\nicefrac{1}{9}$.
\item The judge makes the $Ss$ instruction with probability $\nicefrac{7}{9}$.
\end{itemize}
If a player is told to go, they will know the other player is stopping and will go. If a player is told to stop, they will think the probability the other player is going is $\nicefrac{1}{8}$, so they will be indifferent between stopping and going, so they will not have an incentive to deviate from judge's advice.

If the equilibrium is followed, each player expects the following probability distribution over their returns:

\begin{itemize}
\item They get 10 with probability $\nicefrac{1}{9}$;
\item They get 7 with probability $\nicefrac{1}{9}$; and
\item They get 9 with probability $\nicefrac{7}{9}$
\end{itemize}
That means their expected return is $8\nicefrac{8}{9}$, which is a little larger than the mixed strategy equilibrium. And that's the highest expected return of any correlated equilibrium.

The possibility of attractive correlated equilibria shows us something about the value of \textbf{cheap talk}. Cheap talk is the general term we use for any kind of costless signalling. There are all sorts of interesting results that we can derive about costly signalling, and we'll look at that later in the course. But it is somewhat intuitive that cheap talk is never valuable. After all, it is of no use in games like Prisoners Dilemma.

The existence of games with nice correlated equilibria shows that this intuitive finding is actually false. Players who are allowed cheap talk can work out a correlated equilibrium that they both prefer to strategies that are available without cheap talk.

\section{Coordination Games}

Many games have multiple equilibria. In such cases, it is interesting to work through the means by which one equilibrium rather than another may end up being chosen. Consider, for instance, the following three games (two of which we've seen before).

 \begin{center}
\begin{tabular}{l r | c c}
& & \multicolumn{2}{c}{\textbf{Col}} \\
& & $a$ & $b$  \\ \hline
\multirow{2}{*}{\textbf{Row}}
& $A$ & 1,1 & 0,0 \\
& $B$ & 0,0 & 1,1
\end{tabular}
\end{center}

 \begin{center}
\begin{tabular}{l r | c c}
& & \multicolumn{2}{c}{\textbf{Col}} \\
& & $a$ & $b$  \\ \hline
\multirow{2}{*}{\textbf{Row}}
& $A$ & 2,1 & 0,0 \\
& $B$ & 0,0 & 1,2
\end{tabular}
\end{center}

 \begin{center}
\begin{tabular}{l r | c c}
& & \multicolumn{2}{c}{\textbf{Col}} \\
& & $a$ & $b$  \\ \hline
\multirow{2}{*}{\textbf{Row}}
& $A$ & 5,5 & 0,4 \\
& $B$ & 4,0 & 2,2
\end{tabular}
\end{center}

\noindent In each case, both $Aa$ and $Bb$ are Nash equilibria. 

The top game is a purely cooperative game; the players have exactly the same payoffs in every outcome. It is a model for some real-world games. The two players, Row and Col, have to meet up, and it doesn't matter where. They could either go to location $A$ or $B$. If they go to the same location, they are happy; if not, not. 

In such an abstract presentation of game, it is hard to see how we could select an equilibrium out of the two possibilities. In practice, it often turns out not to be so hard. Thomas Schelling (in \textit{The Strategy of Conflict}) noted that a lot of real-life versions of this game have what he called `focal points'. (These are sometimes called Schelling points now, in his honour.) A focal point is a point that stands out from the other options in some salient respect, and it can be expected that other players will notice that it stands out. Here's a nice example devised by the linguist Christopher Potts (\href{http://semanticsarchive.net/Archive/jExYWZlN/potts-interpretive-economy-mar08.pdf}{Interpretive Economy, Schelling Points, and evolutionary stability}).
\begin{quote}

$A$ and $B$ have to select, without communicating, one of the following nine figures. They each get a reward iff they select the same figure.

\begin{center}
\begin{picture}(140,140)
\drawsquare{0}{0}{40}{40}{40}
\drawsquare{50}{0}{90}{40}{40}
\drawsquare{100}{0}{140}{40}{40}
\drawsquare{0}{50}{40}{90}{40}
\drawsquare{50}{50}{90}{90}{40}
\drawsquare{100}{50}{140}{90}{40}
\drawsquare{0}{100}{40}{140}{40}
\drawsquare{50}{100}{90}{140}{40}
\drawsquare{100}{100}{140}{140}{40}

\pictext{20}{116}{1}
\pictext{70}{116}{2}
\pictext{120}{116}{3}
\pictext{20}{66}{4}
\pictext{70}{66}{5}
\pictext{120}{66}{6}
\pictext{20}{16}{7}
\pictext{70}{16}{8}
\pictext{120}{16}{9}

\end{picture}
\end{center}

\end{quote}
%Schilling on coordination

It's not obvious what they should do there, but compare the situation our players face in the following game.

\begin{quote}

$A$ and $B$ have to select, without communicating, one of the following nine figures. They each get a reward iff they select the same figure.


\begin{center}
\begin{picture}(140,140)
\drawsquare{0}{0}{40}{40}{40}
%\drawsquare{50}{0}{90}{40}{40}
\put(70, 20){\circle{40}}
\drawsquare{100}{0}{140}{40}{40}
\drawsquare{0}{50}{40}{90}{40}
\drawsquare{50}{50}{90}{90}{40}
\drawsquare{100}{50}{140}{90}{40}
\drawsquare{0}{100}{40}{140}{40}
\drawsquare{50}{100}{90}{140}{40}
\drawsquare{100}{100}{140}{140}{40}

\pictext{20}{116}{1}
\pictext{70}{116}{2}
\pictext{120}{116}{3}
\pictext{20}{66}{4}
\pictext{70}{66}{5}
\pictext{120}{66}{6}
\pictext{20}{16}{7}
\pictext{70}{16}{8}
\pictext{120}{16}{9}

\end{picture}
\end{center}

\end{quote}


\noindent We could run experiments to test this, but intuitively, players will do better at the second game than the first. That's because in the second game, they have a focal point to select; one of the options stands out from the crowd.

Schelling tested this by asking people where they would go if they had to meet a stranger in a strange city, and had no way to communicate. The answers suggested that meetups would be much more common than you might expect. People in Paris would go to the Eiffel Tower; people in New York would go to (the information booth at) Grand Central Station, and so on. (This is all very contingent. I don't think Grand Central would still work as a focal point in New York. Or maybe it would work as a focal point for people who regularly caught east-bound regional trains from New York.)

The second game we started this section with is often called `Battle of the Sexes'. The real world example of it that is usually used is that Row and Col are a married couple trying to coordinate on a night out. But for whatever reason, they are at the stage where they simply have to go to one of two venues. Row would prefer that they both go to $A$, while Col would prefer that they both go to $B$. But the worst case scenario is that they go to different things. I won't go over that game much more, because the previous section told you most of what you need to know about it. The game has two asymmetric pure strategy Nash equilibria. It has a symmetric, but quite unfortunate, mixed strategy equilibria. And it has a nice correlated equilibrium, where they flip a coin to decide what to do. That's probably what real couples would do, demonstrating again the importance of correlated equilibria.

The third coordination game I started the section with is called Stag Hunt. The idea behind the game goes back to Rousseau, although you can find similar examples in Hume. The general idea is that $A$ is some kind of cooperative action, and $B$ is a `defecting' action. If players make the same selection, then they do well. If they both cooperate, they do better than if they both defect. But if they don't cooperate, it is better to defect than to cooperate.

Rousseau's example was of a group of hunters out hunting a stag, each of whom sees a hare going by. If they leave the group, they will catch the hare for sure, and have something to eat, but the group hunting the stag will be left with a much lower probability of stag catching. Hume was more explicit than Rousseau that these games come up with both two player and many player versions. Here is a nice many player version (from Ben Polak's OpenYale lectures).

\begin{quote}
Everyone in the group has a choice between Investing and Not Investing. The payoff to anyone who doesn't invest is 0. If 90\% or more of the group Invests, then everyone who Invests gains \$1. If less than 90\% of the group Invests, then everyone who invests loses \$1.
\end{quote}

\noindent Again, there are two Nash equilibria: everyone invests, and everyone does not invest. Hume's hypothesis was that in games like this, the cooperative move (in this case Investing) would be more likely in small groups than large groups. (Experiment: Try this out with groups of friends!)

As well as the two pure strategy Nash equilibria in Stag Hunt, there is a mixed strategy Nash equilibria. We find it, as always, by finding the mixed strategy that makes the other player indifferent between their choices. So assume that Row plays $A$ with probability $x$. Then the expected return for Col from her two options is

\begin{align*}
E(a) &= x \times 5 + (1-x) \times 0 \\
 &= 5x \\
E(b) &= x \times 4 + (1-x) \times 2 \\
 &= 4x + 2 - 2x \\
 &= 2 + 2x
\end{align*}

\noindent So there will be equality only if $5x = 2 + 2x$, i.e., only if $x = \nicefrac{2}{3}$. If that happens, $C$'s expected return from any strategy will be $\nicefrac{10}{3}$. Since the game is completely symmetric, a very similar argument shows that if $C$ plays the mixed strategy of $a$ with probability $\nicefrac{2}{3}$, and $b$ with probability $\nicefrac{1}{3}$, then Row will be indifferent between her choices. So each player cooperating with probability $\nicefrac{2}{3}$ is a Nash equilibrium. 

Unlike it Meeting or Battle of the Sexes, the mixed strategy equilibrium is not worse for each player than either pure strategy equilibria. But it isn't a great outcome either. In Stag Hunt, the mixed equilibrium is better than the defecting equilibrium, but worse than the cooperating equilibrium. 

Say an equilibrium is \textbf{Pareto-preferred} to another equilibrium iff every player would prefer the first equilibrium to the second. An equilibrium is \textbf{Pareto-dominant} iff it is Pareto-preferred to all other equilibria. The cooperative equilibrium is Pareto-dominant in Stag Hunt; neither of the other games have Pareto-dominant equiilbria.

Consider each of the above games from the perspective of a player who does not know what strategy their partner will play. Given a probability distribution over the other player's moves, we can work out which strategy has the higher expected return, or whether the various strategies have the same expected returns as each other. For any strategy $S$, and possible strategy $S^\prime$ of the other player,  let $f$ be a function that maps $S$ into the set of probabilities $x$ such that if the other player plays $S^\prime$ with probability $x$, $S$ has the highest expected return. In two strategy games, we can remove the relativisation to $S^\prime$, since for the purposes we'll go on to, it doesn't matter which $S^\prime$ we use. In somewhat formal language, $f(S)$ is the \textbf{basin of attraction} for $S$; it is the range of probability functions that points towards $S$ being played.

Let $m$ be the usual (Lesbegue) measure over intervals; all that you need to know about this measure is that for any interval $[x, y]$, $m([x, y]) = y - x$; i.e., it maps a continuous interval onto its length. Say $m(f(S))$ is the risk-attractiveness of $S$. Intuitively, this is a measure of how big a range of probability distributions over the other person's play is compatible with $S$ being the best thing to play.

Say that a strategy $S$ is risk-preferred iff $m(f(S))$ is larger than $m(f(S^*))$ for any alternative strategy $S^*$ available to that agent. Say that an equilibrium is \textbf{risk-dominant} iff it consists of risk-preferred strategies for all players.

For simple two-player, two-option games like we've been looking at, all of this can be simplified a lot. An equilibrium is risk-dominant iff each of the moves in it are moves the players would make if they assigned probability $\nicefrac{1}{2}$ to each of the possible moves the other player could make.

Neither the simple Meeting game nor Battle of the Sexes have a risk-dominant equilibrium. An asymmetric Meeting game, however, does. Here is one example.

\begin{center}
\begin{tabular}{l r | c c}
& & \multicolumn{2}{c}{\textbf{Col}} \\
& & $a$ & $b$  \\ \hline
\multirow{2}{*}{\textbf{Row}}
& $A$ & 2,2 & 0,0 \\
& $B$ & 0,0 & 1,1
\end{tabular}
\end{center}

In this case $Aa$ is both Pareto-dominant and risk-dominant. Given that, we might expect that both players would choose $A$. (What if an option is Pareto-dominated, and risk-dominated, but focal? Should it be chosen then? Perhaps; the Eiffel Tower isn't easy to get to, or particularly pleasant once you're there, but seems to be chosen in the Paris version of Schelling's meetup game because of its focality.)

The real interest of risk-dominance comes from  Stag Hunt. In that game, cooperating is Pareto-dominant, but defecting is risk-dominant. The general class of Stag Hunt games is sometimes defined as the class of two-player, two-option games with two equilibria, one of which is Pareto-dominant, and the other of which is risk-dominant. In those games, players have a choice between risk and reward. And at the same time they have to choose between a socially useful choice, and a socially harmful choice. The risky choice is the most socially useful one. But note that the socially useful choice is not necessarily a self-sacrifice. Assuming other players are being cooperative, each player does better by being cooperative than by defecting.

It's at this point that the analogy to society formation becomes, for some philosophers, quite vivid. The thought is that we don't require coercion to get people to stay living in civilised society. Assuming that a civilisation is up and running, each member of the society has a preference for staying in the society to leaving it. But some combination of coercion and public spiritedness might be required to get a society running. If the players are unsure whether the others will behave socially or anti-socially, the rational thing to do is to be anti-social. And in the long run, that's a tragic outcome. I'll leave it to other classes, however, to figure out how much social interactions look like iterated Prisoners Dilemma, and how much they look like iterated Stag Hunt.

\chapter{Games and Time}

So far we've looked just at games where each player makes just one move, and they make it simultaneously. That might not feel like most games that you know about. It isn't how we play, for instance, chess. Instead, most games involve players making multiple moves, and these moves taking place in time. Here is one simple such game. It's not very interesting; you won't have fun games nights playing it. But it is useful to study.

\begin{quote}
\textbf{Five} \\
There are two players, who we'll call $A$ and $B$. First $A$ moves, then $B$, then finally $A$ moves again. Each move involves announcing a number, 1 or 2. $A$ wins if after the three moves, the numbers announced sum to 5. $B$ wins otherwise. 
\end{quote}
This is a simple zero-sum game. The payoff is either 1 to $A$ and 0 to $B$, or 1 to $B$ and 0 to $A$. For simplicity, we'll describe this as $A$ winning or $B$ winning. We'll soon be interested in games with draws, which are a payoff of \nicefrac{1}{2} to each player. But for now we're looking at games that someone wins. 

Before we go on, it's worth thinking about how you would play this game from each player's perspective. The formal approach we'll eventually take is pretty similar, I think, to the way one thinks about the game.

\section{Normal Form and Extensive Form}

Figure \ref{FiveGameChart} is the \textbf{extensive form} representation of \textbf{Five}. We will use $W$ to denote that $A$ wins, and $L$ to denote that $B$ wins. 

\begin{figure}[t]
\begin{center}
\begin{picture}(350, 150)
\pictext{175}{0}{$A$}
\put(175, 12){\circle{4}}\put(173, 13){\line(-2, 1){69}}
\put(177, 13){\line(2, 1){69}}
\pictext{135}{20}{1}
\pictext{215}{20}{2}

\pictext{105}{35}{$B$}
\put(105, 47){\circle*{4}}
\put(105, 47){\line(-1, 1){35}}
\put(105, 47){\line(1, 1){35}}
\pictext{80}{55}{1}
\pictext{130}{55}{2}

\pictext{245}{35}{$B$}
\put(245, 47){\circle*{4}}
\put(245, 47){\line(-1, 1){35}}
\put(245, 47){\line(1, 1){35}}
\pictext{220}{55}{1}
\pictext{270}{55}{2}

\pictext{70}{70}{$A$}
\put(70, 82){\circle*{4}}
\put(70, 82){\line(-1, 2){20}}
\pictext{50}{125}{$L$}
\put(70, 82){\line(1, 2){20}}
\pictext{90}{125}{$L$}
\pictext{55}{95}{1}
\pictext{85}{95}{2}

\pictext{140}{70}{$A$}
\put(140, 82){\circle*{4}}
\put(140, 82){\line(-1, 2){20}}
\pictext{120}{125}{$L$}
\put(140, 82){\line(1, 2){20}}
\pictext{160}{125}{$W$}
\pictext{115}{95}{1}
\pictext{155}{95}{2}

\pictext{210}{70}{$A$}
\put(210, 82){\circle*{4}}
\put(210, 82){\line(-1, 2){20}}
\pictext{190}{125}{$L$}
\put(210, 82){\line(1, 2){20}}
\pictext{230}{125}{$W$}
\pictext{195}{95}{1}
\pictext{225}{95}{2}

\pictext{280}{70}{$A$}
\put(280, 82){\circle*{4}}
\put(280, 82){\line(-1, 2){20}}
\pictext{260}{125}{$W$}
\put(280, 82){\line(1, 2){20}}
\pictext{300}{125}{$L$}
\pictext{265}{95}{1}
\pictext{295}{95}{2}


\end{picture}
\end{center}
\caption{\textbf{Five}}
\label{FiveGameChart}
\end{figure}

To read what's going on here, start at the node that isn't filled in, which in this case is the node at the bottom of the tree. The first move is made by $A$. She chooses to play either 1 or 2. If she plays 1, we go to the left of the chart, and now $B$ has a choice. She plays either 1 or 2, and now again $A$ has a choice. As it turns out, $A$ has the same choice to make whatever $B$ does, though there is nothing essential to this. This is just a two-player game, though there is also nothing essential to this. We could easily have let it be the case that a third player moved at this point. Or we could have made it that who moved at this point depended on which move $B$ had made. The extensive form representation allows for a lot of flexibility in this respect.

At the top of the diagram is a record of who won. In the more general form, we would put the payoffs to each player here. But when we have a simple zero-sum game, it is easier to just record the payoffs to one player. You should check that the description of who wins and loses in each situation is right. In each case, $A$ wins if 2 is played twice, and 1 once.

The extensive form representation form is very convenient for turn taking games. But you might think that it is much less convenient for games where players move simultaneously, as in Prisoners' Dilemma. But it turns out that we can represent that as well, through a small notational innovation. Figure \ref{PDEForm} is the game tree for Prisoners' Dilemma.

\begin{figure}[ht]
\begin{center}
\begin{picture}(350, 110)
\pictext{175}{0}{P1}
\put(175, 12){\circle{4}}\put(173, 13){\line(-2, 1){69}}
\put(177, 13){\line(2, 1){69}}
\pictext{135}{20}{A}
\pictext{215}{20}{B}

\pictext{105}{35}{P2}
\put(105, 47){\circle*{4}}
\put(105, 47){\line(-1, 1){35}}
\pictext{70}{85}{(1, 1)}
\put(105, 47){\line(1, 1){35}}
\pictext{140}{85}{(5, 0)}
\pictext{80}{55}{A}
\pictext{130}{55}{B}

\pictext{245}{35}{P2}
\put(245, 47){\circle*{4}}
\put(245, 47){\line(-1, 1){35}}
\pictext{210}{85}{(0, 5)}
\put(245, 47){\line(1, 1){35}}
\pictext{280}{85}{(3, 3)}
\pictext{220}{55}{A}
\pictext{270}{55}{B}

\multiput(105,47)(5, 0){28}{\line(1, 0){3}}

\end{picture}
\end{center}
\caption{Prisoners' Dilemma}
\label{PDEForm}
\end{figure}
The crucial thing here is the dashed line between the two nodes where P2 (short for Player 2) moves. What this means is that P2 doesn't know which of these nodes she is at when she moves. We normally assume that a player knows what moves are available to her. So we normally only put this kind of dashed line in when it connects two nodes at which the same player moves, and the available moves are the same. When we put in notation like this, we say that the nodes that are connected form an \textbf{information set}. If we don't mark anything, we assume that a node is in a degenerate information set, one that only contains itself.

Strictly speaking, you could regard this tree as a representation of a game where Player 1 goes first, then Player 2 moves, but Player 2 does not know Player 1's move when she moves. But that would be reading, I think, too much into the symbolic representation. What's crucial is that we can represent simultaneous move games in extensive form.

There is a philosophical assumption in this notational convention that I just want to mention. We won't have time to cover it further here, though it is worth thinking about. It is usual to use dashed lines, like I've done, or circles or ovals to represent that for all Player 2 knows, she is at one of the connected nodes. But drawing circles around a set of possibilities is only a good way to represent Player 2's uncertainty if we assume quite a lot about knowledge. In particular, it is only good if we assume that which nodes are epistemically open for Player 2 is independent of which node, within that group, she is at. In formal terms, it amounts to assuming that each player knows exactly what they know, and what they don't know. This is not, in general, a true assumption. Normal people have some false beliefs. Those false beliefs are not pieces of knowledge, but the holders of the false belief do not in general know that they do not amount to knowledge. A more philosophically defensible game theory would not assume that everyone knows what it is that they know. But for now we'll follow orthodoxy and assume we can use representations like these dashed lines to represent uncertainty, even though that requires a lot more self-knowledge than most people have.

These representations using graphs are knows as \textbf{extensive form} representations of games. The representations using tables that we've used previously are known as \textbf{normal form} or \textbf{strategic form} representations. The idea is that any game really can be represented as game of one simultaneous move. The `moves' the players make are the selections of \textbf{strategies}. A strategy in this sense is a plan for what to do in any circumstance whatsoever.

Let's do this for the \textbf{Five}. A strategy for $A$ has three variables. It must specify what she does at the first move, what she does at the second move if $B$ plays 1, and what she does at the second move if $B$ plays 2. (We're assuming here that $A$ will play what she decides to play at the first move. It's possible to drop that assumption, but it results in much more complexity.) So we'll describe $A$'s strategy as $\alpha \beta \gamma$, where $\alpha$ is her move to begin with, $\beta$ is what she does if $B$ plays 1, and $\gamma$ is what she does if $B$ plays 2.  A strategy for $B$ needs to only have two variables: what to do if $A$ plays 1, and what to do if $A$ plays 2. So we'll notate her strategy as $\delta \epsilon$, where $\delta$ is what she does if $A$ plays 1, and $\epsilon$ is what she does if $A$ plays 2. So $A$ has 8 possible strategies, and $B$ has 4 possible strategies. Let's record the giant table listing the outcomes if they play each of those strategies.

\begin{center}
\begin{tabular}{l r | c c c c}
& & \multicolumn{4}{c}{$B$} \\
& & 11 & 12 & 21 & 22 \\ \hline
\multirow{8}{*}{$A$} &
  111 & $L$ & $L$ &  $L$ & $L$\\
& 112 & $L$ & $L$ &  $W$ & $W$\\
& 121 & $L$ & $L$ &  $L$ & $L$\\
& 122 & $L$ & $L$ &  $W$ & $W$\\
& 211 & $L$ & $W$ &  $L$ & $W$\\
& 212 & $L$ & $L$ &  $L$ & $L$\\
& 221 & $W$ & $W$ &  $W$ & $W$\\
& 222 & $W$ & $L$ &  $W$ & $L$
\end{tabular}
\end{center}
There is something quite dramatic about this representation. We can see what $A$ should play. If her strategy is 221, then whatever strategy $B$ plays, $A$ wins. So she should play that; it is a (weakly) dominant strategy. This isn't completely obvious from the extended form graph. 

Here's a related fact. Note that there are only 8 outcomes of the extended form game, but 32 cells in the table. Each outcome on the tree is represented by multiple cells of the table. Let's say we changed the game so that it finishes in a draw, represented by $D$, if the numbers picked sum to 3. That just requires changing one thing on the graph; the $L$ in the top-left corner has to be changed to a $D$. But it requires making many changes to the table.

\begin{center}
\begin{tabular}{l r | c c c c}
& & \multicolumn{4}{c}{$B$} \\
& & 11 & 12 & 21 & 22 \\ \hline
\multirow{8}{*}{$A$} &
  111 & $D$ & $D$ &  $L$ & $L$\\
& 112 & $D$ & $D$ &  $W$ & $W$\\
& 121 & $L$ & $L$ &  $L$ & $L$\\
& 122 & $L$ & $L$ &  $W$ & $W$\\
& 211 & $L$ & $W$ &  $L$ & $W$\\
& 212 & $L$ & $L$ &  $L$ & $L$\\
& 221 & $W$ & $W$ &  $W$ & $W$\\
& 222 & $W$ & $L$ &  $W$ & $L$
\end{tabular}
\end{center}
In part because of this fact, i.e., because every change to the value assignment in the extensive form requires making many changes in the values on the normal form, it isn't a coincidence that there's a row containing nothing but $W$. The following claim about our game can be proved.

\begin{quote}
Assign $W$ and $L$ in any way you like to the eight outcomes of the extended form game. Then draw table that is the normal form representation of the game. It will either have a row containing nothing but $W$, i.e., a winning strategy for $A$, or a column containing nothing but $L$, i.e., a winning strategy for $B$.
\end{quote}
We will prove this in a little bit, but first we will look at how to `solve' games like \textbf{Five}.

\section{Backwards Induction}

The way to think through games like \textbf{Five} is by working from top to bottom. $A$ moves last. Once we get to the point of the last move, there is no tactical decision making needed. $A$ knows what payoff she gets from each move, and she simply will take the highest payoff (assuming rationality).

So let's assume she does that. Let's assume, that is, that $A$ does play her best strategy. Then we know three things about what will happen in the game.

\begin{itemize}
\item If $A$ plays 1, and $B$ plays 2, $A$ will follow with 2.
\item If $A$ plays 2, and $B$ plays 1, $A$ will follow with 2.
\item If $A$ plays 2, and $B$ plays 2, $A$ will follow with 1.
\end{itemize}
Moreover, once we make this assumption there is, in effect, one fewer step in the game. Once $B$ moves, the outcome is determined. So let's redraw the game using that assumption, and just listing payoffs after the second move. This will be Figure \ref{FiveGameChartSecondVersion}

\begin{figure}[ht]
\begin{center}
\begin{picture}(350, 100)
\pictext{175}{0}{$A$}
\put(175, 12){\circle{4}}\put(173, 13){\line(-2, 1){69}}
\put(177, 13){\line(2, 1){69}}
\pictext{135}{20}{1}
\pictext{215}{20}{2}

\pictext{105}{35}{$B$}
\put(105, 47){\circle*{4}}
\put(105, 47){\line(-1, 1){35}}
\put(105, 47){\line(1, 1){35}}
\pictext{80}{55}{1}
\pictext{70}{85}{$L$}
\pictext{130}{55}{2}
\pictext{140}{85}{$W$}

\pictext{245}{35}{$B$}
\put(245, 47){\circle*{4}}
\put(245, 47){\line(-1, 1){35}}
\pictext{210}{85}{$W$}
\put(245, 47){\line(1, 1){35}}
\pictext{280}{85}{$W$}
\pictext{220}{55}{1}
\pictext{270}{55}{2}

\end{picture}
\end{center}
\caption{\textbf{Five} with last move assumed}
\label{FiveGameChartSecondVersion}
\end{figure}
Now we can make the same assumption about $B$. Assume that $B$ will simply make her best move in this (reduced) game. Since $B$ wins if $A$ loses, $B$'s best move is to get to $L$. This assumption then, gives us just one extra constraint.

\begin{itemize}
\item If $A$ plays 1, $B$ will follow with 1.
\end{itemize}
And, once again, we can replace $B$'s actual movement with a payoff of the game under the assumption that $B$ makes the rational move. This gives us an even simpler representation of the game that we see in Figure \ref{FiveGameChartThirdVersion}.

\begin{figure}[ht]
\begin{center}
\begin{picture}(350, 65)
\put(175, 12){\circle{4}}\pictext{175}{0}{$A$}
\put(173, 13){\line(-2, 1){69}}
\pictext{105}{50}{$L$}
\put(177, 13){\line(2, 1){69}}
\pictext{245}{50}{$W$}
\pictext{135}{20}{1}
\pictext{215}{20}{2}

\end{picture}
\end{center}
\caption{\textbf{Five} with last two moves assumed}
\label{FiveGameChartThirdVersion}
\end{figure}
And from this version of the game, we can draw two more conclusions, assuming $A$ is rational.

\begin{itemize}
\item $A$ will play 2 at the first move.
\item $A$ will win.
\end{itemize}
Let's put all of that together. We know $A$ will start with 2, so her strategy will be of the form $2 \beta \gamma$. We also know that $B$ doesn't care which strategy she chooses at that point, so we can't make any further reductions. But we do know that if $A$ plays 2 and $B$ plays 1, $A$ will follow with 2. So $A$'s strategy will be of the form $2 2 \gamma$. And we know that if $A$ plays 2 and $B$ plays 2, then $A$ will play 1. So $A$'s strategy will be 221, as we saw on the table.

Note that the use of backwards induction here hides a multitude of assumptions. We have to assume each player is rational, and each player knows that, and each player knows that, and so on for at least as many iterations as there are steps in the game. If we didn't have those assumptions, it wouldn't be right to simply replace a huge tree with a single outcome. Moreover, we have to make those assumptions be very modally robust. 

We can see this with a slight variant of the game. Let's say this time that the left two outcomes are $D$. So the graph looks like Figure \ref{RevisedFiveGameChart}.

\begin{figure}[ht]
\begin{center}
\begin{picture}(350, 150)
\pictext{175}{0}{$A$}
\put(175, 12){\circle{4}}\put(173, 13){\line(-2, 1){69}}
\put(177, 13){\line(2, 1){69}}
\pictext{135}{20}{1}
\pictext{215}{20}{2}

\pictext{105}{35}{$B$}
\put(105, 47){\circle*{4}}
\put(105, 47){\line(-1, 1){35}}
\put(105, 47){\line(1, 1){35}}
\pictext{80}{55}{1}
\pictext{130}{55}{2}

\pictext{245}{35}{$B$}
\put(245, 47){\circle*{4}}
\put(245, 47){\line(-1, 1){35}}
\put(245, 47){\line(1, 1){35}}
\pictext{220}{55}{1}
\pictext{270}{55}{2}

\pictext{70}{70}{$A$}
\put(70, 82){\circle*{4}}
\put(70, 82){\line(-1, 2){20}}
\pictext{50}{125}{$D$}
\put(70, 82){\line(1, 2){20}}
\pictext{90}{125}{$D$}
\pictext{55}{95}{1}
\pictext{85}{95}{2}

\pictext{140}{70}{$A$}
\put(140, 82){\circle*{4}}
\put(140, 82){\line(-1, 2){20}}
\pictext{120}{125}{$L$}
\put(140, 82){\line(1, 2){20}}
\pictext{160}{125}{$W$}
\pictext{115}{95}{1}
\pictext{155}{95}{2}

\pictext{210}{70}{$A$}
\put(210, 82){\circle*{4}}
\put(210, 82){\line(-1, 2){20}}
\pictext{190}{125}{$L$}
\put(210, 82){\line(1, 2){20}}
\pictext{230}{125}{$W$}
\pictext{195}{95}{1}
\pictext{225}{95}{2}

\pictext{280}{70}{$A$}
\put(280, 82){\circle*{4}}
\put(280, 82){\line(-1, 2){20}}
\pictext{260}{125}{$W$}
\put(280, 82){\line(1, 2){20}}
\pictext{300}{125}{$L$}
\pictext{265}{95}{1}
\pictext{295}{95}{2}


\end{picture}
\end{center}
\caption{\textbf{Five}$^{\prime \prime}$}
\label{RevisedFiveGameChart}
\end{figure}

Now we assume that $A$ makes the optimal move at the last step, so we can replace the top row of outcomes with the outcome that would happen if $A$ moves optimally. This gets us Figure \ref{RevisedFiveGameChartReduced}.

\begin{figure}[ht]
\begin{center}
\begin{picture}(350, 100)
\pictext{175}{0}{$A$}
\put(175, 12){\circle{4}}\put(173, 13){\line(-2, 1){69}}
\put(177, 13){\line(2, 1){69}}
\pictext{135}{20}{1}
\pictext{215}{20}{2}

\pictext{105}{35}{$B$}
\put(105, 47){\circle*{4}}
\put(105, 47){\line(-1, 1){35}}
\put(105, 47){\line(1, 1){35}}
\pictext{80}{55}{1}
\pictext{70}{85}{$D$}
\pictext{130}{55}{2}
\pictext{140}{85}{$W$}

\pictext{245}{35}{$B$}
\put(245, 47){\circle*{4}}
\put(245, 47){\line(-1, 1){35}}
\pictext{210}{85}{$W$}
\put(245, 47){\line(1, 1){35}}
\pictext{280}{85}{$W$}
\pictext{220}{55}{1}
\pictext{270}{55}{2}

\end{picture}
\end{center}
\caption{\textbf{Five}$^{\prime \prime}$ with last move assumed}
\label{RevisedFiveGameChartReduced}
\end{figure}

Now assume that $A$ plays 1 on the first round, then $B$ has to move. From \ref{RevisedFiveGameChartReduced} it looks like $B$ has an easy choice to make. If she plays 1, she gets a draw, if she plays 2, then $A$ wins, i.e., she loses. Since drawing is better than losing, she should play 1 and take the draw.

But why think that playing 2 will lead to $A$ winning? The argument that it did depending on assuming $A$ is perfectly rational. And assuming $B$ is in a position to make this choice, that seems like an unlikely assumption. After all, if $A$ were perfectly rational, she'd have chosen 2, and given herself a chance to force a win. 

Now you might think that even if $A$ isn't perfectly rational, it still is crazy to leave her with an easy winning move. And that's probably a sufficient reason for $B$ to accept the draw, i.e., play 1. But the argument that $B$ should regard playing 2 as equivalent to choosing defeat seems mistaken. $B$ knows that $A$ isn't perfectly rational, and she shouldn't assume perfect rationality from here on out.

We will come back to this point, a lot, in subsequent discussions of backwards induction. Note that it is a point that doesn't really arise in the context of normal form games. There we might wonder about whether common knowledge of rationality is a legitimate assumption at the \textit{start} of the game. But once we've settled that, we don't have a further issue to decide about whether it is still a legitimate assumption at later stages of the game.

\section{Value of Games}

Consider games with the following characteristics.

\begin{itemize}
\item $A$ and $B$ take turns making moves. We will call each point at which they make a move, or at which the game ends, a \textbf{node} of the game.
\item At each move, each player knows what moves have been previously made.
\item At each move, the players have only finitely many possible moves open.
\item The players' preferences over the outcomes are opposite to one another. So if $A$ prefers outcome $o_1$ to $o_2$, then $B$ prefers $o_2$ to $o_1$, and if $A$ is indifferent between $o_1$ and $o_2$, then $B$ is indifferent between them as well.
\item $A$'s preferences over the outcomes are complete; for any two outcomes, she either prefers the first, or prefers the second, or is indifferent between them.
\item There is a finite limit to the total number of possible moves in the game.
\end{itemize}

\noindent The finiteness assumptions entail that there are only finitely many possible outcomes. So we can order the outcomes by $A$'s preferences. (Strictly speaking, we want to order sets of outcomes that are equivalence classes with respect to the relation that $A$ is indifferent between them.) Assign the outcome $A$ least prefers the value 0, the next outcome the value 1, and so on.

Now we recursively define the \textbf{value} of a node as follows.

\begin{itemize}
\item The value of a \textbf{terminal node}, i.e., a node at which the game ends, is the payoff at that node.
\item The value of any node at which $A$ makes a choice is the greatest value of the nodes between which $A$ is choosing to move to.
\item The value of any node at which $B$ makes a choice is the least value of the nodes between which $B$ is choosing to move to.
\end{itemize}

\noindent Finally, we say that the value of the game is the value of the initial node, i.e., the node at which the first choice is made. We can prove a number of things about the value of games. The proof will make crucial use of the notion of a \textbf{subgame}. In any extensive form game, a \textbf{subgame} of a game is the game we get by treating any perfect information node as the initial node of a game, and including the rest of the game `downstream' from there. 

By a perfect information node, I mean a node such that when it is reached, it is common knowledge that it is reached. That's true of all the nodes in most of the games we're looking at, but it isn't true in, for instance, the extensive form version of Prisoners' Dilemma we looked at. Nodes that are in non-degenerate information sets, i.e., information sets that contain other nodes, can't trigger subgames. That's because we typically assume that to play a game, players have to know what game they are playing.

Note that a subgame is really a game, just like a subset is a set. Once we're in the subgame, it doesn't matter a lot how we got there. Indeed, any game we represent is the consequence of some choices by the agent; they are all subgames of the game of life.

\begin{itemize}
\item The value of a game is the value of one of the terminal nodes.
\end{itemize}
We prove this by induction on the length of games. (The length of a game is the \textit{maximum} number of moves needed to reach a terminal node. We've only looked so far at games where every path takes the same number of moves to reach a conclusion, but that's not a compulsory feature of games.) 

If the game has zero moves, then it has just one node, and its value is the value of that node. And that node is a terminal node, so the value is the value of a terminal node. 

Now assume that the claim is true for any game of length $k$ or less, and consider an arbitrary game of length $k + 1$ The first node of the game consists of a choice about which path to go down. So the value of the initial node is the value of one of the subsequent nodes. Once that choice is made, we are in a subgame of length $k$, no matter which choice is made. By the inductive hypothesis, the value of that subgame is the value of one of its terminal nodes. So the value of the game, which is the value of one of the immediate subsequent nodes to the initial node, is the value of one of its terminal nodes. 

\begin{itemize}
\item $A$ can guarantee that the outcome of the game is at least the value of the game.
\end{itemize}
Again, we prove this by induction on the length of games. It is trivial for games of length 0. So assume it is true for all games of length at most $k$, and we'll prove it for games of length $k+1$. The initial node of such a game is either a move by $A$ or a move by $B$. We will consider these two cases separately.

Assume that $A$ makes the first move. Then the value of the initial node is the maximum value of any immediate successor node. So $A$ can select to go to the node with the same value as the value of the game. Then we're in a subgame of length $k$. By the inductive assumption, in that game $A$ can guarantee that the outcome is at least the value of the subgame. And since the value of that node is the value of the subgame, so it is also the value of the initial node, i.e., the value of the initial game. So by choosing that node, and starting that subgame, $A$ can guarantee that the outcome is at least the value of the game.

Now assume that $B$ makes the first move. $B$ can choose the node with the least value of the available choices. Then, as above, we'll be in a subgame in which (by the inductive hypothesis) $B$ can guarantee has an outcome which is at most its value. That is, $B$ can guarantee the outcome of the game is at most the value of the initial node of the subgame. And since $B$ can guarantee that that subgame is played, $B$ can guarantee that the game has an outcome of at most its value.

\begin{itemize}
\item $B$ can guarantee that the outcome of the game is at most the value of the game.
\end{itemize}
The proof of this exactly parallels the previous proof, and the details are left as an exercise.

\bigskip \noindent Let's note a couple of consequences of these theorems.

First, assume that the rationality of each player is common knowledge, and that it is also common knowledge that this will persist throughout the game. Then the kind of backwards induction argument we used is discussing \textbf{Five} will show that the outcome of the game will be the value of the game. That's because if $A$ is rational and knows this much about $B$, the outcome won't be lower than the value, and if $B$ is rational and knows this much about $A$, the outcome won't be greater than the value.

Second, these theorems have many applications to real-life games. Both chess and checkers, for instance, satisfy the conditions we listed. The only condition that is not immediately obvious in each case is that the game ends in finite time. But the rules for draws in each game guarantee that is true. (Exercise: Prove this!) Since these games end with White win, Black win or draw, the value of the game must be one of those three outcomes.

In the case of checkers, we know what the value of the game is. It is a draw. This was proved by the Chinook project at the University of Alberta. We don't yet know what the value of chess is, but it is probably also a draw. Given how many possible moves there are in chess, and in principle how long games can go on, it is hard to believe that chess will be `solved' any time soon. But advances in chess computers may be able to justify to everyone's satisfaction a particular solution, even if we can't prove that is the value of chess.

\section{Subgame Perfect Equilibrium}

Consider the following little game. The two players, call them Player I and Player II, have a choice between two options, call them Good and Bad. The game is a sequential move game; first Player I moves then Player II moves. Each player gets 1 if they choose Good and 0 if they choose Bad. We will refer to this game as \textbf{Easy Game}. Its game tree is figure \ref{EasyGameTree}.

\begin{figure}[ht]
\begin{center}
\begin{picture}(350, 110)
\pictext{175}{0}{I}
\put(175, 12){\circle*{4}}
\put(175, 12){\line(-2, 1){70}}
\put(175, 12){\line(2, 1){70}}
\pictext{135}{20}{$G$}
\pictext{215}{20}{$B$}

\pictext{105}{35}{II}
\put(105, 47){\circle*{4}}
\put(105, 47){\line(-1, 1){35}}
\pictext{70}{85}{(1, 1)}
\put(105, 47){\line(1, 1){35}}
\pictext{140}{85}{(1, 0)}
\pictext{80}{55}{$g$}
\pictext{130}{55}{$b$}

\pictext{245}{35}{II}
\put(245, 47){\circle*{4}}
\put(245, 47){\line(-1, 1){35}}
\pictext{210}{85}{(0, 1)}
\put(245, 47){\line(1, 1){35}}
\pictext{280}{85}{(0, 0)}
\pictext{220}{55}{$g$}
\pictext{270}{55}{$b$}

%\multiput(105,47)(5, 0){28}{\line(1, 0){3}}

\end{picture}
\end{center}
\caption{\textbf{Easy Game}}
\label{EasyGameTree}
\end{figure}
A strategy for Player I in \textbf{Easy Game} is just a choice of one option, Good or Bad. A strategy for Player II is a little more complicated. She has to choose both what to do if Player I chooses Good, and what to do if Player II chooses Bad. We'll write her strategy as $\alpha \beta$, where $\alpha$ is what she does if Player I chooses Good, and $\beta$ is what she does if Player II chooses Bad. (We will often use this kind of notation in what follows. Wherever it is potentially ambiguous, I'll try to explain it. But the notation is very common in works on game theory, and it is worth knowing.)

The most obvious Nash equilibrium of the game is that Player I chooses Good, and Player II chooses Good whatever Player I does. But there is another Nash \eqm, as looking at the strategic form of the game reveals. We'll put Player I on the row and Player II on the column. (We'll also do this from now on unless there is a reason to do otherwise.)

\begin{center}
\begin{tabular}{l r | c c c c}
& & \multicolumn{4}{c}{Player II} \\
& & gg & gb & bg & bb \\ \hline 
\multirow{2}{*}{Player I} &
G & 1, 1 & \textbf{1, 1} & 1, 0 & 1, 0 \\
& B & 0, 1 & 0, 0 & 0, 1 & 0, 0 
\end{tabular}
\end{center}
Look at the cell that I've bolded, where Player I plays Good, and Player II plays Good, Bad. That's a Nash \eqm. Neither player can improve their outcome, given what the other player plays. But it is a very odd Nash \eqm. It would be very odd for Player II to play this, since it risks getting 0 when they can guarantee getting 1.

It's true that Good, Bad is weakly dominated by Good, Good. But as we've already seen, and as we'll see in very similar examples to this soon, there are dangers in throwing out \textit{all} weakly dominated strategies. Many people think that there is something else wrong with what Player II does here.

Consider the sub-game that starts with the right-hand decision node for Play\-er II. That isn't a very interesting game; Player I has no choices, and Player II simply has a choice between 1 and 0. But it is a game. And note that Good, Bad is not a Nash \eqm\ of that game. Indeed, it is a \textit{strictly} dominated strategy in that game, since it involves taking 0 when 1 is freely available.

Say that a strategy pair is a \textbf{subgame perfect \eqm} when it is a Nash \eqm, and it is a Nash \eqm\ of every sub-game of the game. The pair Good and Good, Bad is not subgame perfect, since it is not a Nash \eqm\ of the right-hand subgame.

When we solve extensive form games by backwards induction, we not only find Nash equilibria, but subgame perfect equilibria. Solving this game by backwards induction would reveal that Player II would choose Good wherever she ends up, and then Player I will play Good at the first move. And the only subgame perfect \eqm\ of the game is that Player I plays Good, and Player II plays Good, Good.

The stakes here don't seem to high, since the ruled out Nash equilibrium has the same payout as the subgame perfect equilibrium. But that need not be the case. Sometimes the equilibrium that is not subgame perfect is an \textbf{incredible threat}. We mean \textit{incredible} here in the literal sense of not credible. In the spirit of this, let's consider a game I'll call \textbf{Incredible}. Its game tree is figure \ref{IncredibleGameTree}.

\begin{figure}[ht]
\begin{center}
\begin{picture}(350, 110)
\pictext{175}{0}{I}
\put(175, 12){\circle*{4}}
\put(175, 12){\line(-2, 1){70}}
\put(175, 12){\line(2, 1){70}}
\pictext{135}{20}{L}
\pictext{215}{20}{R}

\pictext{105}{35}{II}
\put(105, 47){\circle*{4}}
\put(105, 47){\line(-1, 1){35}}
\pictext{70}{85}{(4, 3)}
\put(105, 47){\line(1, 1){35}}
\pictext{140}{85}{(1, 4)}
\pictext{80}{55}{$g$}
\pictext{130}{55}{$b$}

\pictext{245}{35}{II}
\put(245, 47){\circle*{4}}
\put(245, 47){\line(-1, 1){35}}
\pictext{210}{85}{(3, 2)}
\put(245, 47){\line(1, 1){35}}
\pictext{280}{85}{(0, 0)}
\pictext{220}{55}{$g$}
\pictext{270}{55}{$b$}

%\multiput(105,47)(5, 0){28}{\line(1, 0){3}}

\end{picture}
\end{center}
\caption{\textbf{Incredible}}
\label{IncredibleGameTree}
\end{figure}
Using backwards induction, we can see that the only subgame perfect equilibrium of the game is $Rbg$. That is, Player I plays $R$, and Player II is disposed to play $b$ if Player I plays Left, and $g$ if Player plays $R$. But writing out the able reveals another Nash equilibria.

\begin{center}
\begin{tabular}{l r | c c c c}
& & \multicolumn{4}{c}{Player II} \\
& & $gg$ & $gb$ & $bg$ & $bb$ \\ \hline 
\multirow{2}{*}{Player I} 
& $L$ & 4,3 & 4,3 & 1,4 & \textbf{1,4} \\
& $R$ & 3,2 & 0,0 & \textbf{3,2} & 0,0
\end{tabular}
\end{center}
I've bolded both the Nash equilibria. And note the one in the top right corner of the table. If Player II is committed to playing $bb$, then the best response Player I can make is to play $L$.

But look how implausible this is in real time. It means that if we get to the right hand node, Player II will choose to get 0, when 2 is available. Of course, we won't actually get to that node. But it requires Player II to have dispositions that make no sense. This is the intuition behind subgame perfect equilibrium. We require that players not just make choices which make sense, but are disposed to do so at all times.

In the rest of this chapter, we'll look at reasons why this might be too strong a constraint on rational players.

\section*{Problems with Backwards Induction}

\newcommand{\tol}[1]{$\langle #1 \rangle$}

When each player only moves once, it is very plausible that each player should play their part of a subgame perfect \eqm\ solution to the game. But this is less compelling when players have multiple moves. We can start to see why by looking at a game that the philosopher Robert Stalnaker has used in a few places to argue against backwards induction reasoning. Again, it is an extensive form game. We will call it \textbf{Cooperative Caterpillar}.

\begin{figure}[ht]
\begin{center}
\begin{picture}(260, 100)
\put(20, 80){\line(1,0){220}}

\put(20, 80){\circle*{4}}
\put(20, 80){\line(0, -1){60}}
\pictext{20}{85}{I}
\pictext{60}{65}{$A_1$}
\pictext{28}{45}{$D_1$}
\pictext{20}{8}{2}

\put(100, 80){\circle*{4}}
\put(100, 80){\line(0, -1){60}}
\pictext{100}{85}{II}
\pictext{140}{65}{$a$}
\pictext{108}{45}{$d$}
\pictext{100}{8}{1}

\put(180, 80){\circle*{4}}
\put(180, 80){\line(0, -1){60}}
\pictext{180}{85}{I}
\pictext{220}{65}{$A_2$}
\pictext{188}{45}{$D_2$}
\pictext{180}{8}{0}

\pictext{248}{75}{3}

\end{picture}
\end{center}
\caption{\textbf{Cooperative Caterpillar}}
\label{PDEForm}
\end{figure}

The game starts in the upper-left corner. There are up to three moves, each of them Across or Down. As soon as one player moves Down, the game ends. It is a common interest game; each player gets the payout at the terminal node.

Since there is only one path to each decision node, a strategy merely has to consist of a set of plans for what to play if that node is reached. Alice's strategy will consist of two capitalised moves, and Bob's strategy will consist of one lower-case move.

If we apply backwards induction, we get that the unique solution of the game is \tol{A_1A_2, a}. Alice would play $A_2$ if it gets that far. Given that, Bob would be better off playing $a$ than $d$ if he gets to play. And given that, Alice is better off playing $A_1$ than $D_1$.

But there are many other Nash equilibria of the game. One of them is \tol{D_1A_2, d}. Given that Player I is playing $D_1$, Player II can play anything without changing her payout; it will be 2 come what may. Given that Player II is playing $d$, Player I is best off playing $D_1$ and taking 2, rather than just getting the 1 that comes from leaving Player II to make a play.

Could this \eqm\ be one that rational players, who know each other to be rational, reach? Stalnaker argues that it could. Assume each player knows that the other player is rational, and is playing that strategy. Given what Player I knows, her strategy choice is clearly rational. She takes 2 rather than the 1 that she would get by playing $A_1$, and she is disposed to take 3 rather than 0 if she gets to the final decision node. So her actual choice is rational, and her disposition is to make another rational choice if we reach the end of the game.

Things are a little trickier for Player II. You might think it is impossible for rational Player II, who knows Player I to be rational, to move $d$. After all, if Player I is rational, then she'll play $A_2$, not $D_2$. And if she plays $A_2$, it is better to play $a$ than $d$. So it looks hard to justify Player II's move. But looks can be deceiving. In fact there isn't anything wrong with Player II's move, as long as he has the right beliefs to justify it. It's very important to distinguish the following two conditionals.

\begin{itemize}
\item If Player I has the choice, she chooses $A_2$ over $D_2$.
\item If Player I were to have the choice, she would choose $A_2$ over $D_2$.
\end{itemize}

\noindent Player II knows that the first, indicative, conditional is true. And indeed it is true. But he doesn't know that the second, subjunctive, conditional is true. After all, if Player I were to have the choice between $A_2$ and $D_2$, she would have, \textit{irrationally}, chosen $A_1$ over $D_1$. And if she had chosen irrationally once, it's possible that she would choose irrationally again.

Here's an analogy that may help explain what's going on. The following set of beliefs is consistent.
\begin{itemize}
\item Any perfectly rational being gets all of the questions on their algebra exam right.
\item Alice is perfectly rational.
\item If Alice had got the second-hardest question on the algebra exam wrong, she would have got the hardest question on the algebra exam wrong as well.
\end{itemize}
 Player II's beliefs are like that. He believes Player I is perfectly rational. He also believes that if Player I were to make an irrational move, she would continue to make irrational moves. That's consistent with belief in perfect rationality, and nothing about the game setup rules out such a belief. He also believes that playing $A_1$ would be irrational. That's correct, given what Player I knows about Player II. Given all those beliefs, playing $d$, if he had the chance, would be rational.

Stalnaker argues that many game theorists have tacitly confused indicative and subjunctive conditionals in reasoning about games like \textbf{Cooperative Caterpillar}. There are other games where similar confusions arguably take place.

\section*{Money Burning Game}
Elchanen Ben-Porath and Eddie Dekel suggested this variant to the Battle of the Sexes game. The variation is in two steps. First, we change the payouts for the basic game to the following. (Note that I'm using a lowercase letter for one of Row's options here; this is to distinguish the $d$ of down from the $D$ of Don't burn.)

\begin{center}
\begin{tabular}{l r | c c}
& & \multicolumn{2}{c}{Player II} \\
& & $l$ & $r$ \\ \hline
\multirow{2}{*}{Player I}
& $u$ & 4, 1 & 0, 0 \\
& $d$ & 0, 0 & 1, 4 \\
\end{tabular}
\end{center}
Then we give Player I the option of publicly burning 2 utils before playing this game. We will use $D$ for Don't burn, and $B$ for burn. So actually each player has four choices. Player I has to choose both $D$ or $B$, then $u$ or $d$. Player 2 has to choose whether to play $l$ or $r$ in each of the two possibilities: first, when $D$ is played, second, when $B$ is played. We'll write $lr$ for the strategy of playing $l$ if $D$, and $r$ if $B$, and so on for the other strategies.

\begin{center}
\begin{tabular}{l r | c c c c}
& & \multicolumn{4}{c}{\textbf{Player II}} \\
& & $ll$ & $lr$ & $rl$ & $rr$ \\ \hline
\multirow{4}{*}{\textbf{Player I}}
& $Du$ & 4, 1 & 4, 1 & 0, 0 & 0, 0\\
& $Dd$ & 0, 0 & 0, 0 & 1, 4 & 1, 4\\
& $Bu$ & 2, 1 & -2, 0 & 2, 1 & -2, 0\\
& $Bd$ & -2, 0 & -1, 4 & -2, 0 & -1, 4
\end{tabular}
\end{center}
Now we can analyse this game a couple of ways. First, we can go through eliminating weakly dominated strategies. Note that $Du$ strictly dominated $Bd$, so we can eliminate it. If $Bd$ is out, then $ll$ weakly dominates $lr$, and $rl$ weakly dominates $rr$. So we can eliminate $lr$ and $rr$. Now $Bu$ weakly dominates $Dd$, relative to the remaining options, so we can eliminate it. Given that just $Du$ and $Bu$ remain, $ll$ weakly dominates all options for Player II, so it is all that is left. And given that $ll$ is all that remains, $Du$ strongly dominates $Bu$. So the iterative deletion of weakly dominated strategies leaves us with \tol{Du, ll} as the unique solution of the game.

Alternatively, we can think the players reason as follows. (The following reasoning is sometimes called `forward induction' reasoning.) Playing $Bd$ is obviously irrational for Player I, since its maximum return is -1, and playing $Du$ or $Dd$ guarantees getting at least 0. So any rational player who plays $B$ must be playing $Bu$. That is, if Player I is rational and plays $B$, she will play $Bu$. Moreover, this is common knowledge among the players. So if Player I plays $B$, she will recognise that Player II will know that she is playing $Bu$. And if Player II knows Player I is playing $u$ after burning, she will play $l$, since that returns her 1 rather than 0. So Player I knows that playing $B$ leads to the 2,1 state; i.e., it returns her 2. But now it is irrational to play $Dd$, since that gets at most 1, and playing $B$ leads to a return of 2. So Player I will play $Bu$ or $Du$. And since the reasoning that leads to this is common knowledge, Player II must play $ll$, since playing $l$ is her best response to a play of $u$, and she knows Player I will play $u$. But if Player II is going to play $ll$, Player I doesn't need to burn; she can just play $Du$.

There is a crucial conditional in the middle of that argument; let's isolate it.

\begin{itemize}
\item If Player I is rational and plays $B$, she will play $Bu$.
\end{itemize}
But that's not what is needed to make the argument work. What Player I needs, and in fact needs to be common knowledge, is the following subjunctive.

\begin{itemize}
\item If Player I is rational then, if she were to play $B$, she would play $Bu$.
\end{itemize}
But in fact there's no reason to believe that. After all, if the forward induction argument is right, then it is \textit{irrational} to burn the money. So if Player I were to play $B$, i.e., burn the money, then she would be doing something irrational. And the fact that Player I is actually rational is consistent with the counterfactual that if she were to do one irrational thing, it would be because she is following an irrational strategy. Note that on the assumption that $Du$ is optimal, then if Player I finds herself having played $B$, it isn't clear that playing $d$ is irrational; it isn't like it is dominated by $u$.

So it isn't true that if Player I were to burn the utils, Player II would know that she is playing $Bu$, and react accordingly by playing $l$. And if that's not true, then there's no reason to think that $Bu$ is better than $Dd$. And if there's no reason to think that, there's no reason to be confident Player 2 will play $ll$.

Robert Stalnaker, in his commentary on the game, goes further and provides a positive model where the players start with a common belief in rationality, and in what each other will do, but in which the players play $Bu$ and $rl$. I'll leave it as an exercise to work out what counterfactuals the agents have to believe to make this rational, but suffice to say that it is a situation where each side has true beliefs about the other, and maximises expected utility given their beliefs.

There is another, relatively simple, \eqm\ to the game. Player I will play $Dd$, and Player II will play $rr$. Player II is certain that Player I will play $d$, and will keep that belief even if Player I irrationally burns 2 utils to start with. Given that certainty, it is rational for Player I to play $Dd$, so Player II's belief is consistent with believing Player I to be rational. And since Player I is playing $Dd$, playing $rr$ is perfectly rational for Player II; indeed it results in her best possible outcome. Moreover, given that Player II is playing $rr$, it makes sense for Player I to play $d$ even if they were, irrationally, to burn the utils. So even though playing $Bd$ would be irrational, since it is strictly dominated, if they were to (irrationally) play $D$, they \textit{should} follow that by playing $d$. There's an important point about the scope of the claim that playing $Bd$ is irrational. It \textit{doesn't} imply that if the player were to irrationally play $B$, it would be irrational to follow with a play of $d$. If Player I was convinced that Player II was playing $rr$, then it would be rational to follow $B$ with $d$.

\section*{Iterated Prisoners' Dilemma}
Let's return to Prisoners Dilemma, using $C$ and $D$ for cooperate and defect. We'll call the players I and II, and use uppercase letters for Player I's strategies, and lower case letters for Player II's strategies. We'll assume that each player knows that the other players are perfectly rational in Stalnaker's sense. (That is, they have true beliefs about the other, and they maximise expected utility given their beliefs.)

There are a couple of arguments that the players must end up at the \eqm\ where every player defects on every move. One of these is an argument by backwards induction. 

At the last move of the game, defecting dominates cooperating. Both players know this. At the second-last move, you might have thought antecedently that there was some benefit to cooperating. After all, it might induce cooperation in the other player. But the only benefit one could get from cooperating was cooperation at the next (i.e., last) move. And no rational player will cooperate on the last move. So, if you're playing with a rational player, there's no benefit to cooperating on the second-last move. But if that's right, and everyone knows it, then there is no benefit to cooperating on the third-last move. The only benefit would be if that would induce cooperation, and it couldn't, since we just proved that any rational player would defect on the second-last move. And so on for any finite length game that you like.

Alternatively, we could look at the game from a strategic perspective. As we've seen in games like \textbf{Burning}, sometimes there are Nash equilibria that don't appear in backwards induction reasoning. But this isn't the case in finitely iterated Prisoners' Dilemma. The only Nash \eqm\ is that both players defect in every circumstance.

This is a little tricky to check since the strategic form of iterated Prisoners' Dilemma gets very complex very quickly. Let's just consider the three round version of the game. Already each player has 128 strategies to choose from. The choice of a strategy involves making 7 distinct choices:

\begin{enumerate*}
\item What to do in the first round.
\item What to do in the second round if the other player cooperates in the first round.
\item What to do in the second round if the other player defects in the first round.
\item What to do in the third round if the other player cooperates in each of the first two rounds.
\item What to do in the third round if the other player cooperates in the first round then defects in the second round.
\item What to do in the third round if the other player defects in the first round then cooperates in the second round.
\item What to do in the third round is the other player defects in each fo the first two rounds.
\end{enumerate*}
Since these 7 choices are distinct, and the player has 2 choices at each point, there are $2^7 = 128$ possible strategies. So the strategic form of the table involves $128 \times 128 = 16384$ cells. Needless to say, we \textit{won't} be putting that table here. (Though it isn't too hard to get a computer to draw the table for you. The four round game, where each player has $2^{15}$ choices, and there are over one billion cells in the decision table, requires more computing power!)

We'll write a strategy as \tol{x_1x_2x_3x_4x_5x_6x_7}, where $x_i$ is 0 if the player's answer to the $i$'th question above is to cooperate, and 1 if it is to defect. So \tol{0000000} is the strategy of always cooperating, \tol{1111111} is the strategy of always defecting, \tol{0010101} is `tit-for-tat', the strategy of cooperating on the first move, then copying the other player's previous move, and so on.

Let \tol{x_1x_2x_3x_4x_5x_6x_7} be any strategy that doesn't always involve defection on the final round, i.e., a strategy where $x_4 + x_5 + x_6 + x_7 < 4$. It is easy enough to verify that such a strategy is weakly dominated by \tol{x_1x_2x_31111}. In some plays of the game, the defection on the final round leads to getting a better outcome. In other plays of the game, \tol{x_1x_2x_3x_4x_5x_6x_7} and \tol{x_1x_2x_31111} have the same play. For instance, \tol{0001000} will do just as well as \tol{0001111} if the opponent plays any strategy of the form \tol{000x_4x_5x_6x_7}, but it can never do better than \tol{0001111}. So if each player is perfectly rational, we can assume their strategy ends with $1111$.

That cuts each player's choices down to 8. Let's do that table. We'll use $C$ and $c$ rather than 0, and $D$ and $d$ rather than 1, so it is easier to distinguish the two players' strategies. When we label the table, we'll leave off the trailing $D$s and $d$'s, since we assume players are defecting on the last round. We'll also leave off the 3 units the players each get in the last round. (In other words, this will look a lot like the table for a \textit{two} round iterated Prisoners' Dilemma, but it is crucial that it is actually a three-round game.)

%\starttab{r c c c c c c c c}
%\gamelab{PDTable} & $ccc$ & $ccd$ & $cdc$ & $cdd$ & $dcc$ & $dcd$ & $ddc$ & $ddd$ \\
%$CCC$ & 14, 14 & 14, 14 & 7, 16 & 7, 16 & 7, 16 & 7, 16 & 0, 18 & 0, 18 \\
%$CCD$ & 14, 14 & 14, 14 & 7, 16 & 7, 16 & 9 & 9 & 3 & 3 \\
%$CDC$ & 16, 7 & 16, 7 & 10, 10 & 10, 10 & 7 & 7 & 0 & 0 \\
%$CDD$ & 16, 7 & 16, 7 & 10, 10 & 10, 10 & 9 & 9 & 3 & 3 \\
%$DCC$ & 16, 7 & 9, 9 & 16, 7 & 9& 10 & 3 & 10 & 3 \\
%$DCD$ & 16, 7 & 9 ,9 & 16, 7 & 9& 12 & 6 & 12 & 6 \\
%$DDC$ & 18, 0 & 12, 3 & 18, 0 & 12 & 10 & 3 & 10 & 3 \\
%$DDD$ & 18, 0 & 12, 3 & 18, 0 & 12& 10 & 6 & 12 & 6 \\
%\fintab

\begin{center}
\begin{tabular}{l r | c c c c c c c c}
& & \multicolumn{8}{c}{\textbf{Player II}} \\
& & $ccc$ & $ccd$ & $cdc$ & $cdd$ & $dcc$ & $dcd$ & $ddc$ & $ddd$ \\ \hline
\multirow{8}{*}{\textbf{Player I}} 
& $CCC$ & 6, 6 & 6, 6 & 3, 8  & 3, 8  & 3, 8 & 3, 8 & 0, 10 & 0, 10\\
& $CCD$ & 6, 6 & 6, 6 & 3, 8  & 3, 8  & 5, 5 & 5, 5 & 1, 6 & 1, 6\\
& $CDC$ & 6, 6 & 8, 3 & 4, 4  & 4, 4  & 3, 8 & 3, 8 & 0, 10 & 0, 10\\
& $CDD$ & 6, 6 & 8, 3 & 4, 4  & 4, 4  & 5, 5 & 5, 5 & 1, 6 & 1, 6\\
& $DCC$ & 6, 6 & 5, 5 & 8, 3  & 5, 5  & 4, 4 & 1, 6 & 4, 4 & 1, 6\\
& $DCD$ & 6, 6 & 5, 5 & 8, 3  & 5, 5  & 6, 1 & 2, 2 & 6, 1 & 2, 2\\
& $DDC$ & 6, 6 & 6, 1 & 10, 0 & 6, 1 & 4, 4 & 1, 6 & 4, 4 & 1, 6\\
& $DDD$ & 6, 6 & 6, 1 & 10, 0 & 6, 1 & 6, 1 & 2, 2 & 6, 1 & 2, 2\\
\end{tabular}
\end{center}
In this game $CCC$ is \textit{strongly} dominated by $CDD$, and $DCC$ is strongly dominated by $DDD$. Similarly $ccc$ and $dcc$ are strongly dominated. So let's delete them.

\begin{center}
\begin{tabular}{l r | c c c c c c}
& & \multicolumn{6}{c}{\textbf{Player II}} \\
& & $ccd$ & $cdc$ & $cdd$ & $dcd$ & $ddc$ & $ddd$ \\ \hline
\multirow{6}{*}{\textbf{Player I}} 
& $CCD$ & 6, 6 & 3, 8  & 3, 8  & 5, 5 & 1, 6 & 1, 6\\
& $CDC$ & 8, 3 & 4, 4  & 4, 4  & 3, 8 & 0, 10 & 0, 10\\
& $CDD$ & 8, 3 & 4, 4  & 4, 4  & 5, 5 & 1, 6 & 1, 6\\
& $DCD$ & 5, 5 & 8, 3  & 5, 5  & 2, 2 & 6, 1 & 2, 2\\
& $DDC$ & 6, 1 & 10, 0 & 6, 1 & 1, 6 & 4, 4 & 1, 6\\
& $DDD$ & 6, 1 & 10, 0 & 6, 1 & 2, 2 & 6, 1 & 2, 2\\
\end{tabular}
\end{center}
Note that each of these is a best response. In particular,
\begin{itemize}
\item $CCD$ is a best response to $dcd$.
\item $CDC$ is a best response to $ccd$.
\item $CDD$ is a best response to $ccd$ and $dcd$.
\item $DCD$ is a best response to $ddc$.
\item $DDC$ is a best response to $cdc$ and $cdd$.
\item $DDD$ is a best response to $cdc$, $cdd$, $ddc$ and $ddd$.
\end{itemize}
Now the only Nash \eqm\ of that table is the bottom-right corner. But there are plenty of strategies that, in a strategic version of the game, would be consistent with common belief in perfect rationality. The players could, for instance, play $CDC$ and $cdc$, thinking that the other players are playing $ccd$ and $CCD$ respectively. Those beliefs would be false, but they wouldn't be signs that the other players are irrational, or that they are thinking the other players are irrational.

But you might suspect that the strategic form of the game and the extensive form are crucially different. The very fact that there are Nash equilibria that are not subgame perfect equilibria suggests that there are tighter constraints on what can be played in an extensive game consistent with rationality and belief in rationality. Stalnaker argues, however, that this isn't right. In particular, any strategy that is (given the right beliefs) perfectly rational in the strategic form of the game is also (given the right beliefs and belief updating dispositions) perfectly rational in the extensive form. We'll illustrate this by working more slowly through the argument that the game play \tol{CDC, cdc} is consistent with common belief in perfect rationality.

The first thing you might worry about is that it isn't clear that $CDC$ is perfectly rational, since it is looks to be weakly dominated by $CDD$, and perfect rationality is inconsistent with playing weakly dominated strategies. But in fact $CDC$ isn't weakly dominated by $CDD$. It's true that on the six columns represented here, $CDC$ never does better than $CDD$. But remember that each of these strategies is short for a \textit{three}-round strategy; they are really short for $CDCDDDD$ and $CDDDDDD$. And $CDCDDDD$ is not weakly dominated by $CDDDDDD$; it does better, for example, against $dcccddd$. That is the strategy is defecting the first round, cooperating the second, then defecting on the third round unless the other player has cooperated on each of the first two rounds. Since $CDC$ does cooperate each of the first two rounds, it gets the advantage of defecting against a cooperator on the final round, and ends up with 8 points, whereas $CDD$ merely ends up with 6.

But why would we think Player II might play $dcccddd$? After all, it is itself a weakly dominated strategy, and perfectly rational beings (like Player II) don't play weakly dominated strategies. But we're not actually thinking Player II \textit{will} play that. Remember, Player I's assumption is that Player II will play $ccddddd$, which is not a weakly dominated strategy. (Indeed, it is tit-for-tat-minus-one, which is a well-known strategy.) What Player I also thinks is that if she's wrong about what Player II will play, then it is possible that Player II is not actually perfectly rational. That's consistent with believing Player II is actually perfectly rational. The assumption of common belief in perfect rationality is not sufficient for the assumption that one should believe that the other player is perfectly rational \textit{no matter what surprises happen}. Indeed, that assumption is barely coherent; some assumptions are inconsistent with perfect rationality.

One might be tempted by a weaker assumption. Perhaps a player should hold on to the belief that the other player is perfectly rational unless they get evidence that is \textit{inconsistent} with that belief. But it isn't clear what could motivate that. In general, when we are surprised, we have to give up something that isn't required by the surprise. If we antecedently believe $p \wedge q$, and learn $\neg (p \wedge q)$, then what we've learned is inconsistent with neither $p$ nor $q$, but we must give up one of those beliefs. Similarly here, it seems we must be prepared to give up a belief in the perfect rationality of the other player in some circumstances when that is not entailed by our surprise. (Note that even if you don't buy the argument of this paragraph, there still isn't a reason to think that perfectly rational players can't play $CDD$. But analysing that possibility is beyond the scope of these notes.)

What happens in the extensive form version of the game? Well, each player cooperates, thinking this will induce cooperation in the other, then each player is surprised by a defection at round two, then at the last round they both defect because it is a one-shot Prisoners' Dilemma. Nothing seems irrational there. We might wonder why neither defected at round one. Well, if they believed that the other player was playing tit-for-tat-minus-one, then it is better to cooperate at round one (and collect 8 points over the first two rounds) than to defect (and collect at most 6). And playing tit-for-tat-minus-one is rational if the other person is going to defect on the first and third rounds, and play tit-for-tat on round two. So as long as Player I thinks that Player II thinks that Player I thinks that she is going to defect on the first and third rounds, and play tit-for-tat on round two, then her cooperation at round one is rational, and consistent with believing that the other player is rational.

But note that we had to attribute an odd belief to Player II. We had to assume that Player II is \textit{wrong} about what Player I will play. It turns out, at least for Prisoners' Dilemma, that this is crucial. If both players are certain that both players are perfectly rational, and have no false beliefs, then the only strategy that can be rationalised is permanent defection. The proof (which I'm leaving out) is in Stalnaker's ``Knowledge, Belief and Counterfactual Reasoning in Games''.

This \textit{isn't} because the no false beliefs principle suffices for backwards induction reasoning in general. In \textbf{Cooperative Caterpillar}, we can have a model of the game where both players are perfectly rational, and have correct beliefs about what the other player will do, and both those things are common belief, and yet the backwards induction solution is not played. In that game $A$ believes, truly, that $B$ will play $d$, and $B$ believes, truly, that $A$ will play $A_1D_2$. And each of these moves is optimal given (true!) beliefs about the other player's strategy.

But Prisoners' Dilemma is special. It isn't that in a game played between players who know each other to have no false beliefs and be perfectly rational that we must end up at the bottom-right corner of the strategic table. But we must end up in a game where every player defects every time. It could be that Player I thinks Player II is playing either $dcd$ or $ddd$, with $ddd$ being much more, and on that basis  she decides to play $dcd$. And Player II could have the converse beliefs. So we'll end up with the play being \tol{DCD, dcd}. But of course that means each player will defect on the first two rounds, and then again on the last round since they are perfectly rational. In such a case the requirement that each player have no false beliefs won't even be enough to get us to Nash \eqm\, since \tol{DCD, dcd} is not a Nash \eqm. But it is enough to get permanent defection.

\chapter{Bayesian Games}

\section{Perfect Bayesian \Eqm}

The core idea behind subgame perfect \eqm\ was that we wanted to eliminate equilibria that relied on `incredible threats'. That is, there are some equilibria such that the first player makes a certain move only because if they make a different move, the second player to move would, on their current strategy, do something that makes things worse for both players. But there's no reason to think that the second player would actually do that.

The same kind of consideration can arise in games where there aren't any subgames. For instance, consider the following game, which we'll call \textbf{Perfect}.

\begin{figure}[ht]
\begin{center}
\begin{picture}(200, 100)
\put(100, 15){\circle{4}}
\pictext{100}{0}{1}

\put(100, 15){\line(-2, 1){70}}
\put(100, 15){\line(0, 1){35}}
\put(100, 15){\line(2, 1){70}}

\pictext{60}{25}{$L$}
\pictext{108}{25}{$M$}
\pictext{140}{25}{$R$}

\multiput(30, 50)(4, 0){18}{\line(1, 0){2}}

\pictext{65}{55}{2}

\put(30, 50){\circle*{4}}
\put(100, 50){\circle*{4}}

\put(30, 50){\line(-1, 2){20}}
\put(30, 50){\line(1, 2){20}}

\put(100, 50){\line(-1, 2){20}}
\put(100, 50){\line(1, 2){20}}

\pictext{15}{60}{$l$}
\pictext{45}{60}{$r$}

\pictext{85}{60}{$l$}
\pictext{115}{60}{$r$}

\pictext{10}{90}{(3, 1)}
\pictext{50}{90}{(0, 0)}

\pictext{80}{90}{(1, 1)}
\pictext{120}{90}{(0, 0)}

\pictext{170}{50}{(2, 2)}
\end{picture}
\end{center}
\caption{\textbf{Perfect}}
\end{figure}
The game here is a little different to one we've seen so far. First Player 1 makes a move, either $L$, $M$ or $R$. If her move is $R$, the game ends, with the (2, 2) payout. If she moves $L$ or $M$, the game continues with a move by Player 2. But crucially, Player 2 does not know what move Player 1 made, so she does not know which node she is at. So there isn't a subgame here to start. Player 2 chooses $l$ or $r$, and then the game ends.

There are two Nash equilibria to \textbf{Perfect}. The obvious \eqm\ is $Ll$. In that \eqm\, Player 1 gets her maximal payout, and Player 2 gets as much as she can get given that the rightmost node is unobtainable. But there is another Nash equilibrium available, namely $Rr$. In that \eqm, Player 2 gets her maximal payout, and Player 1 gets the most she can get given that Player 2 is playing $r$. But note what a strange \eqm\ it is. It relies on the idea that Player 2 would play $r$ were she to be given a move. But that is absurd. Once she gets a move, she has a straight choice between 1, if she plays $l$, and 0, if she plays $r$. So obviously she'll play $l$. 

This is just like the examples that we used to motivate subgame perfect \eqm, but that concept doesn't help us here. So we need a new concept. The core idea is that each player should be modelled as a Bayesian expected utility maximiser. More formally, the following constraints are put on players.

\begin{enumerate*}
\item At each point in the game, each player has a probability distribution over where they are in the game. These probability distributions are correct about the other players' actions. That is, if a player is playing a strategy $S$, everyone has probability 1 that they are playing $S$. If $S$ is a mixed strategy, this might involve having probabilities between 0 and 1 in propositions about which move the other player will make, but players have correct beliefs about other players' strategies.
\item No matter which node is reach, each player is disposed to maximise expected utility on arrival at that node.
\item When a player had a positive probability of arriving at a node, on reaching that node they update by conditionalisation.
\item When a player gave 0 probability to reaching a node (e.g., because the \eqm\ being played did not include that node), they have some disposition or other to form a set of consistent beliefs at that node.
\end{enumerate*}
The last constraint is very weak, but it does enough to eliminate the \eqm\ $Rr$. The constraint implies that when Player 2 moves, she must have some probability distribution $\Pr$ such that there's an $x$ such that $\Pr(L) = x$ and $\Pr(M) = 1-x$. Whatever value $x$ takes, the expected utility of $l$ given $\Pr$ is 1, and the expected utility of $r$ is 0. So being disposed to play $r$ violates the second condition. So $Rr$ is not a Perfect Bayesian \eqm.

It's true, but I'm not going to prove it, that all Perfect Bayesian equilibria are Nash equilibria. It's also true that the converse does not hold, and this we have proved; \textbf{Perfect} is an example.

\section{Signaling Games}

Concepts like Perfect Bayesian \eqm\ are useful for the broad class of games known as signaling games. In a signaling game, Player 1 gets some information that is hidden from player 2. Many applications of these games involve the information being \textit{de se} information, so the information that Player 1 gets is referred to as her \textit{type}. But that's inessential; what is essential is that only one player gets this information. Player 1 then has a choice of move to make. There is a small loss of generality here, but we'll restrict our attention to games where Player 1's choices are independent of her type, i.e., of the information she receives. Player 2 sees Player 1's move (but not, remember, her type) and then has a choice of her own to make. Again with a small loss of generality, we'll restrict attention to cases where Player 2's available moves are independent of what Player 1 does. 

We'll start, in game \textbf{Simple Signal} with a signaling game where the parties' interests are perfectly aligned.

\begin{figure}[ht]
\begin{center}
\begin{picture}(300, 300)

\put(150, 150){\circle{4}}
\pictext{150}{138}{$N$}
\pictext{110}{155}{$p$}
\pictext{190}{155}{$1-p$}
\put(152, 150){\line(1, 0){71}}
\put(77, 150){\line(1, 0){71}}

\put(75, 150){\circle*{4}}
\pictext{65}{145}{1}
\put(75, 152){\line(0, 1){68}}
\pictext{65}{180}{$U$}
\put(75, 148){\line(0, -1){68}}
\pictext{65}{110}{$D$}

\put(225, 150){\circle*{4}}
\pictext{235}{145}{1}
\put(225, 152){\line(0, 1){68}}
\pictext{235}{180}{$U$}
\put(225, 148){\line(0, -1){68}}
\pictext{235}{110}{$D$}

\multiput(75, 220)(12, 0){13}{\line(1, 0){6}}
\pictext{150}{225}{2}

\put(75, 220){\circle*{4}}
\put(75, 220){\line(-1, 1){50}}
\put(75, 220){\line(1, 1){50}}
\pictext{40}{240}{$l$}
\pictext{110}{240}{$r$}
\pictext{25}{275}{1, 1}
\pictext{125}{275}{0, 0}

\put(225, 220){\circle*{4}}
\put(225, 220){\line(-1, 1){50}}
\put(225, 220){\line(1, 1){50}}
\pictext{190}{240}{$l$}
\pictext{260}{240}{$r$}
\pictext{175}{275}{0, 0}
\pictext{275}{275}{1, 1}

\multiput(75, 80)(12, 0){13}{\line(1, 0){6}}
\pictext{150}{60}{2}

\put(75, 80){\circle*{4}}
\put(75, 80){\line(-1, -1){50}}
\put(75, 80){\line(1, -1){50}}
\pictext{40}{50}{$l$}
\pictext{110}{50}{$r$}
\pictext{25}{10}{1, 1}
\pictext{125}{10}{0, 0}

\put(225, 80){\circle*{4}}
\put(225, 80){\line(-1, -1){50}}
\put(225, 80){\line(1, -1){50}}
\pictext{190}{50}{$l$}
\pictext{260}{50}{$r$}
\pictext{175}{10}{0, 0}
\pictext{275}{10}{1, 1}

\end{picture}
\end{center}
\caption{\textbf{Simple Signal}}
\end{figure}

As above, we use an empty disc to signal the initial node of the game tree. In this case, it is the node in the centre of the tree. The first move is made by Nature, again denoted as $N$. Nature assigns a type to Player 1; that is, she makes some proposition true, and reveals it to Player 1. Call that proposition $q$. We'll say that Nature moves left if $q$ is true, and right otherwise. We assume the probability (in some good sense of probability) of $q$ is $p$, and this is known before the start of the game. After Nature moves, Player 1 has to choose $U$p or $D$own. Player 2 is shown Player 1's choice, but not Nature's move. That's the effect of the horizontal dashed lines. If the game reaches one of the upper nodes, Player 2 doesn't know which one it is, and if it reaches one of the lower nodes, again Player 2 doesn't know which it is. Then Player 2 has a make a choice, here simply denoted as $l$eft or $r$ight.

In any game of this form, each player has a choice of four strategies. Player 1 has to choose what to do if $q$ is true, and what to do if $q$ is false. We'll write $\alpha \beta$ for the strategy of doing $\alpha$ if $q$, and $\beta$ if $\neg q$. Since $\alpha$ could be $U$ or $D$, and $\beta$ could be $U$ or $D$, there are four possible strategies. Player 2 has to choose what to do if $U$ is played, and what to do if $D$ is played. We'll write $\gamma \delta$ for the strategy of doing $\gamma$ if $U$ is played, and $\delta$ if $D$ is played. Again, there are four possible choices. (If Player 2 knew what move Nature had made, there would be four degrees of freedom in her strategy choice, so she'd have 16 possible strategies. Fortunately, she doesn't have that much flexibility!)

Any game of this broad form is a signaling game. signaling games differ in (a) the interpretation of the moves, and (b) the payoffs. \textbf{Simple Signal} has a very simple payoff structure. Both players get 1 if Player 2 moves $l$ iff $q$, and 0 otherwise. If we think of $l$ as the formation of a belief that $q$, and $r$ as the formation of the opposite belief, this becomes a simple communication game. The players get a payoff iff Player 2 ends up with a true belief about $q$, so Player 1 is trying to communicate to Player 2 whether $q$ is true. This kind of simple communication game was used by David Lewis in his book \textit{Convention} to show that game theoretic approaches could be fruitful in the study of meaning. The game is perfectly symmetric if $p = \nicefrac{1}{2}$; so as to introduce some asymmetries, I'll work through the case where $p = \nicefrac{3}{5}$.

\textbf{Simple Signal} has a dizzying array of Nash equilibria, even given this asymmetry introducing assumption. They include the following:

\begin{itemize}
\item There are two \textbf{separating} equilibria, where what Player 1 does depends on what Nature does. These are \tol{UD, lr} and \tol{DU, rl}. These are rather nice equilibria; both players are guaranteed to get their maximal payout.
\item There are two \textbf{pooling} equilibria, where what Player 1 does is independent of what Nature does. These are \tol{UU, ll} and \tol{DD, ll}. Given that she gets no information from Player 1, Player 2 may as well guess. And since $\Pr(q) > \nicefrac{1}{2}$, she should guess that $q$ is true; i.e., she should play $l$. And given that Player 2 is going to guess, Player 1 may as well play anything. So these are also equilibria.
\item And there are some \textbf{babbling} equilibria. For instance, there is the \eqm\ where Player 1 plays either $UU$ or $DD$ with some probability $r$, and Player 2 plays $ll$.
\end{itemize}
Unfortunately, these are all Perfect Bayesian equilibria too. For the separating and babbling equilibria, it is easy to see that conditionalising on what Player 1 plays leads to Player 2 maximising expected utility by playing her part of the \eqm. And for the pooling equilibria, as long as the probability of $q$ stays above $\nicefrac{1}{2}$ in any `off-the-\eqm-path' play (e.g., conditional on $D$ in the \tol{UU, ll} \eqm), Player 2 maximises expected utility at every node.

That doesn't mean we have nothing to say. For one thing, the separating equilibria are Pareto-dominant in the sense that both players do better on those equilibria than they do on any other. So that's a non-coincidental reason to think that they will be the equilibria that arise. There are other refinements on Perfect Bayesian equilibria that are more narrowly tailored to signaling games. We'll introduce them by looking at a couple of famous signaling games.

\section{Signaling without Cooperation}

Economists have been interested for several decades in game-theoretic representations of college that give make going to college essentially a signal in a game. For example, consider the following variant of the signaling game, \textbf{College Signal}. It has the following intended interpretation:

\begin{itemize}
\item Player 1 is a student and potential worker, Player 2 is an employer.
\item The student is either bright or dull with probability $p$ of being bright. Nature reveals the type to the student, but only the probability to the employer, so $q$ is that the student is bright.
\item The student has the choice of going to college ($U$) or the beach ($D$).
\item The employer has the choice of hiring the student ($l$) or rejecting them ($r$).
\end{itemize}
In \textbf{College Signal} we make the following extra assumptions about the payoffs.

\begin{itemize}
\item Being hired is worth 4 to the student.
\item Going to college rather than the beach costs the bright student 1, and the dull student 5, since college is much harder work for dullards.
\item The employer gets no benefit from hiring college students as such.
\item Hiring a bright student pays the employer 1, hiring a dull student costs the employer 1, and rejections have no cost or benefit.
\end{itemize}
The resulting game tree is shown in \ref{SignalCollegeTree}. In this game, dull students never prefer to go to college, since even the lure of a job doesn't make up for the pain of actually having to study. So a rational strategy for Player 1 will never be of the form $\alpha U$, since for dull students, college is a dominated option, being dominated by $\alpha D$. But whether bright students should go to college is a trickier question. That is, it is trickier to say whether the right strategy for Player 1 is $UD$ or $DD$, which are the two strategies consistent with eliminating the strictly dominated strategies. (Remember that strictly dominated strategies cannot be part of a Nash \eqm.)

\begin{figure}[ht]
\begin{center}
\begin{picture}(300, 300)

\put(150, 150){\circle{4}}
\pictext{150}{138}{$N$}
\pictext{110}{155}{$p$}
\pictext{190}{155}{$1-p$}
\put(152, 150){\line(1, 0){71}}
\put(77, 150){\line(1, 0){71}}

\put(75, 150){\circle*{4}}
\pictext{65}{145}{1}
\put(75, 152){\line(0, 1){68}}
\pictext{65}{180}{$U$}
\put(75, 148){\line(0, -1){68}}
\pictext{65}{110}{$D$}

\put(225, 150){\circle*{4}}
\pictext{235}{145}{1}
\put(225, 152){\line(0, 1){68}}
\pictext{235}{180}{$U$}
\put(225, 148){\line(0, -1){68}}
\pictext{235}{110}{$D$}

\multiput(75, 220)(12, 0){13}{\line(1, 0){6}}
\pictext{150}{225}{2}

\put(75, 220){\circle*{4}}
\put(75, 220){\line(-1, 1){50}}
\put(75, 220){\line(1, 1){50}}
\pictext{40}{240}{$l$}
\pictext{110}{240}{$r$}
\pictext{25}{275}{3, 1}
\pictext{125}{275}{-1, 0}

\put(225, 220){\circle*{4}}
\put(225, 220){\line(-1, 1){50}}
\put(225, 220){\line(1, 1){50}}
\pictext{190}{240}{$l$}
\pictext{260}{240}{$r$}
\pictext{175}{275}{-1, -1}
\pictext{275}{275}{-5, 0}

\multiput(75, 80)(12, 0){13}{\line(1, 0){6}}
\pictext{150}{60}{2}

\put(75, 80){\circle*{4}}
\put(75, 80){\line(-1, -1){50}}
\put(75, 80){\line(1, -1){50}}
\pictext{40}{50}{$l$}
\pictext{110}{50}{$r$}
\pictext{25}{10}{4, 1}
\pictext{125}{10}{0, 0}

\put(225, 80){\circle*{4}}
\put(225, 80){\line(-1, -1){50}}
\put(225, 80){\line(1, -1){50}}
\pictext{190}{50}{$l$}
\pictext{260}{50}{$r$}
\pictext{175}{10}{4, -1}
\pictext{275}{10}{0, 0}

\end{picture}
\end{center}
\caption{\textbf{College Signal}}
\label{SignalCollegeTree}
\end{figure}

First, consider the case where $p < \nicefrac{1}{2}$. In that case, if Player 1 plays $DD$, then Player 2 gets $2p-1$ from playing either $ll$ or $rl$, and 0 from playing $lr$ or $rr$. So she should play one of the latter two options. If she plays $rr$, then $DD$ is a best response, since when employers aren't going to hire, students prefer to go to the beach. So there is one \textbf{pooling} \eqm, namely \tol{DD, rr}. But what if Player 2 plays $lr$. Then Player 1's best response is $UD$, since bright students prefer college conditional on it leading to a job. So there is also a \textbf{separating} \eqm, namely \tol{UD, lr}. The employer prefers that \eqm, since her payoff is now $p$ rather than 0. And students prefer it, since their payoff is $3p$ rather than 0. So if we assume that people will end up at Pareto-dominant outcomes, we have reason to think that bright students, and only bright students, will go to college, and employers will hire college students, and only college students. And all this is true despite there being no advantage whatsoever to going to college in terms of how good an employee one will be.

Especially in popular treatments of the case, the existence of this kind of model can be used to motivate the idea that college \textit{only} plays a signaling role. That is, some people argue that in the real world college does not make students more economically valuable, and the fact that college graduates have better employment rates, and better pay, can be explained by the signaling function of college. For what it's worth, I highly doubt that is the case. The wage premium one gets for going to college tends to \textit{increase} as one gets further removed from college, although the further out from college you get, the less important a signal college participation is. One can try to explain this fact too on a pure signaling model of college's value, but frankly I think the assumptions needed to do so are heroic. The model is cute, but not really a picture of how actual college works.

So far we assumed that $p < \nicefrac{1}{2}$. If we drop that assumption, and assume instead that $p > \nicefrac{1}{2}$, the case becomes more complicated. Now if Player 1 plays $DD$, i.e., goes to the beach no matter what, Player 2's best response is still to hire them. But note that now a very odd \eqm\ becomes available. The pair \tol{DD, rl} is a Nash \eqm, and, with the right assumptions, a Perfect Bayesian \eqm. This pair says that Player 1 goes to the beach whatever her type, and Player 2 hires only beach goers.

This is a very odd strategy for Player 2 to adopt, but it is a little hard to say just why it is odd. It is clearly a Nash \eqm. Given that Player 2 is playing $rl$, then clearly beach-going dominates college-going. And given that Player 1 is playing $DD$, playing $rl$ gets as good a return as is available to Player 2, i.e., $2p-1$. Could it also be a Perfect Bayesian equilibrium? It could, provided Player 2 has a rather odd update policy. Say that Player 2 thinks that if someone goes to college, they are a dullard with probability greater than $\nicefrac{1}{2}$. That's consistent with what we've said; given that Player 1 is playing $DD$, the probability that a student goes to college is 0. So the conditional probability that a college-goer is bright is left open, and can be anything one likes in Perfect Bayesian \eqm. So if Player 2 sets it to be, say, 0, then the rational reaction is to play $rl$.

But now note what an odd update strategy this is for Player 2. She has to assume that if someone deviates from the $DD$ strategy, it is someone for whom the deviation is strictly dominated. Well, perhaps it isn't crazy to assume that someone who would deviate from an \eqm\ isn't very bright, so maybe this isn't the oddest assumption in this particular context. But a few economists and game theorists have thought that we can put more substantive constraints on probabilities conditional on `off-the-\eqm-path' behaviour. One such constraint, is, roughly, that deviation shouldn't lead to playing a dominated strategy. This is the ``\textbf{intuitive criterion}'' of Cho and Kreps. In this game, all the criterion rules out is the odd pair \tol{DD, rl}. It doesn't rule out the very similiar pair \tol{DD, ll}. But the intuitive criterion makes more substantial constraints in other games. We'll close the discussion of signaling games with such an example, and a more careful statement of the criterion.

\section{The Intuitive Criterion}

The tree in \ref{SignalCanadianTree} represents \textbf{Canadian Signal}, which is also a guessing game. The usual statement of it involves all sorts of unfortunate stereotypes about the type of men who have quiche and/or beer for breakfast, and I'm underwhelmed by it. So I'll run with an example that relies on different stereotypes.

A North American tourist is in a bar. 60\% of the North Americans who pass through that bar are from Canada, the other 40\% are from the US. (This is a little implausible for most parts of the world, but maybe it is very cold climate bar. In Cho and Kreps' version the split is 90/10 not 60/40, but all that matters is which side of 50/50 it is.) The tourist, call her Player 1, knows her nationality, although the barman doesn't. The tourist can ask for the bar TV to be turned to hockey or to baseball, and she knows once she does that the barman will guess at her nationality. (The barman might also try to rely on her accent, but she has a fairly neutral upper-Midwest/central Canada accent.) Here are the tourist's preferences.

\begin{itemize}
\item If she is American, she has a (weak) preference for watching baseball rather than hockey.
\item If she is Canadian, she has a (weak) preference for watching hockey rather than baseball.
\item Either way, she has a strong preference for being thought of as Canadian rather than American. This preference is considerably stronger than her preference over which sport to watch.
\end{itemize}
The barman's preferences are simpler; he prefers to make true guesses to false guesses about the tourist's nationality. All of this is common knowledge. So the decision tree is given below. In this tree, $B$ means asking for baseball, $H$ means asking for hockey, $a$ means guessing the tourist is from the USA, $c$ means guessing she is from Canada.

\begin{figure}[ht]
\begin{center}
\begin{picture}(300, 300)

\put(150, 150){\circle{4}}
\pictext{150}{138}{$N$}
\pictext{110}{155}{$0.6$}
\pictext{190}{155}{$0.4$}
\put(152, 150){\line(1, 0){71}}
\put(77, 150){\line(1, 0){71}}

\put(75, 150){\circle*{4}}
\pictext{65}{145}{1}
\put(75, 152){\line(0, 1){68}}
\pictext{65}{180}{$H$}
\put(75, 148){\line(0, -1){68}}
\pictext{65}{110}{$B$}

\put(225, 150){\circle*{4}}
\pictext{235}{145}{1}
\put(225, 152){\line(0, 1){68}}
\pictext{235}{180}{$H$}
\put(225, 148){\line(0, -1){68}}
\pictext{235}{110}{$B$}

\multiput(75, 220)(12, 0){13}{\line(1, 0){6}}
\pictext{150}{225}{2}

\put(75, 220){\circle*{4}}
\put(75, 220){\line(-1, 1){50}}
\put(75, 220){\line(1, 1){50}}
\pictext{40}{240}{$c$}
\pictext{110}{240}{$a$}
\pictext{25}{275}{3, 1}
\pictext{125}{275}{1, -1}

\put(225, 220){\circle*{4}}
\put(225, 220){\line(-1, 1){50}}
\put(225, 220){\line(1, 1){50}}
\pictext{190}{240}{$c$}
\pictext{260}{240}{$a$}
\pictext{175}{275}{2, -1}
\pictext{275}{275}{0, 1}

\multiput(75, 80)(12, 0){13}{\line(1, 0){6}}
\pictext{150}{60}{2}

\put(75, 80){\circle*{4}}
\put(75, 80){\line(-1, -1){50}}
\put(75, 80){\line(1, -1){50}}
\pictext{40}{50}{$c$}
\pictext{110}{50}{$a$}
\pictext{25}{10}{2, 1}
\pictext{125}{10}{0, -1}

\put(225, 80){\circle*{4}}
\put(225, 80){\line(-1, -1){50}}
\put(225, 80){\line(1, -1){50}}
\pictext{190}{50}{$c$}
\pictext{260}{50}{$a$}
\pictext{175}{10}{3, -1}
\pictext{275}{10}{1, 1}

\end{picture}
\end{center}
\caption{\textbf{Canadian Signal}}
\label{SignalCanadianTree}
\end{figure}
There are two Nash equilibria for this game. One is \tol{HH, ca}; everyone asks for hockey, and the barman guesses Canadian if hockey, American if baseball. It isn't too hard to check this is an \eqm. The Canadian gets her best outcome, so it is clearly an \eqm\ for her. The American gets the second-best outcome, but asking for baseball would lead to a worse outcome. And given that everyone asks for hockey, the best the barman can do is go with the prior probabilities, and that means guessing Canadian. It is easy to see how to extend this to a perfect Bayesian \eqm; simply posit that conditional on baseball being asked for, the probability that the tourist is American is greater than $\nicefrac{1}{2}$.

The other \eqm\ is rather odder. It is \tol{BB, ac}; everyone asks for baseball, and the barman guesses American if hockey, Canadian if baseball. Again, it isn't too hard to see how it is Nash \eqm. The Canadian would rather have hockey, but not at the cost of being thought American, so she has no incentive to defect. The American gets her best outcome, so she has no incentive to defect. And the barman does as well as he can given that he gets no information out of the request, since everyone requests baseball.

Surprisingly, this could also be a perfect Bayesian \eqm. This requires that the barman maximise utility at every node. The tricky thing is to ensure he maxmises utility at the node where hockey is chosen. This can be done provided that conditional on hockey being chosen, the probability of American rises to above  $\nicefrac{1}{2}$. Well, nothing we have said rules that out, so there \textit{exists} a perfect Bayesian \eqm. But it doesn't seem very plausible. Why would the barman adopt just \textit{this} updating disposition?

One of the active areas of research on signaling games is the development of formal rules to rule out intuitively crazy updating dispositions like this. (Small note for the methods folks: at least some of the economists working on this explicitly say that the aim is to develop formal criteria that capture clear intuitions about cases.) I'm just going to note one of the very early attempts in this area, Cho and Kreps' \textit{Intuitive Criterion}.

Start with some folksy terminology that is (I think) novel to me. Let an outcome $o$ of the game be any combination of moves by nature and the players. Call the players'  moves collectively $P$ and nature's move $N$. So an outcome is a pair \tol{P, N}. Divide the players into three types.

\begin{itemize}
\item The \textbf{happy} players in $o$ are those such that given $N$, $o$ maximises their possible returns. That is, for all possible $P^\prime$, their payoff in \tol{P, N} is at least as large as their payoff in \tol{P^\prime, N}.
\item The \textbf{content} players are those who are not happy, but who are such that for no strategy $s^\prime$ which is an alternative to their current strategy $s$, would they do better given what nature and the other players would do if they played $s^\prime$.
\item The \textbf{unhappy} players are those who are neither happy nor content.
\end{itemize}
The standard Nash \eqm\ condition is that no player is unhappy. So assume we have a Nash \eqm\, with only happy and content players. And assume it is also a Perfect Bayesian \eqm. That is, for any `off-\eqm' outcome, each player would have some credence function such that their play maximises utility given that function.

Now add one constraint to those credence functions:

\begin{quote}
\textbf{Intuitive Criterion - Weak Version}

Assume a player has probability 1 that a strategy combination $P$ will be played. Consider their credence function conditional on another player playing $s^\prime$, wh\-ich is different to the strategy $s$ they play in $P$. If it is consistent with everything else the player believes that $s^\prime$ could be played by a player who is content with the actual outcome \tol{P, N}, then give probability 1 to $s^\prime$ being played by a content player.
\end{quote}

\noindent That is, if some deviation from \eqm\ happens, give probability 1 to it being one of the content players, not one of the happy players, who makes the deviation.

That's enough to rule out the odd \eqm\ for the baseball-hockey game. The only way that is a Perfect Bayesian \eqm\ is if the barman responds to a hockey request by \textit{increasing} the probability that the tourist is American. But in the odd \eqm, the American is happy, and the Canadian is merely content. So the Intuitive Criterion says that the barman should give credence 1 to the hockey requester, i.e., the deviator from \eqm, being Canadian. And if so, the barman won't respond to a hockey request by guessing the requestor is American, so the Canadian would prefer to request hockey. And that means the outcome is no longer an \eqm, since the Canadian is no longer even content with requesting baseball.

In games where everyone has only two options, the weak version I've given is equivalent to what Cho and Kreps offers. The official version is a little more complicated. First some terminology of mine, then their terminology next.

\begin{itemize}
\item A player is \textbf{happy to have played} $s$ rather than $s^\prime$ if the payoff for $s$ in $o$ is greater than any possible outcome for $s^\prime$ given $N$.
\item A player is \textbf{content to have played} $s$ rather than $s^\prime$ if they are not happy to have played $s$ rather than $s^\prime$, but the payoff for $s$ in $o$ is as great as the payoff for $s^\prime$ given $N$ and the other players' strategies.
\item A player is \textbf{unhappy to have played} $s$ rather than $s^\prime$ if they are neither happy nor content to have played $s$ rather than $s^\prime$.
\end{itemize}

\noindent If a player is happy to have played $s$ rather than $s^\prime$, and $s$ is part of some \eqm\ outcome $o$, then Cho and Kreps say that $s^\prime$ is an \textbf{equilibrium-dominated} strategy. The full version of the Intuitive Criterion is then:

\begin{quote}
\textbf{Intuitive Criterion - Full Version}

Assume a player has probability 1 that a strategy combination $P$ will be played. Consider their credence function conditional on another player playing $s^\prime$, wh\-ich is different to the strategy $s$ they play in $P$. If it is consistent with everything else the player believes that $s^\prime$ could be played by a player who is merely content to have played $s$ rather than $s^\prime$,  then give probability 1 to $s^\prime$ being played by someone who is content to have played $s$, rather than someone who is happy to have played $s$.
\end{quote}

\noindent This refinement matters in games where players have three or more options. It might be that a player's options are $s_1, s_2$ and  $s_3$, and their type is $t_1$ or $t_2$. In the Nash \eqm\ in question, they play $s_1$. If they are type $t_1$, they are happy to have played $s_1$ rather than $s_2$, but merely content to have played $s_1$ rather than $s_3$. If they are type $t_2$, they are happy to have played $s_1$ rather than $s_3$, but merely content to have played $s_1$ rather than $s_2$. In my terminology above, both players are content rather than happy with the outcome. So the weak version of the Intuitive Criterion wouldn't put any restrictions on what we can say about them conditional on them playing, say $s_2$. But the full version does say something; it says other players should assign probability 1 to the player being type $t_2$ conditional on them playing $s_2$, since the alternative is that $s_2$ is played by a player who is happy they played $s_1$ rather than $s_2$. Similarly, it says that conditional on the player playing $s_3$, other players should assign probability 1 to their being of type $t_1$, since the alternative is to assign positive probability to a player deviating to a strategy they are happy not to play.

The Intuitive Criterion has been the subject of an enormous literature. Goo\-gle Scholar lists nearly 2000 citations for the Cho and Kreps paper alone, and similar rules are discussed in other papers. So there are arguments that it is too weak and arguments too strong, and refinements in all sorts of directions. But we'll leave those. What I most wanted to stress was the \textit{form} a refinement of Perfect Bayesian \eqm\ would take. It is a constraint on the conditional credences of agents conditional on some probability 0 event happening. It's interesting that there are some plausible, even intuitive constraints; in general it seems the economic literature has investigated rationality under 0 probability evidence more than philosophers have.


\part{Voting Theory}

\chapter{Group Decisions}

So far, we've been looking at the way that an individual may make a decision. In practice, we are just as often concerned with group decisions as with individual decisions. These range from relatively trivial concerns (e.g. Which movie shall we see tonight?) to some of the most important decisions we collectively make (e.g. Who shall be the next President?). So methods for grouping individual judgments into a group decision seem important.

Unfortunately, it turns out that there are several challenges facing any attempt to merge preferences into a single decision. In this chapter, we'll look at various approaches that different groups take to form decisions, and how these different methods may lead to different results. The different methods have different strengths and, importantly, different weaknesses. We might hope that there would be a method with none of these weaknesses. Unfortunately, this turns out to be impossible. 

One of the most important results in modern decision theory is the Arrow Impossibility Theorem, named after the economist Kenneth Arrow who discovered it. The Arrow Impossibility Theorem says that there is no method for making group decisions that satisfies a certain, relatively small, list of desiderata. We will set out the theorem, and explore a little what those constraints are.

Then we'll look a bit at real world voting systems, and their different strengths and weaknesses. Different democracies use quite different voting systems to determine the winner of an election. (Indeed, within the United States there is an interesting range of systems used.) And some theorists have promoted the use of yet other systems than are currently used. Choosing a voting system is not quite like choosing a method for making a group decision. We will start off making the assumption that when we're looking at ways to aggregate individual preferences into a group decision, we'll assume that we have clear access to the preferences of individual agents. A voting system is not meant to tally preferences into a decision, it is meant to tally votes. And voters may have reasons (some induced by the system itself) for voting in ways other than their preferences. For instance, many voters in American presidential elections vote for their preferred candidate of the two major candidates, rather than `waste' their vote on a third party candidate.

For now we'll put those problems to one side, and assume that members of the group express themselves honestly when voting. Still, it turns out there are complications that arise for even relatively simple decisions.

\section{Making a Decision}
Seven friends, who we'll imaginatively name $F_1, F_2, ..., F_7$ are trying to decide which restaurant to go to. They have four options, which we'll also imaginatively name $R_1, R_2, R_3, R_4$. The first thing they do is ask which restaurant each person prefers. The results are as follows.

\begin{itemize}
\item $F_1, F_2$ and $F_3$ all vote for $R_1$, so it gets 3 votes
\item $F_4$ and $F_5$ both vote for $R_2$, so it gets 2 votes
\item $F_6$ votes for $R_3$, so it gets 1 vote
\item $F_7$ votes for $R_4$, so it gets 1 vote
\end{itemize}
It looks like $R_1$ should be the choice then. It, after all, has the most votes. It has a `plurality' of the votes - that is, it has the most votes. In most American elections, the candidate with a plurality wins. This is sometimes known as plurality voting, or (for unclear reasons) first-past-the-post or winner-take-all. The obvious advantage of such a system is that it is easy enough to implement.

But it isn't clear that it is the ideal system to use. Only 3 of the 7 friends wanted to go to $R_1$. Possibly the other friends are all strongly opposed to this particular restaurant. It seems unhappy to choose a restaurant that a majority is strongly opposed to, especially if this is avoidable.

So the second thing the friends do is hold a `runoff' election. This is the method used for voting in some U.S. states (most prominently in Georgia and Louisiana) and many European countries. The idea is that if no candidate (or in this case no restaurant) gets a majority of the vote, then there is a second vote, held just between the top two vote getters. Since $R_1$ and $R_2$ were the top vote getters, the choice will just be between those two. When this vote is held the results are as follows.

\begin{itemize}
\item $F_1, F_2$ and $F_3$ all vote for $R_1$, so it gets 3 votes
\item $F_4, F_5, F_6$ and $F_7$ all vote for $R_2$, so it gets 4 votes
\end{itemize}
This is sometimes called `runoff' voting, for the natural reason that there is a runoff. Now we've at least arrived at a result that the majority may not have as their first choice, but which a majority are at least happy to vote for.

But both of these voting systems seem to put a lot of weight on the various friends' first preferences, and less weight on how they rank options that aren't optimal for them. There are a couple of notable systems that allow for these later preferences to count. For instance, here is how the polls in American college sports work. A number of voters rank the best teams from 1 to $n$, for some salient $n$ in the relevant sport. Each team then gets a number of points per ballot, depending on where it is ranked, with $n$ points for being ranked first, $n-1$ points for being ranked second, $n-2$ points for being ranked third, and so on down to 1 point for being ranked $n$'th. The teams' overall ranking is then determined by who has the most points.

In the college sport polls, the voters don't rank every team, only the top $n$, but we can imagine doing just that. So let's have each of our friends rank the restaurants in order, and we'll give 4 points to each restaurant that is ranked first, 3 to each second place, etc. The points that each friend awards are given by the following table.

\begin{center}
\begin{tabular}{r | c c c c c c c c}
 & $F_1$ & $F_2$ & $F_3$ & $F_4$ & $F_5$ & $F_6$ & $F_7$ & Total \\ \hline
$R_1$ & 4 & 4 & 4 & 1 & 1 & 1 & 1 & 16 \\
$R_2$ & 1 & 3 & 3 & 4 & 4 & 2 & 2 & 19 \\
$R_3$ & 3 & 2 & 2 & 3 & 3 & 4 & 3 & 20 \\
$R_4$ & 2 & 1 & 1 & 2 & 2 & 3 & 4 & 15
\end{tabular}
\end{center}
Now we have yet a different choice. By this method, $R_3$ comes out as the best option. This voting method is sometimes called the Borda count. The nice advantage of it is that it lets all preferences, not just first preferences, count. Note that previously we didn't look at all at the preferences of the first three friends, beside noting that $R_1$ is their first choice. Note also that $R_3$ is no one's least favourite option, and is many people's second best choice. These seem to make it a decent choice for the group, and it is these facts that the Borda count is picking up on.

But there is something odd about the Borda count. Sometimes when we prefer one restaurant to another, we prefer it by just a little. Other times, the first is exactly what we want, and the second is, by our lights, terrible. The Borda count tries to approximately measure this - if $X$ strongly prefers $A$ to $B$, then often there will be many choices between $A$ and $B$, so $A$ will get many more points on $X$'s ballot. But this is not necessary. It is possible to have a strong preference for $A$ over $B$ without there being any live option that is `between' them.  In any case, why try to come up with some proxy for strength of preference when we can measure it directly?

That's what happens if we use `range voting'. Under this method, we get each voter to give each option a score, say a number between 0 and 10, and then add up all the scores. This is, approximately, what's used in various sporting competitions that involve judges, such as gymnastics or diving. In those sports there is often some provision for eliminating the extreme scores, but we won't be borrowing that feature of the system. Instead, we'll just get each friend to give each restaurant a score out of 10, and add up the scores. Here is how the numbers fall out.

\begin{center}
\begin{tabular}{r | c c c c c c c c}
 & $F_1$ & $F_2$ & $F_3$ & $F_4$ & $F_5$ & $F_6$ & $F_7$ & Total \\ \hline
$R_1$ & 10 & 10 & 10 & 5 & 5 & 5 & 0 & 45 \\
$R_2$ & 7 & 9 & 9 & 10 & 10 & 7 & 1 & 53\\
$R_3$ & 9 & 8 & 8 & 9 & 9 & 10 & 2 & 55 \\
$R_4$ &8 & 7 & 7 & 8 & 8 & 9 & 10 & 57
\end{tabular}
\end{center}Now $R_4$ is the choice! But note that the friends' individual preferences have not changed throughout. The way each friend would have voted in the previous `elections' is entirely determined by their scores as given in this table. But using four different methods for aggregating preferences, we ended up with four different decisions for where to go for dinner.

I've been assuming so far that the friends are accurately expressing their opinions. If the votes came in just like this though, some of them might wonder whether this is really the case. After all, $F_7$ seems to have had an outsized effect on the overall result here. We'll come back to this when looking at options for voting systems.

\section{Desiderata for Preference Aggregation Mechanisms}
None of the four methods we used so far are obviously crazy. But they lead to four different results. Which of these, if any, is the correct result? Put another way, what is the ideal method for aggregating preferences? One natural way to answer this question is to think about some desirable features of aggregation methods. We'll then look at which systems have the most such features, or ideally have all of them.

One feature we'd like is that each option has a chance of being chosen. It would be a very bad preference aggregation method that didn't give any possibility to, say, $R_3$ being chosen.

More strongly, it would be bad if the aggregation method chose an option $X$ when there was another option $Y$ that everyone preferred to $X$. Using some terminology from the game theory notes, we can express this constraint by saying our method should never choose a Pareto inferior option. Call this the \textbf{Pareto condition}.

We might try for an even stronger constraint. Some of the time, not always but some of the time, there will be an option $C$ such than a majority of voters prefers $C$ to $X$, for every alternative $X$. That is, in a two-way match-up between $C$ and any other option $X$, $C$ will get more votes. Such an option is sometimes called a Condorcet option, after Marie Jean Antoine Nicolas Caritat, the Marquis de Condorcet, who discussed such options. The \textbf{Condorcet condition} on aggregation methods is that a Condorcet option always comes first, if such an option exists.

Moving away from these comparative norms, we might also want our preference aggregation system to be fair to everyone. A method that said $F_2$ is the dictator, and $F_2$'s preferences are the group's preferences, would deliver a clear answer, but does not seem to be particularly fair to the group. There should be \textbf{no dictators}; for any person, it is possible that the group's decision does not match up with their preference.

More generally than that, we might restrict attention to preference aggregation systems that don't pay attention to \textit{who} has various preferences, just to \textit{what} preferences people have. Here's one way of stating this formally. Assume that two members of the group, $v_1$ and $v_2$, swap preferences, so $v_1$'s new preference ordering is $v_2$'s old preference ordering and vice versa. This shouldn't change what the group's decision is, since from a group level, nothing has changed. Call this the \textbf{symmetry} condition.

Finally, we might want to impose a condition that we said is a condition we imposed on independent agents: the \textbf{irrelevance of independent alternatives}. If the group would choose $A$ when the options are $A$ and $B$, then they wouldn't choose $B$ out of any larger set of options that also include $A$. More generally, adding options can change the group's choice, but only to one of the new options.

\section{Assessing Plurality Voting}
It is perhaps a little disturbing to think how few of those conditions are met by plurality voting, which is how Presidents of the USA are elected. (More precisely, it is how the electoral votes in most states are distributed. The President is chosen by a more complicated mechanism.) Plurality voting clearly satisfies the \textbf{Pareto condition}. If everyone prefers $A$ to $B$, then $B$ will get no votes, and so won't win. So far so good. And since any one person might be the only person who votes for their preferred candidate, and since other candidates might get more than one vote, no one person can dictate who wins. So it satisfies \textbf{no dictators}. Finally, since the system only looks at votes, and not at who cast them, it satisfies \textbf{symmetry}.

But it does not satisfy the \textbf{Condorcet condition}. Consider an election with three candidates. $A$ gets 40\% of the vote, $B$ gets 35\% of the vote, and $C$ gets 25\% of the vote. $A$ wins, and $C$ doesn't even finish second. But assume also that everyone who didn't vote for $C$ has her as their second preference after either $A$ or $B$. Something like this may happen if, for instance, $C$ is an independent moderate, and $A$ and $B$ are doctrinaire candidates from the major parties. Then 60\% prefer $C$ to $A$, and 65\% prefer $C$ to $B$. So $C$ is a Condorcet candidate, yet is not elected.

A similar example shows that the system does not satisfy the \textbf{irrelevance of independent alternatives} condition. If $B$ was not running, then presumably $A$ would still have 40\% of the vote, while $C$ would have 60\% of the vote, and would win. One thing you might want to think about is how many elections in recent times would have had the outcome changed by eliminating (or adding) unsuccessful candidates in this way. 

\section{Arrow's Theorem}

Arrow's Theorem is an important mathematical result about voting systems. It needs a bit of setup to state carefully. The theorem says that there is no way to aggregate individual preferences into a group preference ordering that satisfies some intuitive criteria.

To get started, we'll assume that each agent has a \textbf{complete} and \textbf{transitive} preference ordering over the options. If we say $A >_V B$ means that $V$ prefers $A$ to $B$, that $A =_V B$ means that $V$ is indifferent between $A$ and $B$, and that $A \geq_V B$ means that $A >_V B \vee A =_V B$, then these constraints can be expressed as follows.

\begin{description}
\item[Completeness] For any voter $V$ and options $A, B$, either $A \geq_V B$ or $B \geq_V A$
\item[Transitivity] For any voter $V$ and options $A, B$, the following three conditions hold:
	\begin{itemize}
	\item If $A >_V B$ and $B >_V C$ then $A >_V C$
	\item If $A =_V B$ and $B =_V C$ then $A =_V C$
	\item If $A \geq_V B$ and $B \geq_V C$ then $A \geq_V C$
	\end{itemize}
\end{description}
More generally, we assume the \textbf{substitutivity of indifferent options}. That is, if $A =_V B$, then whatever is true of the agent's attitude towards $A$ is also true of the agent's attitude towards $B$. In particular, whatever comparison holds in the agent's mind between $A$ and $C$ holds between $B$ and $C$. (The last two bullet points under transitivity follow from this principle about indifference and the earlier bullet point.)

The effect of these assumptions is that we can represent the agent's preferences by lining up the options from best to worst, with the possibility that we'll have to put two options in one `spot' to represent the fact that the agent values each of them equally.

A \textbf{ranking function} is a function from the preference orderings of the agent to a new preference ordering, which we'll call the preference ordering of the group. We'll use the subscript $_G$ to note that it is the group's ordering we are designing. We'll also assume that the group's preference ordering is complete and transitive. 

There are any number ranking functions that don't look at all like the \textit{group's} preferences in any way. For instance, if the function is meant to work out the results of an election, we could consider the function that takes any input whatsoever, and returns a ranking that simply lists the candidates by age, with the oldest first, the second oldest second, etc. This doesn't seem like it is the group's preferences in any way. Whatever any member of the group thinks, the oldest candidate wins. What Arrow called the \textbf{citizen sovereignty} condition is that for any possible ranking of the candidates, it should be possible to have the group end up with that ranking.

The citizen sovereignty follows from another constraint we might put on ranking functions. If everyone in the group prefers $A$ to $B$, then $A >_G B$, i.e. the group prefers $A$ to $B$. We'll call this the \textbf{Pareto} constraint. It is sometimes called the \textbf{unanimity} constraint, but we'll call it the Pareto condition.

One way to satisfy the Pareto constraint is to pick a particular person, and make them dictator. That is, the function `selects' a person $V$, and says that $A >_G B$ if and only if $A >_V B$. If everyone prefers $A$ to $B$, then $V$ will, so this is consistent with the Pareto constraint. But it also doesn't seem like a way of constructing the group's preferences. So let's say that we'd like a non-dictatorial ranking function. More carefully, for any given person, and any preference ordering they may have, there is some 

The last constraint is one we discussed in the previous chapter: the \textbf{independence of irrelevant alternatives}. Formally, this means that whether $A >_G B$ is true depends only on how the voters rank $A$ and $B$. So changing how the voters rank, say $B$ and $C$, doesn't change what the group says about the $A$, $B$ comparison.

It's sometimes thought that it would be a very good thing if the voting system respected this constraint. Let's say that you believe that if Gary Johnston had not been a candidate in the 2000 U.S. Presidential election, then Hillary Clinton, not Donald Trump, would have won the election. (I have no idea whether this is actually true, but I can see why someone might believe it.) Then you might think it is a little odd that whether Clinton or Trump wins depends on who else is in the election, and not on the voters' preferences between Clinton and Trump. This is a special case of the independence of irrelevant alternatives - you think that the voting system should end up with the result that it would have come up with had there been just those two candidates. If we generalise this motivation a lot, we get the conclusion that third possibilities should be irrelevant.

Unfortunately, we've now got ourselves into an impossible situation. Arrow's theorem says that any ranking function that satisfies the Pareto and independence of irrelevant alternatives constraints, has a dictator in any case where the number of alternatives is greater than 2. When there are only 2 choices, majority rule satisfies all the constraints. But nothing, other than dictatorship, works in the general case.

\section{Cyclic Preferences}
We can see why three option cases are a problem by considering one very simple example. Say there are three voters, $V_1, V_2, V_3$ and three choices $A, B, C$. The agent's rankings are given in the table below. (The column under each voter lists the choices from their first preference, on top, to their least favourite option, on the bottom.)

\begin{center}
\begin{tabular}{c c c}
$V_1$ & $V_2$ & $V_3$ \\ 
$A$ & $B$ & $C$ \\
$B$ & $C$ & $A$ \\
$C$ & $A$ & $B$
\end{tabular}
\end{center}
If we just look at the $A/B$ comparison, $A$ looks pretty good. After all, 2 out of 3 voters prefer $A$ to $B$. But if we look at the $B/C$ comparison, $B$ looks pretty good. After all, 2 out of 3 voters prefer $B$ to $C$. So perhaps we should say $A$ is best, $B$ second best and $C$ worst. But wait! If we just look at the $C/A$ comparison, $C$ looks pretty good. After all, 2 out of 3 voters prefer $C$ to $A$.

It might seem like one natural response here is to say that the three options should be tied. The group preference ranking should just be that $A =_G B =_G = C$. But note what happens if we say that and accept independence of irrelevant alternatives. If we eliminate option $C$, then we shouldn't change the group's ranking of $A$ and $B$. That's what independence of irrelevant alternatives says. So now we'll be left with the following rankings.

\begin{center}
\begin{tabular}{c c c}
$V_1$ & $V_2$ & $V_3$ \\ 
$A$ & $B$ & $A$ \\
$B$ & $A$ & $B$ \\
\end{tabular}
\end{center}
By independence of irrelevant alternatives, we should still have $A =_G B$. But 2 out of 3 voters wanted $A$ over $B$. The one voter who preferred $B$ to $A$ is making it that the group ranks them equally. That's a long way from making them a dictator, but it's our first sign that our constraints give excessive power to one voter. One other thing the case shows is that we can't have the following three conditions on our ranking function.

\begin{itemize}
\item If there are just two choices, then the majority choice is preferred by the group.
\item If there are three choices, and they are symmetrically arranged, as in the table above, then all choices are equally preferred.
\item The ranking function satisfies independence of irrelevant alternatives.
\end{itemize}
I noted after the example that $V_2$ has quite a lot of power. Their preference makes it that the group doesn't prefer $A$ to $B$. We might try to generalize this power. Maybe we could try for a ranking function that worked strictly by consensus. The idea would be that if everyone prefers $A$ to $B$, then $A >_G B$, but if there is no consensus, then $A =_G B$. Since how the group ranks $A$ and $B$ only depends on how individuals rank $A$ and $B$, this method easily satisfies independence of irrelevant alternatives. And there are no dictators, and the method satisfies the Pareto condition. So what's the problem?

Unfortunately, the consensus method described here violates transitivity, so doesn't even produce a group preference ordering in the formal sense we're interested in. Consider the following distribution of preferences.

\begin{center}
\begin{tabular}{c c c}
$V_1$ & $V_2$ & $V_3$ \\ 
$A$ & $A$ & $B$ \\
$B$ & $C$ & $A$ \\
$C$ & $B$ & $C$
\end{tabular}
\end{center}
Everyone prefers $A$ to $C$, so by unanimity, $A >_G C$. But there is no consensus over the $A/B$ comparison. Two people prefer $A$ to $B$, but one person prefers $B$ to $A$. And there is no consensus over the $B/C$ comparison. Two people prefer $B$ to $C$, but one person prefers $C$ to $B$. So if we're saying the group is indifferent between any two options over which there is no consensus, then we have to say that $A =_G B$, and $B =_G C$. By transitivity, it follows that $A =_G C$, contradicting our earlier conclusion that $A >_G C$.

This isn't going to be a formal argument, but we might already be able to see a difficulty here. Just thinking about our first case, where the preferences form a cycle suggests that the only way to have a fair ranking consistent with independence of irrelevant alternatives is to say that the group only prefers options when there is a consensus in favour of that option. But the second case shows that consensus based methods do not in general produce \textit{rankings} of the options. So we have a problem. Arrow's Theorem shows how deep that problem goes.


\section{Proofs of Arrow's Theorem}
The proofs of Arrow's Theorem, though not particularly long, are a little tricky to follow. So we won't go through them in any detail at all. But I'll sketch one proof due to John Geanakopolos of the Cowles Foundation at Yale.\footnote{The proof is available at http://ideas.repec.org/p/cwl/cwldpp/1123r3.html.} Geanakopolos assumes that we have a ranking function that satisfies Pareto and independence of irrelevant alternatives, and aims to show that in this function there must be a dictator.

The first thing he proves is a rather nice lemma. Assume that every voter puts some option $B$ on either the top or the bottom of their preference ranking. Don't assume they all agree: some people hold that $B$ is the very best option, and the rest hold that it is the worst. Geanakopolos shows that in this case the ranking function must put $B$ either at the very top or the very bottom.

To see this, assume that it isn't true. So there are some options $A$ and $C$ such that $A \geq_G B$ and $B \geq_G C$. Now imagine changing each voter's preferences so that $C$ is moved above $A$ while $B$ stays where it is - either on the top or the bottom of that particular voter's preferences. By Pareto, we'll now have $C >_G A$, since everyone prefers $C$ to $A$. But we haven't changed how any person thinks about any comparison involving $B$. So by independence of irrelevant alternatives, $A \geq_G B$ and $B \geq_G C$ must still be true. By transitivity, it follows that $A \geq_G C$, contradicting our conclusion that $C >_G A$.

This is a rather odd conclusion I think. Imagine that we have four voters with the following preferences. 

\begin{center}
\begin{tabular}{c c c c}
$V_1$ & $V_2$ & $V_3$ & $V_4$ \\ 
$B$ & $B$ & $A$ & $C$ \\
$A$ & $C$ & $C$ & $A$ \\
$C$ & $A$ & $B$ & $B$
\end{tabular}
\end{center}
By what we've proven so far, $B$ has to come out either best or worst in the group's rankings. But which should it be? Since half the people love $B$, and half hate it, it seems it should get a middling ranking. One lesson of this is that independence of irrelevant alternatives is a very strong condition, one that we might want to question.

The next stage of Geanakopolos's proof is to consider a situation where at the start everyone thinks $B$ is the very worst option out of some long list of options. One by one the voters change their mind, with each voter in turn coming to think that $B$ is the best option. By the result we proved above, at every stage of the process, $B$ must be either the worst option according to the group, or the best option. $B$ starts off as the worst option, and by Pareto $B$ must end up as the best option. So at one point, when one voter changes their mind, $B$ must go from being the worst option on the group's ranking to being the best option, simply in virtue of that person changing their mind.

We won't go through the rest, but the proof continues by showing that that person has to be a dictator. Informally, the idea is to prove two things about that person, both of which are derived by repeated applications of independence of irrelevant alternatives. First, this person has to retain their power to move $B$ from worst to first whatever the other people think of $A$ and $C$. Second, since they can make $B$ jump all options by changing their mind about $B$, if they move $B$ `halfway', say they come to have the view $A >_V B >_V C$, then $B$ will jump (in the group's ranking) over all options that it jumps over in this voter's rankings. But that's possible (it turns out) only if the group's ranking of $A$ and $C$ is dependent entirely on this voter's rankings of $A$ and $C$. So the voter is a dictator with respect to this pair. A further argument shows that the voter is a dictator with respect to every pair, which shows there must be a dictator.

\section{Voting Systems}

The Arrow Impossibility Theorem shows that we can't have everything that we want in a voting system. In particular, we can't have a voting system that takes as inputs the preferences of each voter, and outputs a preference ordering of the group that satisfies these three constraints.

\begin{enumerate*}
\item \textbf{Unanimity}: If everyone prefers $A$ to $B$, then the group prefers $A$ to $B$.
\item \textbf{Independence of Irrelevant Alternatives}: If nobody changes their mind about the relative ordering of $A$ and $B$, then the group can't change its mind about the relative ordering of $A$ and $B$.
\item \textbf{No Dictators}: For each voter, it is possible that the group's ranking will be different to their ranking
\end{enumerate*}
Any voting system either won't be a function in the sense that we're interested in for Arrow's Theorem, or will violate some of those constraints. (Or both.) But still there could be better or worse voting systems. Indeed, there are many voting systems in use around the world, and serious debate about which is best. In these notes we'll look at the pros and cons of a few different voting systems.

The discussion here will be restricted in two respects. First, we're only interested in systems for making political decisions, indeed, in systems for electing representatives to political positions. We're not interested in, for instance, the systems that a group of friends might use to choose which movie to see, or that an academic department might use to hire new faculty. Some of the constraints we'll be looking at are characteristic of elections in particular, not of choices in general.

Second, we'll be looking only at elections to fill a single position. This is a fairly substantial constraint. Many elections are to fill multiple positions. The way a lot of electoral systems work is that many candidates are elected at once, with the number of representatives each party gets being (roughly) in proportion to the number of people who vote for that party. This is how the parliament is elected in many countries around the world (including, for instance, Mexico, Germany and Spain). Perhaps more importantly, it is basically the norm for new parliaments to have such kind of multi-member constituencies. But the mathematical issues get a little complicated when we look at the mechanisms for selecting multiple candidates, and we'll restrict ourselves to looking at mechanisms for electing a single candidate.

\section{Plurality voting}
By far the most common method used in America, and throughout much of the rest of the world, is plurality voting. Every voter selects one of the candidates, and the candidates with the most votes wins. As we've already noted, this is called plurality, or first-past-the-post, voting.

Plurality voting clearly does not satisfy the independence of irrelevant alternatives condition. We can see this if we imagine that the voting distribution starts off with the table on the left, and ends with the table on the right. (The three candidates are $A$, $B$ and $C$, with the numbers at the top of each column representing the percentage of voters who have the preference ordering listed below it.)

\begin{center}
\begin{tabular}{c c c p{100pt} c c c}
40\% & 35\% & 25\% & & 40\% & 35\% & 25\% \\
\cmidrule(r){1-3}
\cmidrule(r){5-7}
$A$ & $B$ & $C$ & & $A$ & $B$ & $B$ \\
$B$ & $A$ & $B$ & & $B$ & $A$ & $C$ \\
$C$ & $C$ & $A$ & & $C$ & $C$ & $A$
\end{tabular}
\end{center}
All that happens as we go from left-to-right is that some people who previously favoured $C$ over $B$, come to favour $B$ over $C$. Yet this change, which is completely independent of how anyone feels about $A$, is sufficient for $B$ to go from losing the election 40-35 to winning the election 60-40.

This is how we show that a system does not satisfy independent of irrelevant alternatives - coming up with a pair of situations where no voter's opinion about the relative merits of two choices (in this case $A$ and $B$) changes, but the group's ranking of those two choices changes.

One odd effect of this is that whether $B$ wins the election depends not just on how voters compare $A$ and $B$, but on how voters compare $B$ and $C$. One of the consequences of Arrow's Theorem might be taken to be that this kind of thing is unavoidable, but it is worth stopping to reflect on just how pernicious this is to the democratic system. 

Imagine that we are in the left-hand situation, and you are one of the 25\% of voters who like $C$ best, then $B$ then $A$. It seems that there is a reason for you to not vote the way your preferences go; you'll have a better chance of electing a candidate you prefer if you vote, against your preferences, for $B$. So the voting system might encourage voters to not express their preferences adequately. This can have a snowball effect - if in one election a number of people who prefer $C$ vote for $B$, at future elections other people who might have voted for $C$ will also vote for $B$ because they don't think enough other people share their preferences for $C$ to make such a vote worthwhile.

Indeed, if the candidate $C$ themselves strongly prefers $B$ to $A$, but thinks a lot of people will vote for them if they run, then $C$ might even be discouraged from running because it will lead to a worse election result. This doesn't seem like a democratically ideal situation.

Some of these consequences are inevitable consequences of a system that doesn't satisfy independence of irrelevant alternatives. And the Arrow Theorem shows that it is hard to avoid independence of irrelevant alternatives. But some of them seem like serious democratic shortcomings, the effects of which can be seen in American democracy, and especially in the extreme power the two major parties have. (Though, to be fair, a number of other electoral systems that use plurality voting do not have such strong major parties. Indeed, Canada seems to have very strong third parties despite using this system, as does the United Kingdom.)

One clear advantage of plurality voting should be stressed: it is quick and easy. There is little chance that voters will not understand what they have to do in order to express their preferences. And voting is, or at least should be, relatively quick. The voter just has to make one mark on a piece of paper, or press a single button, to vote. When the voter is expected to vote for dozens of offices, as is usual in America (though not elsewhere) this is a serious benefit. In many U.S. elections, we have seen queues hours long of people waiting to vote. Were voting any slower than it actually is, these queues might have been worse.

Relatedly, it is easy to count the votes in a plurality system. You just sort all the votes into different bundles and count the size of each bundle. Some of the other systems we'll be looking at are much harder to count the votes in. This can lead to serious delays in even being able to announce the results of an election.

\section{Runoff Voting}
One solution to some of the problems with plurality voting is runoff voting, which is used in parts of America (including Georgia and Louisiana and, in practice, California) and is very common throughout Europe and South America. The most notable upcoming election using this system is the Presidential election in France. 

The idea is that there are, in general, two elections. At the first election, if one candidate has majority support, then they win. But otherwise the top two candidates go into a runoff. In the runoff, voters get to vote for one of those two candidates, and the candidate with the most votes wins.

This doesn't entirely deal with the problem of a spoiler candidate having an outsized effect on the election, but it makes such cases a little harder to produce. For instance, imagine that there are four candidates, and the arrangement of votes is as follows.

\begin{center}
\begin{tabular}{c c c c}
35\% & 30\% & 20\% & 15\% \\ 
$A$ & $B$ & $C$ & $D$ \\
$B$ & $D$ & $D$ & $C$ \\
$C$ & $C$ & $B$ & $B$ \\
$D$ & $A$ & $A$ & $A$
\end{tabular}
\end{center}
In a plurality election, $A$ will win with only 35\% of the vote.\footnote{This isn't actually that unusual in the overall scope of American elections, especially in primaries.} In a runoff election, the runoff will be between $A$ and $B$, and presumably $B$ will win, since 65\% of the voters prefer $B$ to $A$. But look what happens if $D$ drops out of the election, or all of $D$'s supporters decide to vote more strategically.

\begin{center}
\begin{tabular}{c c c c}
35\% & 30\% & 20\% & 15\% \\ 
$A$ & $B$ & $C$ & $C$ \\
$B$ & $C$ & $B$ & $B$ \\
$C$ & $A$ & $A$ & $A$ \\
\end{tabular}
\end{center}
Now the runoff is between $C$ and $A$, and $C$ will win. $D$ being a candidate means that the candidate most like $D$, namely $C$, loses a race they could have won.

In one respect this is much like what happens with plurality voting. On the other hand, it is somewhat harder to find real life cases that show this pattern of votes. That's in part because it is hard to find cases where there are (a) four serious candidates, and (b) the third and fourth candidates are so close ideologically that they eat into each other's votes and (c) the top two candidates are so close that these third and fourth candidates combined could leapfrog over each of them. Theoretically, the problem about spoiler candidates might look as severe, but it is much less of a problem in practice.

The downside of runoff voting of course is that it requires people to go and vote twice. This can be a major imposition on the time and energy of the voters. More seriously from a democratic perspective, it can lead to an unrepresentative electorate. In American runoff elections, the runoff typically has a much lower turnout than the initial election, so the election comes down to the true party loyalists. At least, this is true in the southern states that have traditional runoff elections. In the Californian system, where the primary becomes a de facto first round of voting it is the other way around. In Europe, the first round has sometimes had a very low turnout, with fringe candidates making it to the final round on the strength of a small but loyal supporter base.

\section{Instant Runoff Voting}
One approach to this problem is to do, in effect, the initial election and the runoff at the same time. In instant runoff voting, every voter lists their preference ordering over their desired candidates. In practice, that means marking `1' beside their first choice candidate, `2' beside their second choice and so on through the candidates. 

When the votes are being counted, the first thing that is done is to count how many first-place votes each candidate gets. If any candidate has a majority of votes, they win. If not, the candidate with the lowest number of votes is eliminated. The vote counter then distributes each ballot for that eliminated candidate to whichever candidate receives the `2' vote on that ballot. If that leads to a candidate having a majority, that candidate wins. If not, the candidate with the lowest number of votes at this stage is eliminated, and their votes are distributed, each voter's vote going to their most preferred candidate of the remaining candidates. This continues until a candidate gets a majority of the votes.

This avoids the particular problem we discussed about runoff voting. In that case, $D$ would have been eliminated at the first round, and $D$'s votes would all have flowed to $C$. That would have moved $C$ about $B$, eliminating $B$. Then with $B$'s preferences, $C$ would have won the election comfortably. But it doesn't remove all problems. In particular, it leads to an odd kind of strategic voting possibility. The following situation does arise, though rarely. Imagine the voters are split the following way.
\begin{center}
\begin{tabular}{c c c}
45\% & 28\% &27\% \\
$A$ & $B$ & $C$ \\
$B$ & $A$ & $B$ \\
$C$ & $C$ & $A$
\end{tabular}
\end{center}
As things stand, $C$ will be eliminated. And when $C$ is eliminated, all of $C$'s votes will be transferred to $B$, leading to $B$ winning. Now imagine that a few of $A$'s voters change the way they vote, voting for $C$ instead of their preferred candidate $A$, so now the votes look like this.
\begin{center}
\begin{tabular}
{c c c c}
43\% & 28\% &27\% & 2\% \\ 
$A$ & $B$ & $C$ & $C$\\
$B$ & $A$ & $B$ & $A$ \\
$C$ & $C$ & $A$ & $B$
\end{tabular}
\end{center}
Now $C$ has more votes than $B$, so $B$ will be eliminated. But $B$'s voters have $A$ as their second choice, so now $A$ will get all the new votes, and $A$ will easily win. Some theorists think that this possibility for strategic voting is a sign that instant runoff voting is flawed.

Perhaps a more serious worry is that the voting and counting system is more complicated. This slows down voting itself, though this is a problem can be partially dealt with by having more resources dedicated to making it possible to vote. The vote count is also somewhat slower. A worse consequence is that because the voter has more to do, there is more chance for the voter to make a mistake. In some jurisdictions, if the voter does not put a number down for each candidate, their vote is invalid, even if it is clear which candidate they wish to vote for. It also requires the voter to have opinions about all the candidates running, and this may include a number of frivolous candidates. But it isn't clear that this is a major problem if it does seem worthwhile to avoid the problems with plurality and runoff voting.

The other systems we'll look at are not in wide use around the world.


\section{Borda Count}
In a Borda Count election, each voter ranks each of the candidates, as in Instant Runoff Voting. Each candidate then receives $n$ points for each first place vote they receive (where $n$ is the number of candidates), $n-1$ points for each second place vote, and so on through the last place candidate getting 1 point. The candidate with the most points wins.

One nice advantage of the Borda Count is that it eliminates the chance for the kind of strategic voting that exists in Instant Runoff Voting, or for that matter any kind of Runoff Voting. It can never make it more likely that $A$ will win by someone changing their vote away from $A$. Indeed, this could only lead to $A$ having fewer votes. This certainly seems to be reasonable.

Another advantage is that many preferences beyond first place votes count. A candidate who is every single voter's second best choice will not do very well under any voting system that gives a special weight to first preferences. But such a candidate may well be in a certain sense the best representative of the whole community.

And a third advantage is that the Borda Count includes a rough approximation of voter's strength of preference. If one voter ranks $A$ a little above $B$, and another votes $B$ many places above $A$, that's arguably a sign that $B$ is a better representative of the two of them than $A$. Although only one of the two prefers $B$, one voter will be a little disappointed that $B$ wins, while the other would be very disappointed if $B$ lost.

These are not trivial advantages. But there are also many disadvantages which explain why no major electoral system has adopted Borda Count voting yet, despite its strong support from some theorists.

First, Borda Count is particularly complicated to implement. It is just as difficult for the voter to as in Instant Runoff Voting; in each case they have to express a complete preference ordering. But it is much harder to count, because the vote counter has to detect quite a bit of information from each ballot. Getting this information from millions of ballots is not a trivial exercise.

Second, Borda Count has a serious problem with `clone candidates'. In plurality voting, a candidate suffers if there is another candidate much like them on the ballot. In Borda Count, a candidate can seriously gain if such a candidate is added. Consider the following situation. In a certain electorate, of say 100,000 voters, 60\% of the voters are Republicans, and 40\% are Democrats. But there is only one Republican, call them $R$, on the ballot, and there are 2 Democrats, $D1$ and $D2$ on the ballot. Moreover, $D2$ is clearly a worse candidate than $D1$, but the Democrats still prefer the Democrat to the Republican. Since the district is overwhelmingly Republican, intuitively the Republican should win. But let's work through what happens if 60,000 Republicans vote for $R$, then $D1$, then $D2$, and the 40,000 Democrats vote $D1$ then $D2$ then $R$. In that case, $R$ will get $60,000 \times 3 + 40,000 \times 1 = 220,000$ points, $D1$ will get $60,000 \times 2 + 40,000 \times 3 = 240,000$ points, and $D2$ will get $60,000 \times 1 + 40,000 \times 2 = 140,000$ points, and $D1$ will win. Having a `clone' on the ticket was enough to push $D1$ over the top.

On the one hand, this may look a lot like the mirror image of the `spoiler' problem for plurality voting. But in another respect it is much worse. It is hard to get someone who is a lot ideologically like your opponent to run in order to improve your electoral chances. It is much easier to convince someone who already wants you to win to add their name to the ballot in order to improve your chances. In practice, this would either lead to an arms race between the two parties, each trying to get the most names onto the ballot, or very restrictive (and hence undemocratic) rules about who was even allowed to be on the ballot, or, most likely, both.

The third problem comes from thinking through the previous problem from the point of view of a Republican voter. If the Republican voters realise what is up, they might vote tactically for $D2$ over $D1$, putting $R$ back on top. In a case where the electorate is as partisan as in this case, this might just work. But this means that Borda Count is just as susceptible to tactical voting as other systems; it is just that the tactical voting often occurs downticket. (There are more complicated problems, that we won't work through, about what happens if the voters mistakenly judge what is likely to happen in the election, and tactical voting backfires.)

Finally, it's worth thinking about whether the supposed major virtue of Borda Count, the fact that it considers all preferences and not just first choices, is a real gain. The core idea behind Borda Count is that all preferences should count equally. So the difference between first place and second place in a voter's affections counts just as much as the difference between third and fourth. But for many elections, this isn't how the voters themselves feel. I suspect many people reading this have strong feelings about who was the best candidate in the past Presidential election. I suspect very few people had strong feelings about who was the third best versus fourth best candidate. This is hardly a coincidence; people identify with a party that is their first choice. They say, ``I'm a Democrat'' or ``I'm a Green'' or ``I'm a Republican''. They don't identify with their third versus fourth preference. Perhaps voting systems that give primary weight to first place preferences are genuinely reflecting the desires of the voters. 

\section{Approval Voting}
In plurality voting, every voter gets to vote for one candidate, and the candidate with the most votes wins. Approval voting is similar, except that each voter is allowed to vote for as many candidates as they like. The votes are then added up, and the candidate with the most votes wins. Of course, the voter has an interest in not voting for too many candidates. If they vote for all of the candidates, this won't advantage any candidate; they may as well have voted for no candidates at all.

The voters who are best served by approval voting, at least compared to plurality voting, are those voters who wish to vote for a non-major candidate, but who also have a preference between the two major candidates. Under approval voting, they can vote for the minor candidate that they most favor, and also vote for the the major candidate who they hope will win. Of course, runoff voting (and Instant Runoff Voting) also allow these voters to express a similar preference. Indeed, the runoff systems allow the voters to express not only two preferences, but express the order in which they hold those preferences. Under approval voting, the voter only gets to vote for more than one candidate, they don't get to express any ranking of those candidates.

But arguably approval voting is easier on the voter. The voter can use a ballot that looks just like the ballot used in plurality voting. And they don't have to learn about preference flows, or Borda Counts, to understand what is going on in the voting. Currently there are many voters who vote for, or at least appear to try to vote for, multiple candidates. This is presumably inadvertent, but approval voting would let these votes be counted, which would refranchise a number of voters. Approval voting has never been used as a mass electoral tool, so it is hard to know how quick it would be to count, but presumably it would not be incredibly difficult.\footnote{There is one circumstance when something like Approval Voting is used in the United States. It is not uncommon when there are $n > 1$ candidates to be elected to a position, like school board members, that each voter can select $n$ candidates, and the top $n$ vote getters are elected. This isn't true Approval Voting, since there is still a cap on how many people you can vote for.}

One striking thing about approval voting is that it is not a function from voter preferences to group preferences. Hence it is not subject to the Arrow Impossibility Theorem. It isn't such a function because the voters have to not only rank the candidates, they have to decide where on their ranking they will `draw the line' between candidates that they will vote for, and candidates that they will not vote for. Consider the following two sets of voters. In each case candidates are listed from first preference to last preference, with stars indicating which candidates the voters vote for.


\begin{center}
\begin{tabular}{c c c p{100pt} c c c}
40\% & 35\% & 25\% & & 40\% & 35\% & 25\% \\
\cmidrule(r){1-3}
\cmidrule(r){5-7}
$*A$ & $*B$ & $*C$ & & $*A$ & $*B$ & $*C$ \\
$B$ & $A$ & $B$ & & $B$ & $A$ & $*B$ \\
$C$ & $C$ & $A$ & & $C$ & $C$ & $A$
\end{tabular}
\end{center}
In the election on the left-hand-side, no voter takes advantage of approval voting to vote for more than one candidate. So $A$ wins with 40\% of the vote. In the election on the right-hand-side, no one's preferences change. But the 25\% who prefer $C$ also decide to vote for $B$. So now $B$ has 60\% of the voters voting for them, as compared to 40\% for $A$ and 25\% for $C$, so $B$ wins.

This means that the voting system is not a function from voter preferences to group preferences. If it were a function, fixing the group preferences would fix who wins. But in this case, without a single voter changing their preference ordering of the candidates, a different candidate won. Since the Arrow Impossibility Theorem only applies to functions from voter preferences to group preferences, it does not apply to Approval Voting.

\section{Range Voting}
In Range Voting, every voter gives each candidate a score. Let's say that score is from 0 to 10. The name `Range' comes from the range of options the voter has. In the vote count, the score that each candidate receives from each voter is added up, and the candidate with the most points wins.

In principle, this is a way for voters to express very detailed opinions about each of the candidates. They don't merely rank the candidates, they measure how much better each candidate is than all the other candidates. And this information is then used to form an overall ranking of the various candidates.

In practice, it isn't so clear this would be effective. Imagine that a voter $V$ thinks that candidate $A$ would be reasonably good, and candidate $B$ would be merely OK, and that no other candidates have a serious chance of winning. If $V$ was genuinely expressing their opinions, they might think that $A$ deserves an 8 out of 10, and $B$ deserves a 5 out of 10. But $V$ wants $A$ to win, since $V$ thinks $A$ is the better candidate. And $V$ knows that what will make the biggest improvement in $A$'s chances is if they score $A$ a 10 out of 10, and $B$ a 0 out of 10. That will give $A$ a 10 point advantage, whereas they may only get a 3 point advantage if the voter voted sincerely.

It isn't unusual for a voter to find themselves in $V$'s position. So we might suspect that although Range Voting will give the voters quite a lot of flexibility, and give them the chance to express detailed opinions, it isn't clear how often it would be in a voter's interests to use these options.

And Range Voting is quite complex, both from the perspective of the voter and of the vote counter. There is a lot of information to be gleaned from each ballot in Range Voting. This means the voter has to go to a lot of work to fill out the ballot, and the vote counter has to do a lot of work to process all that information. This means that Range Voting might be very slow, both in terms of voting and counting. And if voters have a tactical reason for not wanting to fill in detailed ballots, this might mean it's a lot of effort for not a lot of reward, and that we should stick to somewhat simpler vote counting methods.

\part{Group Minds}

\end{document}